\documentclass[11pt,oneside]{book}

\usepackage[margin=1in,headheight=15pt]{geometry}
\usepackage[T1]{fontenc}
\usepackage{lmodern}
\usepackage{microtype}
\usepackage{mathtools,amssymb,amsthm}
\usepackage{tikz}
\usetikzlibrary{arrows.meta,calc,decorations.pathreplacing,patterns,positioning,shapes.geometric}
\usepackage[most]{tcolorbox}
\usepackage{enumitem}
\usepackage{booktabs,tabularx,array,multicol}
\usepackage{xcolor}
\usepackage{fancyhdr}
\usepackage{hyperref}

\definecolor{ink}{HTML}{162332}
\definecolor{navy}{HTML}{17324D}
\definecolor{teal}{HTML}{167D82}
\definecolor{gold}{HTML}{D69E2E}
\definecolor{coral}{HTML}{D8664A}
\definecolor{mist}{HTML}{EDF5F4}
\definecolor{paper}{HTML}{FBFAF7}
\definecolor{muted}{HTML}{5D6B78}

\pagecolor{paper}
\color{ink}
\hypersetup{colorlinks=true,linkcolor=teal,urlcolor=coral,citecolor=teal,pdfauthor={Course design by OpenAI Codex},pdftitle={Patterns, Proofs, and Primes}}
\setlength{\parindent}{0pt}
\setlength{\parskip}{5pt}
\setlist{itemsep=2pt,topsep=5pt}
\setlist[enumerate,1]{label=\textbf{\arabic*.},leftmargin=*}

\pagestyle{fancy}
\fancyhf{}
\fancyhead[L]{\small\color{muted}\nouppercase{\leftmark}}
\fancyhead[R]{\small\color{muted}Patterns, Proofs, and Primes}
\fancyfoot[C]{\color{teal}\thepage}
\renewcommand{\headrulewidth}{0.3pt}
\renewcommand{\chaptermark}[1]{\markboth{#1}{}}

\newtheoremstyle{course}{6pt}{6pt}{\itshape}{} {\bfseries\color{navy}}{.}{.5em}{}
\theoremstyle{course}
\newtheorem{theorem}{Theorem}[chapter]
\newtheorem{lemma}[theorem]{Lemma}
\newtheorem{proposition}[theorem]{Proposition}
\newtheorem{corollary}[theorem]{Corollary}
\theoremstyle{definition}
\newtheorem{definition}[theorem]{Definition}
\newtheorem{example}[theorem]{Example}
\newtheorem{question}[theorem]{Guiding question}

\newtcolorbox{bigidea}[1][]{enhanced,breakable,colback=mist,colframe=teal,boxrule=.8pt,arc=2mm,left=5mm,borderline west={3pt}{0pt}{teal},title={Big idea},fonttitle=\bfseries,#1}
\newtcolorbox{checkpoint}[1][]{enhanced,breakable,colback=white,colframe=gold,boxrule=.7pt,arc=2mm,title={Checkpoint},fonttitle=\bfseries\color{ink},#1}
\newtcolorbox{warningbox}[1][]{enhanced,breakable,colback=coral!5,colframe=coral,boxrule=.7pt,arc=2mm,title={Common trap},fonttitle=\bfseries,#1}
\newtcolorbox{labbox}[1][]{enhanced,breakable,colback=navy!4,colframe=navy,boxrule=.7pt,arc=2mm,title={Explore},fonttitle=\bfseries,#1}
\newtcolorbox{exercisebox}[1][]{enhanced,breakable,colback=white,colframe=navy!45,boxrule=.6pt,arc=1mm,title={Exercises},fonttitle=\bfseries\color{navy},#1}

\newcommand{\Z}{\mathbb Z}
\newcommand{\N}{\mathbb N}
\newcommand{\Q}{\mathbb Q}
\newcommand{\R}{\mathbb R}
\newcommand{\F}{\mathbb F}
\newcommand{\divides}{\mid}
\newcommand{\ndivides}{\nmid}
\newcommand{\congr}{\equiv}
\newcommand{\floor}[1]{\left\lfloor #1\right\rfloor}
\newcommand{\abs}[1]{\left|#1\right|}
\newcommand{\angles}[1]{\left\langle #1\right\rangle}
\newcommand{\exlevel}[1]{\hfill\textcolor{teal}{\small[#1]}}
\newcommand{\chaptertag}[1]{\par\vspace{-1.4em}{\color{teal}\large\sffamily #1}\par\vspace{1em}}

\begin{document}

\hypersetup{pageanchor=false}
\begin{titlepage}
\begin{tikzpicture}[remember picture,overlay]
  \fill[navy] (current page.south west) rectangle (current page.north east);
  \foreach \r in {1.2,2.0,2.8,3.6,4.4}{
    \draw[teal!55,line width=1.2pt] ($(current page.north east)+(-3.2,-3.0)$) circle (\r cm);
  }
  \foreach \a in {0,30,...,330}{
    \fill[gold] ($(current page.north east)+(-3.2,-3.0)+(\a:2.8)$) circle (2.2pt);
  }
  \draw[coral,line width=3pt] ($(current page.south west)+(0,2.0)$) -- ($(current page.south east)+(0,2.0)$);
\end{tikzpicture}
\vspace*{4.2cm}
{\color{gold}\sffamily\bfseries\Large AN INTRODUCTION TO NUMBER THEORY\par}
\vspace{.5cm}
{\color{white}\fontsize{36}{42}\selectfont\bfseries Patterns, Proofs,\par and Primes\par}
\vspace{1cm}
{\color{white!78}\Large A visual, inquiry-driven course with exercises\par}
\vfill
{\color{white!75}\large One-semester course packet\par}
{\color{white!55}\normalsize Euclid $\to$ continued fractions $\to$ congruences $\to$ groups $\to$ prime counting\par}
\end{titlepage}

\hypersetup{pageanchor=true}
\frontmatter
\chapter*{How to use this course}
\addcontentsline{toc}{chapter}{How to use this course}

\begin{bigidea}
Number theory is an unusually good place to learn what proof feels like. The objects are familiar, experiments are cheap, and elementary questions quickly demand new ways of thinking. This course repeatedly follows the same rhythm:
\[
\text{compute} \longrightarrow \text{notice} \longrightarrow \text{conjecture} \longrightarrow \text{prove} \longrightarrow \text{reuse}.
\]
\end{bigidea}

\section*{Audience and prerequisites}

This packet is designed for a first proof-oriented course: strong secondary students, first- or second-year undergraduates, or mathematically mature non-majors. Students should be comfortable with algebra, functions, and mathematical induction. No calculus or abstract algebra is required. A short ``proof studio'' in the first week can supply any missing experience with direct proof, contradiction, and induction.

\section*{Learning goals}

By the end, students should be able to:
\begin{enumerate}
  \item use the Euclidean algorithm both computationally and as a source of proofs;
  \item prove unique factorization from Euclid's lemma and use prime exponents to prove $\sqrt2$ irrational;
  \item generate primes with the sieve of Eratosthenes and justify its stopping rule;
  \item encode continued-fraction convergents by unimodular $2\times2$ matrix products, connect them to extended-Euclidean coefficients, and use them as alternating bounds;
  \item explain rational base-$b$ periodicity and prove the Lagrange--Galois characterization of periodic continued fractions using fixed points, conjugation, and reduced surds;
  \item prove the symmetry of square-root periods and identify exactly
  which convergents solve the positive and negative Pell equations;
  \item calculate in $\Z/n\Z$, solve linear congruences, and use the Chinese remainder theorem;
  \item explain Fermat's little theorem through cyclic symmetry and through finite groups;
  \item prove basic identities and divisibility statements for binomial coefficients;
  \item understand why central binomial coefficients control the density of primes;
  \item write short proofs that clearly separate evidence, claim, and justification.
\end{enumerate}

\section*{A 14-week route}

The design assumes two 75-minute meetings per week. ``Studio'' means that at least half the meeting is devoted to student work at boards or in small groups.

\begin{center}
\small
\begin{tabularx}{\textwidth}{>{\bfseries}p{.07\textwidth}p{.39\textwidth}X}
\toprule
Week & Core mathematics & Suggested studio or assessment \\
\midrule
1 & Divisibility; gcd; Euclidean algorithm & Tiling model; proof diagnostic \\
2 & B\'ezout identity; linear Diophantine equations & Extended Euclid relay \\
3 & Primes; unique factorization; $\sqrt2\notin\Q$; the sieve & Sieve studio; repair a flawed proof; Problem Set 1 \\
4 & Finite continued fractions; matrix products; extended Euclid; calendar design & Match remainders with convergent coefficients \\
5 & Alternating convergents; Lagrange and Galois; Pell equations & Fixed points and reduced-state cycles; Problem Set 2 \\
6 & Congruences; units; linear congruences & Clock arithmetic and affine ciphers \\
7 & Chinese remainder theorem & CRT construction studio; short quiz \\
8 & Necklace colorings and Fermat's theorem & Orbit counting with physical beads; Problem Set 3 \\
9 & Groups, subgroups, cyclic groups & Cayley-table investigation \\
10 & Lagrange's theorem; group proof of Fermat & Compare two proofs; Problem Set 4 \\
11 & Binomial coefficients and Pascal's triangle & Identity stations \\
12 & $p$-adic valuations; Kummer's carry idea & Color divisibility patterns; Problem Set 5 \\
13 & Central binomial coefficient; Chebyshev functions & Prime-counting proof workshop \\
14 & Chebyshev estimate; synthesis & Capstone presentations or oral exam \\
\bottomrule
\end{tabularx}
\end{center}

Chapter~3 now gives full elementary periodicity proofs. A compact lecture
route is to introduce linear fractional transformations alongside the
matrix recurrence, do the fixed-point and prefix-conjugation calculations
in class, and assign the three-step finite-state argument as guided reading.
Return to the reduced-predecessor formula to prove Galois's criterion;
Exercises~20--26 provide practice. Section~\ref{sec:pell} then proves
the square-root palindrome and the exact Pell remainder identity,
including the period-parity rule and generation of all positive
solutions. Exercises~27--31 reinforce those arguments.
Section~\ref{sec:archimedes-cattle}, Exercises~32--34, and the first
notebook develop Archimedes' cattle problem as an optional application;
its modular computation can also be revisited after congruences.
The metric-theory and astronomical asides remain optional.

\section*{Assessment design}

One workable balance is 45\% problem sets, 20\% two short proof interviews, 15\% computational or visual labs, and 20\% capstone. Problem sets should mix three modes:
\begin{itemize}
  \item \textbf{Fluency}: focused calculations that make later proofs readable.
  \item \textbf{Reasoning}: multi-step proofs, often with a hint or intermediate claim.
  \item \textbf{Exploration}: generate data, state a conjecture, then test its limits.
\end{itemize}

The exercises in each chapter carry these labels. A typical weekly assignment is two fluency problems, two reasoning problems, and one exploration. The appendices supply a short matrix toolkit, continued-fraction arithmetic, an astronomical Venus case study, selected hints and solutions, and reference sheets. An instructor can remove the solution pages when distributing assignments.

\section*{Teaching principles}

\begin{itemize}
  \item \textbf{Let the picture precede the notation.} Tiling, winding around a clock, rotating a necklace, and counting lattice paths give students something to reason about before symbols compress the idea.
  \item \textbf{Prove one theorem twice.} Fermat's little theorem is the deliberate centerpiece: orbit counting supplies a concrete proof, while groups reveal the reusable structure.
  \item \textbf{Make proof choices visible.} Ask not only ``Is this correct?'' but ``Why did this definition unlock that step?''
  \item \textbf{Use computation as a microscope, not an oracle.} The optional Sage notebook helps students find examples and conjectures; written work must still explain why a pattern continues.
\end{itemize}

\section*{Dependency map}

\begin{center}
\begin{tikzpicture}[
  node distance=8mm and 11mm,
  box/.style={draw=navy!55,rounded corners=2mm,fill=white,minimum height=9mm,align=center,font=\small,inner xsep=4mm},
  arr/.style={-{Stealth[length=2mm]},thick,teal}
]
\node[box] (euclid) {Euclidean\\algorithm};
\node[box,right=of euclid] (fta) {Unique\\factorization};
\node[box,below=of euclid] (cf) {Continued\\fractions};
\node[box,right=of fta] (mod) {Modular\\arithmetic};
\node[box,below=of mod] (neck) {Necklace\\orbits};
\node[box,right=of mod] (group) {Finite\\groups};
\node[box,right=of neck] (flt) {Fermat's little\\theorem};
\node[box,below=of fta] (binom) {Binomial\\coefficients};
\node[box,below=of flt] (cheb) {Chebyshev's\\estimate};
\draw[arr] (euclid)--(fta);
\draw[arr] (euclid)--(cf);
\draw[arr] (fta)--(mod);
\draw[arr] (mod)--(neck);
\draw[arr] (mod)--(group);
\draw[arr] (neck)--(flt);
\draw[arr] (group)--(flt);
\draw[arr] (fta)--(binom);
\draw[arr] (binom)--(cheb);
\draw[arr] (fta) to[bend left=18] (cheb);
\end{tikzpicture}
\end{center}

\begin{checkpoint}
\textbf{A good opening prompt.} Write $391$ and $299$ on the board. Ask: ``What is the largest square tile that covers a $391\times299$ rectangle without cutting? How could you be certain?'' Do not define gcd until the class has invented the subtraction or division process.
\end{checkpoint}

\tableofcontents

\mainmatter
\chapter{The algorithm hidden in a rectangle}
\chaptertag{Divisibility, gcd, and B\'ezout's identity}

\begin{question}
What is the largest square tile that exactly covers a $391\times299$ rectangle?
\end{question}

If a square has integer side length $d$, then it fits precisely only when $d\divides391$ and $d\divides299$. We are looking for the \emph{greatest common divisor}, written $\gcd(391,299)$.

\section{Divisibility as structure}

For integers $a,b$, we say $a\divides b$ when $b=ak$ for some $k\in\Z$. Divisibility is not an operation like division; it is a yes-or-no relation.

\begin{proposition}[Division algorithm]
If $a\in\Z$ and $b>0$, there are unique integers $q,r$ such that
\[
a=bq+r,\qquad 0\le r<b.
\]
\end{proposition}

The quotient $q$ says how many full copies of $b$ fit inside $a$; the remainder $r$ is what is left. The crucial observation is
\[
\gcd(a,b)=\gcd(b,r)\quad\text{whenever }a=bq+r.
\]
Indeed, an integer divides both $a$ and $b$ exactly when it divides both $b$ and $a-bq=r$.

\section{Euclid's algorithm}

Repeat the division algorithm until the remainder is zero:
\begin{align*}
391&=1\cdot299+92,\\
299&=3\cdot92+23,\\
92&=4\cdot23+0.
\end{align*}
The last nonzero remainder is $23$, so $\gcd(391,299)=23$.

\begin{center}
\begin{tikzpicture}[scale=.72]
  \fill[navy!7] (0,0) rectangle (7.82,5.98);
  \draw[navy,thick] (0,0) rectangle (7.82,5.98);
  \fill[teal!25] (0,0) rectangle (5.98,5.98);
  \draw[teal,thick] (0,0) rectangle (5.98,5.98);
  \foreach \y in {0,1.84,3.68}{
    \fill[gold!35] (5.98,\y) rectangle (7.82,{\y+1.84});
    \draw[gold!80!black] (5.98,\y) rectangle (7.82,{\y+1.84});
  }
  \fill[coral!30] (5.98,5.52) rectangle (6.44,5.98);
  \foreach \x in {6.44,6.90,7.36}{
    \fill[coral!30] (\x,5.52) rectangle ({\x+.46},5.98);
    \draw[coral] (\x,5.52) rectangle ({\x+.46},5.98);
  }
  \draw[coral] (5.98,5.52) rectangle (6.44,5.98);
  \node[below] at (3.91,0) {$391$};
  \node[rotate=90,above] at (0,2.99) {$299$};
  \node[teal!70!black] at (2.99,2.99) {$299\times299$};
  \node[gold!60!black,rotate=90] at (6.9,2.76) {three $92\times92$ squares};
  \node[coral!80!black,above] at (6.9,5.98) {four $23\times23$ squares};
\end{tikzpicture}

\small The same remainders appear when successively removing the largest possible squares.
\end{center}

\begin{theorem}[Euclidean algorithm]
For positive integers $a\ge b$, repeated division terminates, and the last nonzero remainder is $\gcd(a,b)$.
\end{theorem}

\begin{proof}
Each remainder is a nonnegative integer smaller than the previous positive remainder, so the process must stop. Replacing $(a,b)$ by $(b,r)$ does not change the common divisors. At the last line $r_{k-1}=q r_k$, hence the common divisors are precisely the divisors of $r_k$.
\end{proof}

\section{Running the movie backward}

The remainders do more than find the gcd. Reverse the equations:
\begin{align*}
23&=299-3\cdot92\\
  &=299-3(391-299)\\
  &=4\cdot299-3\cdot391.
\end{align*}
Thus $23$ is an integer linear combination of $391$ and $299$.

\begin{theorem}[B\'ezout's identity]
For integers $a,b$, not both zero, there are integers $x,y$ such that
\[
ax+by=\gcd(a,b).
\]
Moreover, $\gcd(a,b)$ is the smallest positive integer of the form $ax+by$.
\end{theorem}

\begin{corollary}
The equation $ax+by=c$ has integer solutions exactly when $\gcd(a,b)\divides c$.
\end{corollary}

\begin{proof}
If $ax+by=c$, every common divisor of $a,b$ divides $c$. Conversely, multiply a B\'ezout identity for the gcd by $c/\gcd(a,b)$.
\end{proof}

\begin{warningbox}
The gcd is positive, but B\'ezout coefficients need not be. Restricting $x,y$ to positive integers makes the statement false.
\end{warningbox}

\begin{labbox}
Choose two 3-digit integers. Make a two-column record: Euclid's divisions on the left, backward substitutions on the right. How large can the coefficients become? The optional Sage notebook automates the arithmetic but also asks you to certify its output by substitution.
\end{labbox}

\begin{exercisebox}
\begin{enumerate}
  \item Compute $\gcd(2024,748)$ by Euclid's algorithm. \exlevel{Fluency}
  \item Express $\gcd(527,341)$ in the form $527x+341y$. \exlevel{Fluency}
  \item Prove: if $d\divides a$ and $d\divides b$, then $d\divides (ma+nb)$ for all $m,n\in\Z$. \exlevel{Reasoning}
  \item Find all integer solutions to $35x+22y=1$. Describe how one solution generates all the others. \exlevel{Reasoning}
  \item Show that $\gcd(a,b)=\gcd(a+kb,b)$ for every $k\in\Z$. Explain why this is the engine of Euclid's algorithm. \exlevel{Reasoning}
  \item A rectangle has integer sides $a>b$. Repeatedly remove the largest square available. Prove that the final square has side $\gcd(a,b)$. \exlevel{Exploration}
  \item Consecutive Fibonacci numbers make Euclid's algorithm take many steps. Test this claim, then prove $\gcd(F_{n+1},F_n)=1$. \exlevel{Exploration}
\end{enumerate}
\end{exercisebox}


\chapter{Every integer has an address}
\chaptertag{Primes and the fundamental theorem of arithmetic}

Positive integers factor. But why should the factorization into primes be unique? The arithmetic step that makes uniqueness possible is modest enough to miss.

\begin{definition}
An integer $p>1$ is \emph{prime} if its only positive divisors are $1$ and $p$. An integer $n>1$ that is not prime is \emph{composite}.
\end{definition}

\section{Euclid's lemma}

\begin{lemma}[Euclid]
If $p$ is prime and $p\divides ab$, then $p\divides a$ or $p\divides b$.
\end{lemma}

\begin{proof}
Suppose $p\ndivides a$. Since the only positive divisors of $p$ are $1,p$, we have $\gcd(p,a)=1$. B\'ezout gives $px+ay=1$. Multiply by $b$:
\[
pbx+aby=b.
\]
Both terms on the left are divisible by $p$, so $p\divides b$.
\end{proof}

\begin{warningbox}
The word ``prime'' is essential. We have $6\divides 2\cdot3$, but $6$ divides neither factor.
\end{warningbox}

Applying the lemma repeatedly shows: if a prime divides $a_1a_2\cdots a_k$, it divides at least one factor.

\section{Existence and uniqueness}

\begin{theorem}[Fundamental theorem of arithmetic]
Every integer $n>1$ can be written as a product of primes. Apart from rearranging the factors, this product is unique.
\end{theorem}

\begin{proof}[Existence]
Use strong induction on $n$. If $n$ is prime, we are done. If it is composite, write $n=ab$ with $1<a,b<n$. By induction, both $a$ and $b$ are products of primes, so $n$ is too.
\end{proof}

\begin{proof}[Uniqueness]
Suppose
\[
p_1p_2\cdots p_r=q_1q_2\cdots q_s
\]
with all $p_i,q_j$ prime. Since $p_1$ divides the right side, Euclid's lemma says $p_1=q_j$ for some $j$. Cancel this common factor. Repeat. No factor can remain on only one side, so the lists agree up to order.
\end{proof}

\section{An exponent-parity proof of irrationality}

Unique factorization can rule out rational numbers without trying to compute
their decimal expansions.

\begin{theorem}
The number $\sqrt2$ is irrational.
\end{theorem}

\begin{proof}
Suppose instead that $\sqrt2=a/b$ for integers $a,b$ with $b\ne0$. Squaring
and clearing the denominator gives
\[
  a^2=2b^2.
\]
By the fundamental theorem of arithmetic, the exponent of $2$ in the prime
factorization of $a^2$ is even: every exponent occurring in $a$ has been
doubled. The exponent of $2$ in $2b^2$ is odd: it is one plus twice the
exponent occurring in $b$. Uniqueness says that equal positive integers must
have the same exponent of every prime, so this is impossible.
\end{proof}

\begin{center}
\begin{tikzpicture}[
  node distance=12mm,
  parity/.style={draw=navy!55,rounded corners=2mm,fill=white,minimum width=48mm,
    minimum height=13mm,align=center},
  arr/.style={-{Stealth[length=2.2mm]},thick,coral}
]
  \node[parity] (left) {$a^2$ has an \textbf{even}\\number of factors $2$};
  \node[parity,right=of left] (right) {$2b^2$ has an \textbf{odd}\\number of factors $2$};
  \draw[arr] (left)--node[above,font=\small\bfseries] {cannot be equal} (right);
\end{tikzpicture}

\small The contradiction is not merely that one side is even; it is that the
two sides contain incompatible numbers of factors $2$.
\end{center}

The same proof works with any prime in place of $2$. More generally, a rational
number can be a perfect square only when every prime exponent in its reduced
numerator and denominator is even.

\begin{labbox}[title={Aside: existence without a witness}]
Here is a famous one-line use of the law of excluded middle. Set
\[
  x=(\sqrt2)^{\sqrt2}.
\]
Either $x$ is rational or it is irrational. In the first case, the two
irrational numbers $a=b=\sqrt2$ satisfy $a^b=x\in\mathbb Q$. In the second
case, the two irrational numbers $a=x$ and $b=\sqrt2$ satisfy
\[
  a^b=\left((\sqrt2)^{\sqrt2}\right)^{\sqrt2}=2.
\]
Thus, classically, there exist irrational $a$ and $b$ for which $a^b$ is
rational.

The proof is valid in classical mathematics, but it is not a constructive
proof: it uses the disjunction ``$x$ is rational or $x$ is not rational''
without deciding which alternative holds, so it does not identify which of
the two proposed pairs is a witness. Constructive and intuitionistic
mathematics do not in general accept this unrestricted use of excluded
middle.

Unique factorization gives an elementary explicit witness. Take
\[
  a=\sqrt2,
  \qquad
  b=\log_{\sqrt2}3.
\]
Then $a^b=3$. Moreover $b$ is irrational: if $b=p/q$ with $q>0$, then
raising $(\sqrt2)^{p/q}=3$ to the power $2q$ would give
\[
  2^p=3^{2q},
\]
contradicting unique prime factorization. This time the proof supplies the
pair that works. Much more deeply, the Gelfond--Schneider theorem implies that
$(\sqrt2)^{\sqrt2}$ is transcendental, so it retrospectively selects the
second branch of the classical argument.
\end{labbox}

Collecting repeated primes gives a coordinate system
\[
n=2^{e_2}3^{e_3}5^{e_5}\cdots,
\]
where only finitely many exponents are nonzero. Multiplication adds exponent vectors; gcd takes coordinatewise minima; lcm takes coordinatewise maxima.

\begin{center}
\begin{tikzpicture}[x=1.2cm,y=1.0cm]
  \draw[-{Stealth},navy] (0,0)--(4.5,0) node[right] {$e_2$};
  \draw[-{Stealth},navy] (0,0)--(0,3.7) node[above] {$e_3$};
  \foreach \x in {0,...,4} \foreach \y in {0,...,3} \fill[navy!25] (\x,\y) circle (1.3pt);
  \fill[teal] (3,2) circle (4pt) node[above right] {$72=2^3 3^2$};
  \fill[coral] (2,3) circle (4pt) node[above right] {$108=2^2 3^3$};
  \fill[gold!80!black] (2,2) circle (4pt) node[below left] {$\gcd=36$};
  \draw[dashed,gold!80!black] (2,2)--(3,2);
  \draw[dashed,gold!80!black] (2,2)--(2,3);
  \fill[navy] (3,3) circle (4pt) node[above right] {$\operatorname{lcm}=216$};
\end{tikzpicture}

\small Prime exponents turn gcd and lcm into geometry.
\end{center}

\begin{proposition}
If $n=\prod p^{e_p}$ and $m=\prod p^{f_p}$, then
\[
\gcd(n,m)=\prod p^{\min(e_p,f_p)},\qquad
\operatorname{lcm}(n,m)=\prod p^{\max(e_p,f_p)}.
\]
In particular, $\gcd(n,m)\operatorname{lcm}(n,m)=nm$.
\end{proposition}

\section{The sieve of Eratosthenes}

Unique factorization also explains a practical way to find every prime up to a
chosen bound $N$. Write down the integers $2,3,\dots,N$. Circle the first
uncrossed number, then cross out its larger multiples. Repeat. The circled
numbers are precisely the primes.

\begin{center}
\newcommand{\sieveprime}[1]{\tikz[baseline=(n.base)]\node[
  draw=teal,fill=teal!13,rounded corners=1mm,minimum size=6.5mm,
  inner sep=0pt,font=\small\bfseries,text=teal!70!black] (n) {#1};}
\newcommand{\sievegone}[1]{\tikz[baseline=(n.base)]{
  \node[minimum size=6.5mm,inner sep=0pt,font=\small,text=muted] (n) {#1};
  \draw[coral!65,line width=.45pt] (n.south west)--(n.north east);}}
\newcommand{\sieveone}[1]{\tikz[baseline=(n.base)]\node[
  minimum size=6.5mm,inner sep=0pt,font=\small,text=muted!70] (n) {#1};}
\renewcommand{\arraystretch}{1.25}
\begin{tabular}{*{10}{c}}
\sieveone{1}&\sieveprime{2}&\sieveprime{3}&\sievegone{4}&\sieveprime{5}&\sievegone{6}&\sieveprime{7}&\sievegone{8}&\sievegone{9}&\sievegone{10}\\
\sieveprime{11}&\sievegone{12}&\sieveprime{13}&\sievegone{14}&\sievegone{15}&\sievegone{16}&\sieveprime{17}&\sievegone{18}&\sieveprime{19}&\sievegone{20}\\
\sievegone{21}&\sievegone{22}&\sieveprime{23}&\sievegone{24}&\sievegone{25}&\sievegone{26}&\sievegone{27}&\sievegone{28}&\sieveprime{29}&\sievegone{30}\\
\sieveprime{31}&\sievegone{32}&\sievegone{33}&\sievegone{34}&\sievegone{35}&\sievegone{36}&\sieveprime{37}&\sievegone{38}&\sievegone{39}&\sievegone{40}\\
\sieveprime{41}&\sievegone{42}&\sieveprime{43}&\sievegone{44}&\sievegone{45}&\sievegone{46}&\sieveprime{47}&\sievegone{48}&\sievegone{49}&\sievegone{50}
\end{tabular}

\small Sieve to $50$: the boxed survivors are prime; every composite has been crossed out.
\end{center}

\begin{proposition}[Why the sieve stops early]
If $n>1$ is composite, then it has a prime divisor at most $\sqrt n$.
Consequently, to find all primes up to $N$, it is enough to cross out multiples
of primes $p\le\sqrt N$.
\end{proposition}

\begin{proof}
Write $n=ab$ with $1<a\le b$. Then $a\le\sqrt n$. A prime factor $p$ of $a$
therefore satisfies $p\le a\le\sqrt n$. Thus if an integer $n\le N$ survives
all sieving primes through $\sqrt N$, it cannot be composite.
\end{proof}

\begin{checkpoint}
For the display above, only the multiples of $2,3,5,$ and $7$ need to be
crossed out. The next prime is $11>\sqrt{50}$. Notice also that one may begin
crossing out at $p^2$: every smaller multiple of $p$ has already been handled
by an earlier prime.
\end{checkpoint}

\section{There is no last prime}

\begin{theorem}[Euclid]
There are infinitely many primes.
\end{theorem}

\begin{proof}
If $p_1,\dots,p_k$ were all primes, consider $N=p_1\cdots p_k+1$. Some prime divides $N$, but none of the listed primes does: division by any $p_i$ leaves remainder $1$. Contradiction.
\end{proof}

This proof does not say that $N$ is prime. It says that $N$ has a prime factor missing from the proposed complete list.

\begin{exercisebox}
\begin{enumerate}
  \item Factor $7560$ and use the exponent vectors to find the number of its positive divisors. \exlevel{Fluency}
  \item Compute $\gcd(2^7 3^2 11,2^4 3^5 5)$ and the lcm without multiplying first. \exlevel{Fluency}
  \item Prove that if $\gcd(a,b)=1$ and $a\divides c$, $b\divides c$, then $ab\divides c$. \exlevel{Reasoning}
  \item Reproduce the exponent-parity proof that $\sqrt2$ is irrational, then prove that $\sqrt p$ is irrational for every prime $p$. State exactly where unique factorization is used. \exlevel{Reasoning}
  \item Prove that every composite $n>1$ has a prime divisor at most $\sqrt n$. Turn this into a primality test. \exlevel{Reasoning}
  \item Use the sieve of Eratosthenes to list every prime at most $100$. Explain why, after sieving by $2,3,5,$ and $7$, no further crossing-out step is needed. \exlevel{Fluency}
  \item Is $n!+1$ always prime? Find the first counterexample, then explain why Euclid's proof is unaffected. \exlevel{Exploration}
  \item Prove that infinitely many primes are congruent to $3$ modulo $4$ by adapting Euclid's idea to $4p_1\cdots p_k-1$. You may use the fact that a product of numbers congruent to $1$ modulo $4$ is again $1$ modulo $4$. \exlevel{Challenge}
  \item Verify directly that $a=\sqrt2$ and $b=\log_{\sqrt2}3$ are irrational and satisfy $a^b=3$. What does this explicit witness add to the excluded-middle proof? \exlevel{Reasoning}
\end{enumerate}
\end{exercisebox}

\chapter{Fractions inside fractions}
\chaptertag{Continued fractions and best approximations}

The quotients from Euclid's algorithm are not bookkeeping debris. Read them forward and they produce a new numerical language.

\section{Euclid, folded up}

From
\[
415=4\cdot93+43,\quad 93=2\cdot43+7,\quad
43=6\cdot7+1,\quad 7=7\cdot1,
\]
we obtain
\[
\frac{415}{93}=4+\cfrac1{2+\cfrac1{6+\cfrac1{7}}}=[4;2,6,7].
\]

\begin{definition}
A finite simple continued fraction is
\[
[a_0;a_1,\dots,a_n]
=a_0+\cfrac1{a_1+\cfrac1{\ddots+\cfrac1{a_n}}},
\]
where $a_0\in\Z$ and $a_i\in\N$ for $i\ge1$.
\end{definition}

The truncations $[a_0;\dots,a_k]=p_k/q_k$ are called \emph{convergents}. They can be computed without rebuilding the nested fraction:
\[
p_k=a_kp_{k-1}+p_{k-2},\qquad q_k=a_kq_{k-1}+q_{k-2},
\]
with $p_{-2}=0,p_{-1}=1,q_{-2}=1,q_{-1}=0$.

\section{Linear fractional transformations}
\label{sec:mobius}

A \emph{linear fractional transformation}, also called a M\"obius
transformation, is a function of the form
\[
 T_M(t)=\frac{at+b}{ct+d},
 \qquad M=\begin{pmatrix}a&b\\c&d\end{pmatrix},
 \qquad ad-bc\ne0.
\]
We will use integer or rational coefficients. The formula is defined when
$ct+d\ne0$. The determinant condition excludes a constant function and
ensures that the transformation can be undone.

Why does a $2\times2$ matrix act on a single number? Represent $t$ by the
column $(t,1)^T$, multiply, and take the ratio of the two coordinates:
\[
 \begin{pmatrix}a&b\\c&d\end{pmatrix}
 \begin{pmatrix}t\\1\end{pmatrix}
 =\begin{pmatrix}at+b\\ct+d\end{pmatrix}
 \quad\longmapsto\quad \frac{at+b}{ct+d}.
\]
Thus multiplying $M$ by any nonzero scalar does not change $T_M$. In this
chapter, $T_M(t)$ means this ratio, not ordinary multiplication of a matrix
by a scalar. One may also include a point $\infty$: when $c\ne0$,
$T_M(\infty)=a/c$ and $T_M(-d/c)=\infty$; when $c=0$, $\infty$ stays fixed.

\begin{proposition}[Composition is matrix multiplication]
For invertible $2\times2$ matrices $M,N$,
\[
 T_M\circ T_N=T_{MN},\qquad (T_M)^{-1}=T_{M^{-1}}.
\]
\end{proposition}
\begin{proof}
If $N=\begin{psmallmatrix}e&f\\g&h\end{psmallmatrix}$, substitution gives
\[
 T_M(T_N(t))
 =\frac{a(et+f)+b(gt+h)}{c(et+f)+d(gt+h)}
 =\frac{(ae+bg)t+(af+bh)}{(ce+dg)t+(cf+dh)}.
\]
These four coefficients are the entries of $MN$. Taking $N=M^{-1}$
proves the inverse rule. Explicitly, solving $y=(at+b)/(ct+d)$ gives
$t=(dy-b)/(a-cy)$.
\end{proof}

As with functions, the rightmost factor acts first. In particular,
\[
 D(a)=\begin{pmatrix}a&1\\1&0\end{pmatrix}
 \quad\text{acts by}\quad T_{D(a)}(t)=a+\frac1t.
\]
Therefore
\[
 [a_0;a_1,\ldots,a_k,t]=T_{D(a_0)D(a_1)\cdots D(a_k)}(t).
\]
Here the final entry $t$ is an entire remaining tail, not necessarily an
integer. The nesting of a continued fraction has become a product of
matrices.

\section{The matrix engine}

Package one continued-fraction digit into the matrix
\[
  D(a)=\begin{pmatrix}a&1\\1&0\end{pmatrix}.
\]
Repeated matrix multiplication stores every convergent at once:
\begin{equation}\label{eq:cf-matrix-product}
 D(a_0)D(a_1)\cdots D(a_k)
 =\begin{pmatrix}p_k&p_{k-1}\\q_k&q_{k-1}\end{pmatrix}.
\end{equation}
Indeed, multiplying the right-hand side by $D(a_{k+1})$ produces exactly the
two convergent recurrences. The first column advances to the next convergent,
while the old first column becomes the second.

Each digit matrix has determinant $-1$, not $+1$. Thus it lies in
$\operatorname{GL}_2(\Z)$, and a product with an even number of factors lies
in $\operatorname{SL}_2(\Z)$. More important than the sign is
\emph{unimodularity}: every product has determinant $\pm1$ and therefore has
an inverse with integer entries. Appendix~\ref{app:matrices} develops the
small amount of matrix algebra behind this observation.  Appendix~\ref{app:gosper}
then uses the same M\"obius machinery to perform exact arithmetic directly on
continued-fraction streams.

For example,
\[
 D(4)D(2)D(6)D(7)
 =\begin{pmatrix}415&58\\93&13\end{pmatrix},\qquad
 \begin{pmatrix}415&58\\93&13\end{pmatrix}^{-1}
 =\begin{pmatrix}13&-58\\-93&415\end{pmatrix}.
\]
The entries of the inverse already display the certificate
$1=13\cdot415-58\cdot93$. This is the first payoff of invertibility; the next
section explains why it is exactly the extended Euclidean algorithm.

\section{Extended Euclid in the convergents}

The convergents do not merely reuse Euclid's quotients: their numerators and
denominators are exactly the coefficients produced by the extended Euclidean
algorithm.

\begin{theorem}[Remainders from convergents]
Let $A>B>0$, and write the Euclidean divisions as
\[
 r_{-2}=A,\qquad r_{-1}=B,\qquad
 r_{k-2}=a_k r_{k-1}+r_k\quad(k=0,1,\dots,n),
\]
where $r_{n-1}=\gcd(A,B)$ and $r_n=0$. If
\[
  \frac AB=[a_0;a_1,\dots,a_n]
  \quad\text{and}\quad
  \frac{p_k}{q_k}=[a_0;a_1,\dots,a_k],
\]
then
\[
  \boxed{r_k=(-1)^k(q_kA-p_kB)}\qquad(0\le k\le n).
\]
In particular, the penultimate convergent gives a B\'ezout identity for the gcd.
\end{theorem}

\begin{proof}
For $k=0$, the formula reads $r_0=A-a_0B$, the first Euclidean division. For
$k=1$,
\[
  r_1=B-a_1r_0=-a_1A+(a_0a_1+1)B=-q_1A+p_1B.
\]
If the formula holds for the preceding two remainders, then
\begin{align*}
r_k&=r_{k-2}-a_kr_{k-1}\\
&=(-1)^k\bigl((a_kq_{k-1}+q_{k-2})A
 -(a_kp_{k-1}+p_{k-2})B\bigr)\\
&=(-1)^k(q_kA-p_kB),
\end{align*}
using the convergent recurrences. This proves the formula by induction.
\end{proof}

For $A/B=415/93$, the two algorithms can be read in the same table:
\begin{center}
\small
\begin{tabular}{cccc}
\toprule
$k$ & $a_k$ & convergent $p_k/q_k$ & signed linear combination $r_k$ \\
\midrule
$0$ & $4$ & $4/1$ & $415-4(93)=43$ \\
$1$ & $2$ & $9/2$ & $-(2(415)-9(93))=7$ \\
$2$ & $6$ & $58/13$ & $13(415)-58(93)=1$ \\
$3$ & $7$ & $415/93$ & $-(93(415)-415(93))=0$ \\
\bottomrule
\end{tabular}
\end{center}

Thus the convergent just before $415/93$ contains the certificate
\[
  1=13\cdot415-58\cdot93.
\]

\begin{center}
\begin{tikzpicture}[node distance=11mm and 8mm]
\tikzset{term/.style={draw=teal,rounded corners,fill=mist,minimum width=12mm,minimum height=8mm,font=\small},
arrow/.style={-{Stealth},navy!65,thick}}
\node[term] (a0) {$4$};
\node[term,right=of a0] (a1) {$2$};
\node[term,right=of a1] (a2) {$6$};
\node[term,right=of a2] (a3) {$7$};
\node[below=13mm of a0] (c0) {$4/1$};
\node[below=13mm of a1] (c1) {$9/2$};
\node[below=13mm of a2] (c2) {$58/13$};
\node[below=13mm of a3] (c3) {$415/93$};
\draw[arrow] (a0)--(c0); \draw[arrow] (a1)--(c1);
\draw[arrow] (a2)--(c2); \draw[arrow] (a3)--(c3);
\draw[teal!55,thick] (c0)--(c1)--(c2)--(c3);
\node[muted,below=6mm of c2] {Each convergent remembers all earlier partial quotients.};
\end{tikzpicture}
\end{center}

\begin{theorem}[Determinant identity]
Consecutive convergents satisfy
\[
p_kq_{k-1}-p_{k-1}q_k=(-1)^{k-1}.
\]
Consequently, each $p_k/q_k$ is in lowest terms and
\[
\abs{\frac{p_k}{q_k}-\frac{p_{k-1}}{q_{k-1}}}=\frac1{q_kq_{k-1}}.
\]
\end{theorem}

\begin{proof}
Take determinants in \eqref{eq:cf-matrix-product}. The determinant on the
left is $(-1)^{k+1}$, while the determinant on the right is
$p_kq_{k-1}-p_{k-1}q_k$. Since $(-1)^{k+1}=(-1)^{k-1}$, the identity follows.
\end{proof}

\section{Infinite continued fractions}

For an irrational real number $x$, repeatedly take integer parts:
\[
a_0=\floor{x},\qquad x_1=\frac1{x-a_0},\qquad a_1=\floor{x_1},\quad\dots
\]
This produces $x=[a_0;a_1,a_2,\dots]$.

\begin{example}[The silver staircase]
Let $x=\sqrt2$. Since
\[
\frac1{\sqrt2-1}=\sqrt2+1=2+(\sqrt2-1),
\]
the same remainder repeats, so
\[
\sqrt2=[1;\overline{2}]=[1;2,2,2,\dots].
\]
Its first convergents are
\[
1,\quad \frac32,\quad\frac75,\quad\frac{17}{12},\quad\frac{41}{29},\dots
\]
They alternately fall below and above $\sqrt2$.
\end{example}

\begin{theorem}[The alternating squeeze]
Let $c_k=p_k/q_k$ be the zero-indexed convergents of an irrational number
$x=[a_0;a_1,a_2,\dots]$. Then
\[
  c_0<c_2<c_4<\cdots<x<\cdots<c_5<c_3<c_1.
\]
Thus, when the convergents are counted starting with the first, the odd-numbered
ones increase and stay below $x$, while the even-numbered ones decrease and
stay above $x$.
\end{theorem}

\begin{proof}[The M\"obius derivative trick]
Fix $k$ and leave the tail variable. The continued fraction with prefix through
$a_k$ is the fractional linear, or M\"obius, map
\[
  M_k(t)=[a_0;a_1,\dots,a_k,t]
  =\frac{p_kt+p_{k-1}}{q_kt+q_{k-1}}.
\]
The determinant identity gives
\[
  M_k'(t)=
  \frac{p_kq_{k-1}-p_{k-1}q_k}{(q_kt+q_{k-1})^2}
  =\frac{(-1)^{k+1}}{(q_kt+q_{k-1})^2}.
\]
The complete tail $\alpha_{k+1}=[a_{k+1};a_{k+2},\dots]$ is finite and positive,
and
\[
  x=M_k(\alpha_{k+1}),\qquad c_k=M_k(\infty).
\]
When $k$ is even, $M_k$ is decreasing, so $c_k<x$; when $k$ is odd,
$M_k$ is increasing, so $x<c_k$. The same comparison with
$c_{k+2}=M_k([a_{k+1};a_{k+2}])$ shows that $c_k<c_{k+2}$ for even $k$ and
$c_{k+2}<c_k$ for odd $k$. These are exactly the two monotone chains displayed
above.
\end{proof}

\begin{bigidea}
Every new convergent switches sides, but every second convergent moves inward.
Continued fractions therefore trap $x$ between a rising sequence of lower
bounds and a falling sequence of upper bounds.
\end{bigidea}

\begin{center}
\begin{tikzpicture}[x=.95cm,y=1cm]
  \draw[-{Stealth},navy] (-.3,0)--(8.05,0);
  \foreach \x/\lab/\col in {0/1/teal,1.3/{7/5}/teal,2.6/{41/29}/teal,
    3.8/{\sqrt2}/gold,5.0/{99/70}/coral,6.3/{17/12}/coral,7.6/{3/2}/coral}{
    \draw[\col,thick] (\x,-.11)--(\x,.11);
    \node[\col!80!black,above=2pt] at (\x,.11) {$\lab$};
  }
  \draw[-{Stealth},teal,thick] (0,-.42)--(1.18,-.42);
  \draw[-{Stealth},teal,thick] (1.3,-.42)--(2.48,-.42);
  \draw[-{Stealth},coral,thick] (7.6,-.42)--(6.42,-.42);
  \draw[-{Stealth},coral,thick] (6.3,-.42)--(5.12,-.42);
\end{tikzpicture}

\small Schematic, not to scale: both parity subsequences move inward toward $\sqrt2$.
\end{center}

\begin{theorem}[Approximation by convergents]
If $p_k/q_k$ is a convergent of an irrational $x$, then
\[
\abs{x-\frac{p_k}{q_k}}<\frac1{q_kq_{k+1}}<\frac1{q_k^2}.
\]
Moreover, if $0<q<q_{k+1}$, then no fraction $p/q$ is closer to $x$ than both neighboring convergents in the appropriate interval. In this precise sense, convergents are best approximations with small denominators.
\end{theorem}

The determinant identity is the engine: distinct reduced fractions $a/b,c/d$ are at least $1/(bd)$ apart. There is simply not enough room for another small-denominator fraction between a convergent and $x$.

\begin{theorem}[Legendre's approximation criterion]
Let $x$ be irrational and let $p/q$ be in lowest terms with $q>0$. If
\[
  \abs{x-\frac pq}<\frac{1}{2q^2},
\]
then $p/q$ is a convergent of $x$.
\end{theorem}

The theorem is a partial converse to the approximation theorem: an
\emph{exceptionally} good rational approximation cannot have arisen by
accident. A standard proof chooses the two consecutive convergents whose
denominators surround $q$ and writes the integer vector $(p,q)$ in their
unimodular basis. If $(p,q)$ is not one of the basis vectors, the determinant
$\pm1$ forces it to be at least $1/(2q^2)$ away from $x$. This is another place
where the integral inverse of the convergent matrix does real work.

\begin{labbox}[title={Aside: continued fractions keep the calendar},breakable=false]
A calendar must use whole days. Write the tropical year as $365+\alpha$ days,
where $\alpha\approx0.2421988$. A $q$-year cycle with $p$ leap days therefore
asks for $p/q\approx\alpha$. Now
\[
 \alpha=[0;4,7,1,3,5,\dots],
 \qquad\text{with convergents}\qquad
 \frac14,\ \frac7{29},\ \frac8{33},\ \frac{31}{128},\
 \frac{163}{673},\dots
\]
The first convergent is the Julian rule: one leap day every four years. The
convergent $31/128=0.2421875$ is remarkably accurate, but a $128$-year civil
cycle is awkward.

The Gregorian rule adds a practical constraint: organize corrections by
centuries. Writing $q=100Q$ changes the approximation problem:
\[
 q=100Q:\qquad \frac pQ\approx100\alpha\approx24.21988,
 \qquad [24;4]=\frac{97}{4}.
\]
Thus the first nonintegral convergent of $100\alpha$ gives $97$ leap days in
$400$ years and the mean length
\[
  365+\frac{97}{400}=365.2425.
\]
Operationally, the rule counts $100$ multiples of $4$, removes the four
century years, and restores the one multiple of $400$: $100-4+1=97$.

Notice that $97/400$ is not a convergent of $\alpha$ itself; it emerges after
the century constraint. This is a mathematical reconstruction, not a claim
about the reformers' historical method. It shows that ``best'' depends on which
denominators an application permits.
\end{labbox}

\begin{labbox}[title={Aside: continued fractions are roughly ``base 11''}]
Khinchin's theorem says that for almost every real number,
\[
  (a_1a_2\cdots a_n)^{1/n}\longrightarrow
  K_0\approx2.685452,
\]
the universal \emph{Khinchin constant}. For efficiency, L\'evy's companion
theorem is more direct:
\[
 q_n^{1/n}\longrightarrow
 L=e^{\pi^2/(12\log2)}\approx3.275823.
\]
Since convergent error is on the scale $q_n^{-2}$, it follows that
\[
 \abs{x-p_n/q_n}^{-1/n}\longrightarrow
 L^2=e^{\pi^2/(6\log2)}\approx10.731.
\]
That is the ``base 11'' joke: one generic continued-fraction term buys about
$\log_{10}(10.731)\approx1.0306$ decimal digits. Lochs' theorem makes the
comparison exact: $N$ known decimal digits determine about
$(6\log2\log10/\pi^2)N\approx0.97027N$ continued-fraction terms.

``Almost every'' matters: a partial quotient is not literally a fixed-base
digit, and quadratic surds such as $\sqrt2$ have exceptional periodic growth.
\end{labbox}

\begin{checkpoint}[title={Experiment: $\pi$ looks typical; $e$ does not},breakable=false]
Write
\[
 G_n(x)=(a_1a_2\cdots a_n)^{1/n},
 \qquad B_n(x)=q_n^{2/n}.
\]
Khinchin and L\'evy predict, for almost every $x$, that
$G_n(x)\to2.685452\dots$ and $B_n(x)\to10.731016\dots$. For
\[
 \pi=[3;7,15,1,292,1,1,1,2,1,3,\dots],
\]
a computation gives the following rounded values:
\begin{center}
\small
\begin{tabular}{rccc}
\toprule
$n$ & $G_n(\pi)$ & $B_n(\pi)$ & $\log_{10}B_n(\pi)$ \\
\midrule
$100$     & $2.68315$ & $10.68768$ & $1.02888$ \\
$1{,}000$ & $2.66270$ & $10.54434$ & $1.02302$ \\
$10{,}000$ & $2.66320$ & $10.57847$ & $1.02442$ \\
$100{,}000$ & $2.67956$ & $10.67835$ & $1.02850$ \\
\midrule
almost-everywhere target & $2.68545$ & $10.73102$ & $1.03064$ \\
\bottomrule
\end{tabular}
\end{center}
The wobble is real: rare large partial quotients, such as the early $292$,
produce visible jumps. A published computation using $17{,}001{,}303$ terms
found $G_n(\pi)=2.686393$. This is compelling evidence, but no theorem is known
that $\pi$ satisfies either limiting law.

The number $e$ is a useful contrast, not a second generic-looking example:
\[
 e=[2;1,2,1,1,4,1,1,6,1,1,8,\dots].
\]
For $n=3m$, its first $n$ noninitial partial quotients have product $2^m m!$.
Stirling's formula therefore gives
\[
 G_{3m}(e)=(2^m m!)^{1/(3m)}\sim(2m/e)^{1/3}\longrightarrow\infty.
\]
Also $q_{3m}\geq a_1\cdots a_{3m}=2^m m!$, so $B_{3m}(e)\to\infty$ as well.
Thus $e$ is provably exceptional for both laws, while $\pi$ is numerically
persuasive but theoretically mysterious.
\end{checkpoint}

\begin{warningbox}[breakable=false]
Legendre, Lagrange, and Galois all appear in this subject. Legendre's theorem
recognizes very good rational approximations. Lagrange's theorem characterizes
eventually periodic continued fractions of quadratic irrationals; Galois's
refinement says exactly when the period begins immediately.
\end{warningbox}

\section{A parallel periodicity: digits in base \texorpdfstring{$b$}{b}}

Continued fractions are not the first place students meet repeating expansions.
Let $0<m<q$ and write $x=m/q$ in base $b\ge2$. Long division produces digits
$d_1,d_2,\dots$ and remainders $r_0,r_1,\dots$ by
\[
  r_0=m,\qquad br_n=qd_{n+1}+r_{n+1},\qquad 0\le r_{n+1}<q.
\]
There are only $q$ possible remainders. If a remainder becomes $0$, the
expansion terminates; otherwise some remainder repeats, and from that point
the deterministic long-division process repeats. Thus every rational number
has an eventually periodic base-$b$ expansion.

The converse uses only a geometric series. If a block of $\ell$ digits with
integer value $N$ repeats immediately, then
\[
  0.\overline{d_1\cdots d_\ell}_b
  =\frac{N}{b^\ell}+\frac{N}{b^{2\ell}}+\cdots
  =\frac{N}{b^\ell-1}.
\]
An initial nonrepeating block only adds a rational number with denominator a
power of $b$.

\begin{theorem}[When the repeating block begins immediately]
If the reduced fraction $m/q$ has a purely periodic base-$b$ expansion, then
$\gcd(q,b)=1$. Conversely, if $\gcd(q,b)=1$, its expansion is purely periodic.
\end{theorem}

\begin{proof}[The direction available now]
A purely periodic expansion has the form $N/(b^\ell-1)$. Its reduced
denominator divides $b^\ell-1$, which is relatively prime to $b$.
\end{proof}

The reverse implication is deliberately postponed until Chapter~6. Euler's
theorem will give
\[
  b^{\varphi(q)}\congr1\pmod q,
\]
so $q\divides b^{\varphi(q)}-1$; rewriting $m/q$ with that denominator creates
a repeating block (possibly with leading zeros). In decimal, for example,
\[
  \frac17=0.\overline{142857},\qquad
  \frac18=0.125000\dots.
\]
The first expansion is purely periodic because $\gcd(7,10)=1$; the second is
only eventually periodic.

\section{Periodic tails are fixed points}
\label{sec:periodic-fixed-points}

The matrix language answers a question that the nested notation conceals:
what kind of number can a repeating continued fraction represent?
Throughout this discussion the continued fractions are \emph{simple}:
$a_0\in\Z$ and $a_n\ge1$ for $n\ge1$.

\begin{definition}
An infinite continued fraction is \emph{eventually periodic} if its digits
repeat after a finite prefix. It is \emph{purely periodic} if repetition
starts at the first digit:
\[
 [\overline{b_0,\ldots,b_{\ell-1}}]
 =[b_0;b_1,\ldots,b_{\ell-1},b_0,b_1,\ldots].
\]
In a purely periodic simple continued fraction, every $b_i$ is positive,
including $b_0$. Rational numbers have finite simple continued fractions;
we do not call those eventually periodic here.
\end{definition}

\begin{definition}
A \emph{quadratic irrational} (or quadratic surd) is an irrational real
root of a quadratic polynomial with integer coefficients. Its other root
is its \emph{conjugate}, written $x'$. Equivalently,
\[
 x=u+v\sqrt D,\qquad x'=u-v\sqrt D
 \quad (u,v\in\Q,\ v\ne0),
\]
where $D$ is a positive nonsquare integer. Replacing $\sqrt D$ by
$-\sqrt D$ respects addition, multiplication, and division by nonzero
numbers. In particular, $(T_M(x))'=T_M(x')$ for rational $M$.
\end{definition}

\begin{lemma}[A nontrivial fixed point is quadratic]
\label{lem:mobius-fixed-point}
If an invertible rational matrix
$M=\begin{psmallmatrix}a&b\\c&d\end{psmallmatrix}$ is not a scalar
multiple of the identity and fixes an irrational number $t$ under its
linear fractional action, then $t$ is a quadratic irrational.
\end{lemma}
\begin{proof}
Multiplying $t=(at+b)/(ct+d)$ by its denominator gives
\begin{equation}\label{eq:mobius-fixed-point}
 ct^2+(d-a)t-b=0.
\end{equation}
The coefficients are not all zero: that would say $c=b=0$ and $d=a$,
making $M$ scalar. Moreover, $c$ cannot be zero, since a nonzero linear
polynomial with rational coefficients cannot vanish at an irrational
number. Thus this is a genuine quadratic equation. Clear denominators to
obtain integer coefficients.
\end{proof}

\begin{warningbox}
The qualification \emph{nontrivial} matters. The identity transformation
fixes every real number, so $t=T_I(t)$ by itself says nothing about $t$.
For a period matrix there is no such problem: its lower-left entry is
strictly positive.
\end{warningbox}

\begin{theorem}[Lagrange]
\label{thm:lagrange-periodicity}
A real number has an eventually periodic simple continued fraction if and
only if it is a quadratic irrational.
\end{theorem}

\begin{proof}[Proof that periodicity implies a quadratic equation]
First let $x=[\overline{b_0,\ldots,b_{\ell-1}}]$. Removing one complete
period leaves the same infinite tail, so
\[
 x=T_N(x),\qquad
 N=D(b_0)\cdots D(b_{\ell-1})
   =\begin{pmatrix}A&B\\C&D_0\end{pmatrix}.
\]
Because every $b_i\ge1$, we have $B,C>0$. Also
$\det N=(-1)^\ell$. The fixed-point equation is
\[
 Cx^2+(D_0-A)x-B=0.
\]
The alternating-convergent argument also applies to an infinite
positive digit sequence: it gives a limit and an error at most
$1/(q_kq_{k+1})$. That limit cannot be rational. Indeed, the
convergents are in lowest terms and $q_k\to\infty$, so they are
eventually distinct from any fixed rational $p/q$. Their distance
from $p/q$ is then at least $1/(q q_k)$, contradicting the error
bound once $q_{k+1}>q$. Thus $x$ is quadratic.

Now allow a prefix:
\[
 y=[a_0;\ldots,a_{r-1},\overline{b_0,\ldots,b_{\ell-1}}]
   =T_M(x),\qquad M=D(a_0)\cdots D(a_{r-1}).
\]
If the prefix is empty, take $M=I$. Since $T_M$ is invertible,
substituting $x=T_{M^{-1}}(y)$ into $x=T_N(x)$ gives
\begin{equation}\label{eq:prefix-conjugation}
 \boxed{\ y=T_{MNM^{-1}}(y).\ }
\end{equation}
Both $M$ and $N$ have determinant $\pm1$, so $K=MNM^{-1}$ is an
integer matrix. It is not scalar: if $K=\lambda I$, then
$N=M^{-1}KM=\lambda I$, contrary to $C>0$.

Finally, write $M=\begin{psmallmatrix}a&b\\c&d\end{psmallmatrix}$.
The number $y$ is irrational. Otherwise the inverse formula
$x=T_{M^{-1}}(y)$ would make $x$ rational. Its denominator cannot vanish
at $y=T_M(x)$: direct substitution gives
$a-cy=(ad-bc)/(cx+d)\ne0$.
The fixed-point lemma applied to $K$ therefore gives a genuine quadratic
equation for $y$ as well.
\end{proof}

\begin{bigidea}
A repeating tail is a fixed point. Attaching a finite prefix
\emph{conjugates the transformation that fixes it}. This is the missing
algebraic step between a purely periodic tail and an eventually periodic
number; it requires only matrix multiplication and inversion.
\end{bigidea}

\begin{example}[A prefix changes the equation, not its degree]
Take
\[
 x=[\overline{1,2}],\qquad
 N=D(1)D(2)=\begin{pmatrix}3&1\\2&1\end{pmatrix}.
\]
Then $2x^2-2x-1=0$, so $x=(1+\sqrt3)/2$. Prefix the digit $1$:
\[
 y=[1;\overline{1,2}]=1+\frac1x,\qquad
 M=D(1)=\begin{pmatrix}1&1\\1&0\end{pmatrix}.
\]
The conjugated matrix is
\[
 MNM^{-1}
 =\begin{pmatrix}1&1\\1&0\end{pmatrix}
  \begin{pmatrix}3&1\\2&1\end{pmatrix}
  \begin{pmatrix}0&1\\1&-1\end{pmatrix}
 =\begin{pmatrix}2&3\\1&2\end{pmatrix}.
\]
Thus $y=(2y+3)/(y+2)$, giving $y^2=3$. Since $y>0$, $y=\sqrt3$.
The same matrix will reappear as a generator of Pell solutions.
\end{example}

\section{Why every quadratic irrational repeats}
\label{sec:quadratic-finite-states}

We now prove the other direction of Lagrange's theorem. There are two
tasks: find a bounded set of integer states, and show that the continued
fraction eventually enters it. The conjugate supplies exactly the
inequalities we need.

\begin{definition}
A quadratic irrational $x$ is \emph{reduced} if
\[
 x>1\qquad\text{and}\qquad -1<x'<0.
\]
\end{definition}

\subsection*{Step 1: every quadratic irrational has a reduced tail}
Write $x_0=x$, $a_n=\floor{x_n}$, and
$x_{n+1}=1/(x_n-a_n)$. For $n\ge1$, $x_n>1$ and $a_n\ge1$.
Conjugating the recurrence gives
\[
 x_{n+1}'=\frac1{x_n'-a_n}.
\]
Consequently, once $x_n'<0$ with $n\ge1$, the next conjugate belongs to
$(-1,0)$. Every later complete quotient is reduced too.

Why must a negative conjugate appear? Invert the matrix for the first
$n$ digits and then conjugate. For $n\ge2$ this gives
\[
 x_n'=-\frac{q_{n-2}x'-p_{n-2}}{q_{n-1}x'-p_{n-1}}.
\]
The convergents $p_k/q_k$ tend to $x$, and $x'\ne x$. Hence, for all
sufficiently large $k$, the numbers
$x'-p_k/q_k$ have the same sign as $x'-x$.
The numerator and denominator in the displayed ratio therefore have
the same sign for large $n$, since the $q_k$ are positive.
It follows that $x_n'<0$ eventually, as required.
This argument also applies when the original $x$ is negative or less
than $1$.

\subsection*{Step 2: integer states evolve by an exact recurrence}
Choose an equation $Ax^2+Bx+C=0$ with integer coefficients and $A\ne0$.
After changing all three signs if necessary, the quadratic formula reads
\[
 x=\frac{-B+\sqrt D}{2A},
 \qquad D=B^2-4AC>0.
\]
Here $D$ is a nonsquare integer; it need not be squarefree. Put
$P_0=-B$ and $Q_0=2A$. Then $Q_0\ne0$ and
$Q_0\divides D-P_0^2$, because $D-P_0^2=-4AC$.
The sign change above is important: $Q_0$ may initially be negative.

Suppose $x_n=(P_n+\sqrt D)/Q_n$ with integers $P_n,Q_n$ satisfying
$Q_n\divides D-P_n^2$. Rationalizing $1/(x_n-a_n)$ yields
\begin{equation}\label{eq:quadratic-state-recurrence}
 \boxed{\quad
 P_{n+1}=a_nQ_n-P_n,\qquad
 Q_{n+1}=\frac{D-P_{n+1}^2}{Q_n}.
 \quad}
\end{equation}
Indeed,
\[
 \frac1{x_n-a_n}
 =\frac{Q_n}{\sqrt D-P_{n+1}}
 =\frac{P_{n+1}+\sqrt D}{(D-P_{n+1}^2)/Q_n}.
\]
Since $P_{n+1}\congr-P_n\pmod{Q_n}$, the new $Q_{n+1}$ is an integer.
It is nonzero because $D$ is not a square. The equality
$D-P_{n+1}^2=Q_nQ_{n+1}$ supplies the divisibility condition for the
next step. Thus the representation persists forever with the same $D$.

\subsection*{Step 3: reduced states lie in a finite box}
For a reduced complete quotient,
\[
 x_n=\frac{P_n+\sqrt D}{Q_n}>1,\qquad
 -1<\frac{P_n-\sqrt D}{Q_n}<0.
\]
Their difference is $2\sqrt D/Q_n>1$, so
$0<Q_n<2\sqrt D$. The conjugate inequalities then give
$\sqrt D-Q_n<P_n<\sqrt D$. In particular,
\begin{equation}\label{eq:reduced-state-bounds}
 \boxed{\quad |P_n|<\sqrt D,\qquad 0<Q_n<2\sqrt D.\quad}
\end{equation}
Only finitely many integer pairs satisfy these bounds.
By Step 1 all sufficiently late states lie in this box, and by Step 2
each state determines its next state, including the digit
$a_n=\floor{(P_n+\sqrt D)/Q_n}$.
Two states must repeat. Every later state, and hence every later digit,
then repeats with the same period. This completes the proof of
Theorem~\ref{thm:lagrange-periodicity}.

\begin{example}[The states for $\sqrt3$]
Starting with $(P_0,Q_0)=(0,1)$ and $D=3$ gives
\begin{center}
\begin{tabular}{ccccc}
\toprule
$n$ & $P_n$ & $Q_n$ & $x_n$ & $a_n$\\
\midrule
$0$ & $0$ & $1$ & $\sqrt3$ & $1$\\
$1$ & $1$ & $2$ & $(1+\sqrt3)/2$ & $1$\\
$2$ & $1$ & $1$ & $1+\sqrt3$ & $2$\\
$3$ & $1$ & $2$ & $(1+\sqrt3)/2$ & $1$\\
\bottomrule
\end{tabular}
\end{center}
The initial state is not reduced. The next two are reduced and form a
cycle, giving $\sqrt3=[1;\overline{1,2}]$.
\end{example}

\section{Galois: when repetition starts immediately}
\label{sec:galois}

The fixed-point equation tells us that periodic continued fractions are
quadratic. It does \emph{not} tell us which quadratic irrationals have
no initial nonrepeating segment. That distinction is Galois's theorem.

\begin{theorem}[Galois's pure-periodicity criterion]
\label{thm:galois}
The simple continued fraction of a quadratic irrational $x$ is purely
periodic if and only if $x$ is reduced:
\[
 x>1\qquad\text{and}\qquad -1<x'<0.
\]
\end{theorem}
\begin{proof}
Suppose first that $x$ is purely periodic with period length $\ell$.
The fixed-point argument shows that it is quadratic, and Step 1 above
shows that all sufficiently late complete quotients are reduced.
But $x_{k\ell}=x_0=x$ for every $k$. Taking $k$ large proves that
$x$ itself is reduced.

Conversely, suppose $x$ is reduced. All its complete quotients remain
reduced, and the finite-state proof gives $x_i=x_j$ for some $i<j$.
To conclude that the period starts at $0$, we need one more fact:
\emph{a reduced complete quotient has at most one reduced predecessor}.

Indeed, suppose the continued-fraction step sends a reduced $u$ to $z$,
and let $a=\floor u$. Then $z=1/(u-a)$, so
\[
 -\frac1{z'}=a-u'.
\]
Since $-1<u'<0$, we have $a<-1/z'<a+1$. Therefore
\begin{equation}\label{eq:reduced-predecessor}
 \boxed{\quad
 a=\floor{-1/z'},\qquad u=a+\frac1z.
 \quad}
\end{equation}
Both the previous digit and the previous reduced number are determined
by $z$. If $i>0$, equality $x_i=x_j$ therefore forces
$x_{i-1}=x_{j-1}$. Repeating the argument gives $x_0=x_{j-i}$.
The deterministic continued-fraction algorithm now repeats from its
first step, proving pure periodicity.
\end{proof}

\begin{bigidea}
Finitely many states give \emph{eventual} periodicity.
Recovering the previous reduced state removes the possible initial
segment and gives \emph{pure} periodicity. These are different steps.
For instance, the finite map $0\mapsto1$, $1\mapsto1$ has an initial
segment when started at $0$; finiteness alone does not exclude one.
\end{bigidea}

\begin{example}[Two inequalities, both necessary]
The number $\sqrt2=[1;\overline2]$ is not reduced, because
$-\sqrt2<-1$. It is eventually but not purely periodic.
In contrast,
\[
 1+\sqrt2=[\overline2],
 \qquad \frac{1+\sqrt5}{2}=[\overline1]
\]
are reduced and purely periodic.
The number $\sqrt2/2$ has conjugate in $(-1,0)$ but is less than $1$,
so it is not purely periodic either.

A displayed prefix does not necessarily represent a genuine preperiod:
\[
 [2;\overline{1,2}]=1+\sqrt3=[\overline{2,1}].
\]
Here the conjugate $1-\sqrt3$ lies in $(-1,0)$, so Galois correctly
detects pure periodicity despite the first way of writing the digits.
\end{example}

\begin{corollary}[The conjugate reverses the period]
If $x=[\overline{b_0,\ldots,b_{\ell-1}}]$, then
\[
 -\frac1{x'}=[\overline{b_{\ell-1},\ldots,b_0}].
\]
\end{corollary}
\begin{proof}
Each $D(b_i)$ is symmetric, so reversing the digits transposes their
product $N=\begin{psmallmatrix}A&B\\C&D_0\end{psmallmatrix}$.
Its fixed-point equation is $Cx^2+(D_0-A)x-B=0$. Substituting
$x'=-1/z$ gives
\[
 Bz^2+(D_0-A)z-C=0,
\]
which is the fixed-point equation for $N^T$. Since $B,C>0$, this
quadratic has exactly one positive root. Both $-1/x'$ and the
reversed periodic continued fraction are that root.
\end{proof}

\begin{checkpoint}[title={Three names, three different conclusions}]
\textbf{Legendre}: an exceptionally close rational approximation must
be a convergent.
\quad\textbf{Lagrange}: the numbers with eventually periodic simple
continued fractions are exactly the quadratic irrationals.
\quad\textbf{Galois}: among them, pure periodicity means exactly
$x>1$ and $-1<x'<0$.
The first direction of Lagrange is the fixed-point calculation;
Galois's additional content is the precise conjugate criterion.
\end{checkpoint}


\section{Periodicity meets Pell's equation}
\label{sec:pell}

Fix a nonsquare integer $D\ge2$. \emph{Pell's equation} asks for positive
integers $p,q$ such that
\[
 p^2-Dq^2=1.
\]
We will also study the equation with right-hand side $-1$. There are two
questions to settle: why must a solution appear among the convergents of
$\sqrt D$, and which convergents actually work? An approximation estimate
answers the first; an exact identity answers the second.

\subsection*{Every Pell solution is a convergent}
Suppose $p^2-Dq^2=\pm1$, with $p,q>0$. Any common divisor of $p,q$ would
divide $1$, so $p/q$ is reduced. Factoring a difference of squares gives
\begin{equation}\label{eq:pell-error}
 \abs{\sqrt D-\frac pq}
 =\frac{1}{q^2(\sqrt D+p/q)}
 <\frac{1}{2q^2}.
\end{equation}
For the strict inequality, note that
$(p/q)^2\ge D-1\ge1$, while $\sqrt D>1$.
Legendre's criterion now makes $p/q$ a convergent of $\sqrt D$.
This does not yet prove that a solution exists: we still need a convergent
whose difference of squares is exactly $1$ or $-1$.

\subsection*{The palindrome in a square-root period}
Write $a_0=\floor{\sqrt D}$. The first complete quotient is
\[
 x_1=\frac1{\sqrt D-a_0}>1,
 \qquad x_1'=-\frac1{\sqrt D+a_0}\in(-1,0).
\]
Galois's theorem therefore makes $x_1$ purely periodic. Let $L$ be its
\emph{least} period length, so that
$\sqrt D=[a_0;\overline{a_1,\ldots,a_L}]$.

\begin{theorem}[The square-root palindrome]
\label{thm:sqrt-palindrome}
The last digit of the least period is $a_L=2a_0$, and
\[
 a_j=a_{L-j}\qquad(1\le j\le L-1).
\]
Thus
\[
 \boxed{\quad
 \sqrt D=[a_0;\overline{a_1,\ldots,a_{L-1},2a_0}],
 \quad (a_1,\ldots,a_{L-1})\text{ is a palindrome}.
 \quad}
\]
For $L=1$ the palindrome is the empty block.
\end{theorem}
\begin{proof}
Apply Corollary~\ref{cor:period-reversal} to $x_1$:
\[
 -\frac1{x_1'}=\sqrt D+a_0
   =[\overline{a_L,a_{L-1},\ldots,a_1}].
\]
Its integer part is $2a_0$, so $a_L=2a_0$. On the other hand,
\[
 \sqrt D+a_0=2a_0+\frac1{x_1}
   =[2a_0;\overline{a_1,\ldots,a_L}].
\]
These are two expansions of the same irrational number. The
continued-fraction algorithm determines its digits uniquely. Comparing
the next $L-1$ digits gives
$(a_1,\ldots,a_{L-1})=(a_{L-1},\ldots,a_1)$.
\end{proof}

For example,
\[
 \sqrt{23}=[4;\overline{1,3,1,8}],\qquad
 \sqrt{14}=[3;\overline{1,2,1,6}],\qquad
 \sqrt2=[1;\overline2].
\]
The final $2a_0$ marks the end of the symmetric block; it is not itself
part of the palindrome asserted by the theorem.

There is a useful consequence for the integer states
$x_n=(P_n+\sqrt D)/Q_n$ of
\eqref{eq:quadratic-state-recurrence}, starting with $(P_0,Q_0)=(0,1)$.
Since $x_{L+1}=x_1$,
\[
 x_L=a_L+\frac1{x_{L+1}}=2a_0+(\sqrt D-a_0)=\sqrt D+a_0.
\]
Consequently $(P_L,Q_L)=(a_0,1)$, and the same state occurs at every
positive multiple of $L$. Conversely, if $n\ge1$ and $Q_n=1$, reducedness
gives $-1<P_n-\sqrt D<0$, forcing $P_n=a_0$. Thus $x_n=x_L$.
The $L$ states in a least cycle are distinct, so
\begin{equation}\label{eq:pell-state-return}
 \boxed{\qquad Q_n=1\quad\Longleftrightarrow\quad L\divides n
 \qquad(n\ge1).\qquad}
\end{equation}

\subsection*{An exact identity for every convergent}
\begin{lemma}[The complete quotient measures the Pell remainder]
\label{lem:pell-remainder}
For every $n\ge1$,
\begin{equation}\label{eq:pell-remainder}
 \boxed{\quad p_{n-1}^2-Dq_{n-1}^2=(-1)^nQ_n.\quad}
\end{equation}
\end{lemma}
\begin{proof}
Keep the entire tail in the matrix product:
\[
 \sqrt D
 =\frac{p_{n-1}x_n+p_{n-2}}{q_{n-1}x_n+q_{n-2}},
 \qquad x_n=\frac{P_n+\sqrt D}{Q_n}.
\]
Substitute the second expression into the first and multiply through.
Because $\sqrt D$ is irrational, the rational coefficients and the
coefficients of $\sqrt D$ must agree separately. They give
\begin{align*}
 p_{n-1}&=P_nq_{n-1}+Q_nq_{n-2},\\
 Dq_{n-1}&=P_np_{n-1}+Q_np_{n-2}.
\end{align*}
Multiply the first equality by $p_{n-1}$ and the second by $q_{n-1}$,
then subtract:
\begin{align*}
 p_{n-1}^2-Dq_{n-1}^2
 &=Q_n(p_{n-1}q_{n-2}-p_{n-2}q_{n-1})\\
 &=(-1)^nQ_n.
\end{align*}
The last step is the determinant of the product of the $n$ digit
matrices $D(a_0),\ldots,D(a_{n-1})$. The initial conventions
$p_{-1}=1,q_{-1}=0$ also cover $n=1$.
\end{proof}

\begin{theorem}[Exactly which convergents solve Pell]
\label{thm:pell-convergents}
For each $m\ge1$,
\begin{equation}\label{eq:pell-period}
 p_{mL-1}^2-Dq_{mL-1}^2=(-1)^{mL}.
\end{equation}
These are all positive integer solutions to the two equations
$p^2-Dq^2=\pm1$. In particular:
\begin{itemize}
 \item If $L$ is even, $+1$ occurs at every index $mL-1$, and the
 equation with right-hand side $-1$ has no integer solution.
 \item If $L$ is odd, $-1$ occurs for odd $m$ and $+1$ for even $m$.
\end{itemize}
The least positive solution of the $+1$ equation is
$(p_{L-1},q_{L-1})$ for even $L$ and $(p_{2L-1},q_{2L-1})$ for odd $L$.
\end{theorem}
\begin{proof}
All $Q_n$ with $n\ge1$ are positive. The lemma says that a convergent
solves one of the two equations exactly when $Q_n=1$.
By \eqref{eq:pell-state-return}, this happens exactly when $n=mL$;
\eqref{eq:pell-period} and its sign follow. Legendre's criterion, already
applied above, ensures that no positive solution lies outside the
convergents. Finally, the recurrence makes $p_0,p_1,\ldots$ strictly
increasing, so the first admissible index gives the least $p$, and hence
the least $q$ as well. If a solution of the $-1$ equation had either
coordinate zero, it would require $p^2=-1$ or $Dq^2=1$, both impossible;
changing signs reduces every other integer solution to a positive one.
\end{proof}

\begin{warningbox}[title={Stop before the final digit}]
The first candidate is
\[
 \frac{p_{L-1}}{q_{L-1}}=[a_0;a_1,\ldots,a_{L-1}],
\]
which stops \emph{before} the final $2a_0$. There are $L$ digits here,
counting $a_0$. For $\sqrt3=[1;\overline{1,2}]$, the candidate is
$[1;1]=2$, and $2^2-3\cdot1^2=1$. Including the next digit gives
$[1;1,2]=5/3$, whose remainder is $25-27=-2$.
\end{warningbox}

\begin{example}[An odd period: $D=13$]
The state calculation gives
$\sqrt{13}=[3;\overline{1,1,1,1,6}]$, with $L=5$.
The lemma explains every entry of the last column below.
\begin{center}
\begin{tabular}{ccrrr}
\toprule
$k$ & $a_k$ & $p_k/q_k$ & $Q_{k+1}$ & $p_k^2-13q_k^2$\\
\midrule
0 & 3 & $3/1$ & 4 & $-4$\\
1 & 1 & $4/1$ & 3 & $3$\\
2 & 1 & $7/2$ & 3 & $-3$\\
3 & 1 & $11/3$ & 4 & $4$\\
4 & 1 & $18/5$ & 1 & $-1$\\
5 & 6 & $119/33$ & 4 & $4$\\
6 & 1 & $137/38$ & 3 & $-3$\\
7 & 1 & $256/71$ & 3 & $3$\\
8 & 1 & $393/109$ & 4 & $-4$\\
9 & 1 & $649/180$ & 1 & $1$\\
\bottomrule
\end{tabular}
\end{center}
Thus the first negative solution is $(18,5)$, while the first positive
one is $(649,180)$. Squaring the former gives the latter:
\[
 (18+5\sqrt{13})^2=649+180\sqrt{13}.
\]
\end{example}

\subsection*{The period matrix becomes a Pell matrix}
Set $p=p_{L-1}$, $q=q_{L-1}$ and
$M=D(a_0)\cdots D(a_{L-1})$. At the returning state $P_L=a_0,Q_L=1$,
the two coefficient identities in the lemma give
\[
 M\begin{pmatrix}1&a_0\\0&1\end{pmatrix}
 =\begin{pmatrix}p&a_0p+p_{L-2}\\q&a_0q+q_{L-2}\end{pmatrix}
 =\begin{pmatrix}p&Dq\\q&p\end{pmatrix}.
\]
The shift by $a_0$ accounts for the fact that $M$ sends the tail
$\sqrt D+a_0$ to $\sqrt D$. The resulting matrix fixes $\sqrt D$ and
has determinant $p^2-Dq^2=(-1)^L$. For odd $L$, square it to obtain
determinant $+1$.

\subsection*{Why one fundamental solution generates all the others}
For $z=r+s\sqrt D$, define its \emph{norm} by
$N(z)=zz'=r^2-Ds^2$. Conjugation respects multiplication, so
$N(zw)=N(z)N(w)$. If $(u,v)$ is the least positive solution of the $+1$
equation, put $\varepsilon=u+v\sqrt D>1$. Its inverse is
$u-v\sqrt D$, again with integer coefficients.

\begin{proposition}[All positive Pell solutions are powers]
\label{prop:pell-powers}
Every positive solution of $p^2-Dq^2=1$ has the form
\[
 p+q\sqrt D=\varepsilon^j\qquad(j=1,2,\ldots).
\]
\end{proposition}
\begin{proof}
Powers have integer coefficients, norm $1$, and positive coefficients
for $j\ge1$. Conversely, let $\alpha=p+q\sqrt D>1$ have norm $1$.
Choose $j\ge0$ with $\varepsilon^j\le\alpha<\varepsilon^{j+1}$ and put
$\beta=\alpha\varepsilon^{-j}=r+s\sqrt D$. Then $r,s$ are integers,
$N(\beta)=1$, and $1\le\beta<\varepsilon$. Since $\beta'=1/\beta$,
\[
 r=\frac{\beta+\beta^{-1}}2>0,
 \qquad s=\frac{\beta-\beta^{-1}}{2\sqrt D}\ge0.
\]
If $\beta>1$, the function $t-t^{-1}$ is strictly increasing for $t>0$,
so $0<s<v$. This would be a smaller positive Pell solution, a
contradiction. Therefore $\beta=1$ and $\alpha=\varepsilon^j$.
Since $q>0$, we have $j\ge1$.
\end{proof}

Equivalently, multiplication by $\varepsilon$ is the matrix recurrence
\[
 \begin{pmatrix}p^{(j+1)}\\q^{(j+1)}\end{pmatrix}
 =\begin{pmatrix}u&Dv\\v&u\end{pmatrix}
  \begin{pmatrix}p^{(j)}\\q^{(j)}\end{pmatrix},
 \qquad (p^{(0)},q^{(0)})=(1,0).
\]
The superscript counts powers; the subscript on $p_k/q_k$ counts
continued-fraction convergents. For $D=2$, $\varepsilon=3+2\sqrt2$ gives
$(3,2),(17,12),(99,70)$.

\begin{bigidea}
Periodicity gives the exact return $Q_{mL}=1$. The determinant identity
then turns that return into $p_{mL-1}^2-Dq_{mL-1}^2=(-1)^{mL}$.
Legendre proves completeness, and powers of the least positive solution
give a second way to list every positive solution. Each step has a
different job.
\end{bigidea}

\begin{checkpoint}
Compute $99/70$ as a decimal. Its accuracy follows from
$99^2-2\cdot70^2=1$, which forces
\[
 \abs{\sqrt2-99/70}=\frac1{70^2(\sqrt2+99/70)}.
\]
For a second check, use the period of $\sqrt{23}$ above to find the
fundamental solution $(24,5)$ without searching through denominators.
\end{checkpoint}

\section{The Cattle of the Sun: Archimedes' challenge}
\label{sec:archimedes-cattle}

The cattle problem attributed to Archimedes asks for a herd belonging to
the Sun, with bulls and cows in four colors: white, black, yellow, and
dappled. Its proportions can be written as linear equations. Two further
conditions ask for a square and a triangle. Together they lead to a Pell
equation with a remarkably large first solution.

This optional application begins with the
\href{https://math.nyu.edu/Archimedes/Cattle/Statement.html}{ancient problem's conditions}.
For a detailed modern exposition, see Ilan Vardi,
\href{https://doi.org/10.1080/00029890.1998.12004887}{``Archimedes' Cattle Problem,''
\emph{American Mathematical Monthly} \textbf{105} (1998), 305--319};
an \href{https://www.maths.tcd.ie/~lebed/Cattle\%20problem\%20Ilan\%20Vardi.pdf}{accessible preprint}
is also available. The calculations below can be reproduced with exact
integer arithmetic in the first notebook.

\subsection*{Eight unknowns, seven linear conditions}
Let $W,B,Y,S$ count white, black, yellow, and spotted (dappled) bulls;
let $w,b,y,s$ count cows of the corresponding colors. The conditions are
\begin{align*}
 W&=\left(\frac12+\frac13\right)B+Y,
 &B&=\left(\frac14+\frac15\right)S+Y,
 &S&=\left(\frac16+\frac17\right)W+Y,\\
 w&=\left(\frac13+\frac14\right)(B+b),
 &b&=\left(\frac14+\frac15\right)(S+s),\\
 y&=\left(\frac16+\frac17\right)(W+w),
 &s&=\left(\frac15+\frac16\right)(Y+y).
\end{align*}
The cow equations use an entire herd of the specified color, including
its bulls. The two additional conditions concern bulls only:
\begin{equation}\label{eq:cattle-shapes}
 W+B=\text{a square},\qquad
 Y+S=\frac{h(h+1)}2\quad\text{for some integer }h\ge1.
\end{equation}

Eliminate $B,Y,S$ from the first three equations to obtain
\[
 B=\frac{267}{371}W,\qquad
 Y=\frac{297}{742}W,\qquad
 S=\frac{790}{1113}W.
\]
The four cow equations form a cycle. Substituting around that cycle
expresses $w$ as a known multiple of $W$ plus
$\frac7{12}\frac9{20}\frac{11}{30}\frac{13}{42}w$.
Since this coefficient is less than $1$, $w$, and then every cow count,
is uniquely determined by $W$. Clearing denominators gives
\begin{center}
\begin{tabular}{lrr}
\toprule
Color & Bulls & Cows\\
\midrule
White   & $10\,366\,482t$ & $7\,206\,360t$\\
Black   & $7\,460\,514t$  & $4\,893\,246t$\\
Yellow  & $4\,149\,387t$  & $5\,439\,213t$\\
Dappled & $7\,358\,060t$  & $3\,515\,820t$\\
\bottomrule
\end{tabular}
\end{center}
Substitution verifies all seven equations, and the eight coefficients
have gcd $1$. Therefore all eight counts are integers exactly when
$t$ is an integer: a B\'ezout combination of those counts equals $t$.
Every positive solution of the linear part has $t\ge1$, and the total
number of cattle is
\begin{equation}\label{eq:cattle-total-scale}
 T=50\,389\,082t.
\end{equation}

\subsection*{The square and the triangle produce Pell}
The square condition is
\[
 W+B=17\,826\,996t
 =4\cdot4\,456\,749t,
 \qquad4\,456\,749=3\cdot11\cdot29\cdot4657.
\]
The four factors in the last product are distinct primes (trial division
checks $4657$). Unique factorization therefore says that this is a square
exactly when
\[
 t=4\,456\,749r^2\qquad(r\ge1).
\]
Indeed, every prime in the displayed squarefree product must have odd
exponent in $t$, and every other prime must have even exponent.
Now the triangular condition becomes
\[
 11\,507\,447t=\frac{h(h+1)}2,
 \qquad
 51\,285\,802\,909\,803r^2=\frac{h(h+1)}2.
\]
Multiplying by $8$ and completing the square gives
\begin{equation}\label{eq:cattle-large-pell}
 \boxed{\quad X^2-410\,286\,423\,278\,424r^2=1,
 \qquad X=2h+1.\quad}
\end{equation}
Conversely, a positive solution of this equation has odd $X$, since the
coefficient of $r^2$ is divisible by $8$. Thus $h=(X-1)/2$ is an integer
and reconstructs the triangle. Together with the table, this proves an
equivalence, not just a necessary condition.

The coefficient is large, but contains a large square factor:
\[
 410\,286\,423\,278\,424=4\,729\,494\cdot9314^2.
\]
Put $Z=9314r$. The equivalent smaller-coefficient problem is
\begin{equation}\label{eq:cattle-small-pell}
 \boxed{\quad X^2-4\,729\,494Z^2=1,
 \qquad 9314\divides Z.\quad}
\end{equation}
Keeping the divisibility condition is essential: an arbitrary solution
of the smaller Pell equation need not give an integer $r$.

\subsection*{A short period, followed by a divisibility test}
The integer-state algorithm gives $a_0=2174$ and least period length
$L=92$ for $\sqrt{4\,729\,494}$. The final digit is $4348$; the preceding
$91$ digits form a palindrome. Since $92$ is even, the convergent
$p_{91}/q_{91}$ gives the fundamental solution
\begin{align*}
 u&=109931986732829734979866232821433543901088049,\\
 v&=50549485234315033074477819735540408986340.
\end{align*}
These values can be checked by the exact equality
$u^2-4\,729\,494v^2=1$. Write
$\varepsilon=u+v\sqrt{4\,729\,494}$ and
\[
 \varepsilon^j=X_j+Z_j\sqrt{4\,729\,494}.
\]
Proposition~\ref{prop:pell-powers} proves that this lists every positive
solution of the smaller Pell equation. We need the first $j$ for which
$9314\divides Z_j$.

We can find it without ever forming the enormous powers. Reduce the
multiplication recurrence modulo $9314$. Here
\[
 u\congr9063,\quad v\congr7708,\quad
 4\,729\,494\congr7296\pmod{9314}.
\]
The following finite computation checks every exponent up to the first
success, using only small integers:
\begin{quote}\begin{minipage}{\linewidth}\small
\begin{verbatim}
modulus = 9314
u, v, d = 9063, 7708, 7296
a, b = 1, 0
for j in range(1, 2330):
    a, b = ((u*a + d*v*b) % modulus,
            (v*a + u*b) % modulus)
    if b == 0:
        print(j, a, b)
        break
\end{verbatim}
\end{minipage}\end{quote}
Its output is \texttt{2329 9313 0}; no earlier exponent passes the test.
Thus $\varepsilon^{2329}$ produces the smallest admissible $r$ and
hence the smallest herd. This is a reproducible finite calculation,
distinct from the general Pell proofs above. Vardi explains a shorter
divisibility test using a quadratic finite field.

\subsection*{An enormous answer with a compact description}
Let $X+Z\sqrt{4\,729\,494}=\varepsilon^{2329}$. Then the least total is
\begin{equation}\label{eq:cattle-exact-total}
 T=50\,389\,082\cdot4\,456\,749
       \left(\frac{Z}{9314}\right)^2.
\end{equation}
This is an exact specification using integer arithmetic. The first
notebook computes $X,Z$ by repeated squaring, verifies the seven
proportions and both shape conditions, and checks
\[
 10^{206544}\le T<10^{206545}.
\]
The decimal expansion of the smallest herd's size therefore has
$206\,545$ digits:
\[
 T\approx7.760271406\times10^{206544}.
\]
Its size is a mathematical consequence of the conditions, rather than
a plausible census of animals.

Conjugation also gives a compact real-number expression, of the kind
developed in Vardi's exposition. Since
\[
 Z=\frac{\varepsilon^{2329}-\varepsilon^{-2329}}
          {2\sqrt{4\,729\,494}},
\]
substitution into \eqref{eq:cattle-exact-total} yields
\[
 T=\frac{25\,194\,541}{184\,119\,152}
       \left(\varepsilon^{4658}+\varepsilon^{-4658}-2\right).
\]
An expression can be short even when its decimal expansion is long.

\begin{bigidea}
The linear conditions leave one scale factor. Prime-exponent parity
makes that factor a fixed number times a square. The triangular
condition becomes Pell, and a divisibility test selects the allowed
powers of its fundamental solution. The same elementary tools that
solve $p^2-2q^2=1$ organize the cattle problem.
\end{bigidea}


\begin{exercisebox}
\begin{enumerate}
  \item Find the continued fraction of $355/113$. List all convergents. \exlevel{Fluency}
  \item Compute the first six convergents of $[1;\overline{1}]$. What familiar sequence appears? \exlevel{Fluency}
  \item Prove the determinant identity by induction. Use it to prove $\gcd(p_k,q_k)=1$. \exlevel{Reasoning}
  \item Show that $[a_0;\dots,a_n]=[a_0;\dots,a_n-1,1]$ when $a_n>1$. Explain the only ambiguity in finite expansions. \exlevel{Reasoning}
  \item Derive $\sqrt3=[1;\overline{1,2}]$ and compute four convergents. \exlevel{Reasoning}
  \item Write the tropical year as $365+\alpha$ with
  $\alpha\approx0.2421988$. Compute the first four nonzero convergents of
  $\alpha$ and identify the Julian rule. Then impose a century-based cycle by
  approximating $100\alpha$; explain how $97/4$ becomes the Gregorian rule
  $97/400$. Compare it with $31/128$. \exlevel{Exploration}
  \item Investigate which convergents to $\sqrt2$ solve $p^2-2q^2=1$ and which solve $-1$. State and prove a pattern using the recurrence. \exlevel{Challenge}
  \item Let $A/B=[a_0;\dots,a_n]$ and use the convergent recurrences to prove
  $r_k=(-1)^k(q_kA-p_kB)$. Apply the penultimate convergent of $527/341$ to
  recover a B\'ezout identity for $\gcd(527,341)$. \exlevel{Reasoning}
  \item For $x=\sqrt3$, compute enough convergents to display three lower
  bounds and three upper bounds. Verify that the lower bounds increase and the
  upper bounds decrease, then explain why the M\"obius derivative has the sign
  needed for each chain. \exlevel{Reasoning}
  \item Prove the matrix product formula \eqref{eq:cf-matrix-product} by
  induction. Explain why its inverse always has integer entries. \exlevel{Reasoning}
  \item Long-divide $1/7$ and $5/12$ in base $10$, recording the remainder at
  every step. Use repeated remainders to predict the preperiod and period.
  \exlevel{Fluency}
  \item Prove directly from a geometric series that every eventually periodic
  base-$b$ expansion is rational. \exlevel{Reasoning}
  \item After studying Euler's theorem, return to prove that $\gcd(q,b)=1$
  implies the base-$b$ expansion of $m/q$ is purely periodic. Give an explicit
  upper bound for a period length. \exlevel{Reasoning}
  \item Verify that $355/113$ satisfies Legendre's inequality for $\pi$ using
  a calculator value of $\pi$. What conclusion does the theorem certify?
  \exlevel{Exploration}
  \item Let $t=[\overline{1,2}]$. Use the matrix $D(1)D(2)$ to find a quadratic
  equation for $t$, and identify the positive root. \exlevel{Reasoning}
  \item Check Galois's inequalities for $1+\sqrt2$ and
  $(1+\sqrt5)/2$. Explain why their continued fractions are purely periodic.
  \exlevel{Reasoning}
  \item Let $p^2-Dq^2=\pm1$ with $p,q>0$. Derive
  \eqref{eq:pell-error}, verify its strict Legendre bound, and conclude that
  $p/q$ is a convergent of $\sqrt D$. \exlevel{Reasoning}
  \item Compute the period of $\sqrt3$. Use its length to predict whether one
  period solves the $+1$ or $-1$ Pell equation. Then use powers of
  $\begin{psmallmatrix}2&3\\1&2\end{psmallmatrix}$ to generate three solutions
  of $p^2-3q^2=1$. \exlevel{Exploration}
  \item Define $G_n(x)=(a_1\cdots a_n)^{1/n}$ and $B_n(x)=q_n^{2/n}$.
  Use the Sage notebook to reproduce several rows of the $\pi$ table above.
  Explain why agreement with the Khinchin--L\'evy constants is evidence but
  not proof. Then derive the displayed formula for $G_{3m}(e)$ directly from
  the pattern in the continued fraction of $e$. \exlevel{Exploration}
  \item Compute $D(1)D(2)$ and $D(2)D(1)$, and write the linear
  fractional transformation associated to each. Explain the order of
  composition and find each transformation's positive fixed point.
  \exlevel{Fluency}
  \item Let $x=[\overline{1,2}]$ and $N=D(1)D(2)$. For $j=1,2$, set
  $M=D(j)$ and compute $MNM^{-1}$. Use its fixed-point equation to
  identify $y=[j;\overline{1,2}]$. Which of these numbers is purely
  periodic? Write its period without an initial prefix.
  \exlevel{Reasoning}
  \item Prove that an invertible rational linear fractional transformation
  sends every irrational real number to an irrational real number.
  Explain why a pole cannot occur at an irrational input. Then identify
  the missing hypothesis in the assertion that $t=T_M(t)$ always makes
  $t$ quadratic. Give a counterexample. \exlevel{Reasoning}
  \item Use Galois's criterion to decide which of
  \[
   \sqrt2,\quad 1+\sqrt2,\quad \frac{1+\sqrt5}{2},\quad
   \frac{1+\sqrt3}{2},\quad \frac{\sqrt2}{2},\quad -\sqrt2
  \]
  have purely periodic continued fractions. Explain why all six are
  eventually periodic. \exlevel{Fluency}
  \item Starting from $D=13$ and $(P_0,Q_0)=(0,1)$, use
  \eqref{eq:quadratic-state-recurrence} to compute the continued fraction
  of $\sqrt{13}$. Record $(P_n,Q_n,a_n)$ until a state first repeats,
  and verify the reduced-state bounds after the first step.
  \exlevel{Exploration}
  \item Recover the unique reduced predecessor of
  $z=(1+\sqrt3)/2$ using \eqref{eq:reduced-predecessor}.
  For $D=3$, enumerate the integer pairs satisfying
  \eqref{eq:reduced-state-bounds} that actually represent reduced surds
  and satisfy $Q\divides3-P^2$. Show how the continued-fraction step
  permutes them. Why would finiteness without reversibility not suffice
  to prove pure periodicity? \exlevel{Reasoning}
  \item Let $x=[\overline{1,2,3}]$. Compute its period matrix and its
  positive fixed point. Transpose the matrix to obtain the equation for
  $[\overline{3,2,1}]$, and verify directly that this number is $-1/x'$.
  \exlevel{Reasoning}
  \item Compute the integer states and least period for $\sqrt{23}$.
  Identify the palindrome, find the fundamental Pell solution, and
  calculate what happens if you include the final digit $2a_0$ in the
  first candidate. \exlevel{Fluency}
  \item For $D=6$, calculate four convergents and the corresponding
  $Q_{k+1}$. Verify \eqref{eq:pell-remainder} in each row. Use
  \eqref{eq:pell-state-return} to classify every convergent that solves
  either sign of Pell. \exlevel{Reasoning}
  \item Decide whether $p^2-Dq^2=-1$ is solvable for
  $D=2,3,5,7,13,14$. Give the least positive solution when it exists.
  For $D=3$, also give a one-line impossibility proof modulo $3$.
  \exlevel{Reasoning}
  \item Let $p^2-3q^2=1$ with $p,q>0$. Show that multiplying
  $p+q\sqrt3$ by $2-\sqrt3$ gives $(1,0)$ from $(2,1)$,
  and otherwise a positive solution with smaller second coordinate.
  Use this descent to prove directly that all positive solutions are
  powers of $2+\sqrt3$. \exlevel{Reasoning}
  \item Use the integer-state algorithm to find the least period of
  $\sqrt{61}$. Find the first negative and positive Pell solutions,
  and explain why the theorem proves that each is the first of its
  sign. Use exact integers, not a decimal approximation to $\sqrt{61}$.
  \exlevel{Exploration}
  \item Verify the cattle table at $t=1$, including its gcd and total.
  Explain why the seven linear equations have only one free scale
  factor, and why the gcd certifies that this scale is an integer.
  Does the table at $t=1$ solve the full problem? \exlevel{Reasoning}
  \item For positive integers $c,r$, prove that $cr^2$ is triangular
  exactly when $X^2-8cr^2=1$ has a positive odd solution $X$.
  Apply this to the cattle problem, and prove both directions of the
  passage from \eqref{eq:cattle-large-pell} to
  \eqref{eq:cattle-small-pell}. \exlevel{Reasoning}
  \item Reproduce the modular search for the cattle problem in the
  first notebook. Explain why its first hit proves minimality. Using
  $\varepsilon^{2329}\congr-1\pmod{9314}$ and the absence of an
  earlier hit, prove that the admissible exponents are exactly the
  positive multiples of $2329$. Why does $4658$ appear in the total's
  compact formula? \exlevel{Exploration}
\end{enumerate}
\end{exercisebox}

\chapter{Arithmetic on a clock}
\label{chap:modular}
\chaptertag{Extended Euclid, Chinese remainders, and interpolation}

\begin{bigidea}
The extended gcd algorithm is the computational engine of modular arithmetic.
It tells us when division is possible and supplies the inverse when it is.
We will use those inverses to reconstruct an integer from its remainders,
then recognize the same construction in polynomial interpolation.
Finally, a reversible recurrence will explain why Fibonacci residues repeat.
\end{bigidea}

\begin{definition}
For $n\ge2$, we write $a\congr b\pmod n$ if $n\divides(a-b)$. The equivalence class of $a$ is written $[a]$ or $\bar a$.
\end{definition}

Congruence remembers the remainder and forgets the number of complete turns around the clock.

\begin{center}
\begin{tikzpicture}[scale=1.0]
  \draw[navy!35,thick] (0,0) circle (1.65);
  \foreach \k in {0,...,11}{
    \coordinate (P\k) at ({90-30*\k}:1.65);
    \fill[teal] (P\k) circle (2pt);
    \node at ({90-30*\k}:1.94) {$\k$};
  }
  \draw[-{Stealth[length=3mm]},coral,line width=1.5pt] (P9) arc[start angle=180,end angle=30,radius=1.65];
  \node[coral,align=left,anchor=west] at (2.6,.2) {$9+5\congr2\pmod{12}$\\one complete turn disappears};
\end{tikzpicture}
\end{center}

\section{Legal arithmetic}

If $a\congr b\pmod n$ and $c\congr d\pmod n$, then
\[
a+c\congr b+d\pmod n,\qquad ac\congr bd\pmod n.
\]
Powers are therefore legal too. Division is different.

These rules do not depend on the chosen representatives. For example, if
$a=b+kn$ and $c=d+\ell n$, then
$ac-bd=n(b\ell+dk+k\ell n)$. Thus addition and multiplication are
well-defined operations on the residue classes $\Z/n\Z$. Reduce after
each step or at the end: the resulting residue class is the same.

\section{Extended gcd: division comes with a certificate}
\label{sec:modular-egcd}

\begin{definition}
A residue $a$ is a \emph{unit modulo $n$} if there exists $b$ with $ab\congr1\pmod n$. Such a $b$ is an inverse of $a$.
\end{definition}

\begin{theorem}[Invertibility criterion]
$a$ has an inverse modulo $n$ exactly when $\gcd(a,n)=1$.
\end{theorem}

\begin{proof}
The congruence $ab\congr1\pmod n$ means $ab+nk=1$ for some $k$, so any
common divisor of $a,n$ divides $1$. Conversely, when $\gcd(a,n)=1$,
extended Euclid produces $u,v$ such that $au+nv=1$. Reducing this
\emph{integer identity} modulo $n$ gives $au\congr1$: the coefficient
$u$, reduced modulo $n$, is the required inverse.
\end{proof}

The theorem is an algorithm, not a direction to try every residue until
one works. Here is a forward version of the backward substitutions from
Chapter~1. First reduce $a$ so that $0\le a<n$. Keep three entries in each row:
\[
 r_i=ns_i+at_i.
\]
Start with $(r_0,s_0,t_0)=(n,1,0)$ and $(r_1,s_1,t_1)=(a,0,1)$.
Whenever $r_i\ne0$, divide $r_{i-1}=q_i r_i+r_{i+1}$ and update
\[
 (r_{i+1},s_{i+1},t_{i+1})
 =(r_{i-1},s_{i-1},t_{i-1})-q_i(r_i,s_i,t_i).
\]
The same subtraction in all three columns preserves the displayed
identity. The last nonzero row gives the gcd and its B\'ezout coefficients.

\begin{example}[An inverse without guessing]
Find $17^{-1}\pmod{43}$. The divisions are
\[
43=2\cdot17+9,\quad17=9+8,\quad9=8+1,\quad8=8\cdot1.
\]
The coefficient table records $r=43s+17t$ throughout:
\begin{center}
\begin{tabular}{rrrl}
\toprule
$r$ & $s$ & $t$ & Certificate \\
\midrule
43 & 1 & 0 & $43=1\cdot43+0\cdot17$ \\
17 & 0 & 1 & $17=0\cdot43+1\cdot17$ \\
9 & 1 & $-2$ & $9=43-2\cdot17$ \\
8 & $-1$ & 3 & $8=-43+3\cdot17$ \\
1 & 2 & $-5$ & $1=2\cdot43-5\cdot17$ \\
0 & $-17$ & 43 & $0=-17\cdot43+43\cdot17$ \\
\bottomrule
\end{tabular}
\end{center}
Hence $17^{-1}\congr-5\congr38\pmod{43}$. Check the result directly:
$17\cdot38=646=15\cdot43+1$.
\end{example}

\begin{checkpoint}
In a certificate $au+nv=1$, which coefficient is the inverse of $a$
modulo $n$? Read the coefficient of $a$, even if the table began with
$n$ in its first row. The final multiplication check catches a swapped
coefficient or a sign error.
\end{checkpoint}

\begin{warningbox}
Cancellation can fail. Modulo $6$, $2\cdot1\congr2\cdot4$, but $1\not\congr4$. You may cancel $a$ modulo $n$ only when $a$ is a unit, or after carefully reducing the modulus.
\end{warningbox}

For a prime $p$, every nonzero residue is a unit. We write
$\F_p=\Z/p\Z$ and call it a \emph{field}: addition and multiplication obey
the usual commutative arithmetic laws, and every nonzero element permits
division. The rational and real numbers are other examples of fields.
For composite $n$, $\Z/n\Z$ need not permit division by a nonzero residue.

\section{Solving a linear congruence with the same algorithm}

\begin{proposition}[Linear congruence]
The congruence $ax\congr b\pmod n$ has a solution exactly when $d=\gcd(a,n)$ divides $b$. When it does, there are exactly $d$ incongruent solutions modulo $n$.
\end{proposition}

\begin{proof}[Construction and completeness]
Run extended Euclid once to obtain $au+nv=d$. If $ax\congr b\pmod n$,
then $b=ax-nk$ for an integer $k$, so $d\divides b$ is necessary.
If $d\divides b$, multiply the certificate by $b/d$ to obtain a solution
\[
 x_0=u\frac bd.
\]
If $d=n$, then both $a$ and $b$ are multiples of $n$, so all $n$ residues
are solutions. Otherwise $n/d\ge2$, and any two solutions satisfy
$(a/d)(x-x_0)\congr0\pmod{n/d}$. Since $a/d$ is a unit modulo $n/d$,
this is equivalent to $x\congr x_0\pmod{n/d}$. Therefore the complete list
modulo $n$ is
\[
 x_0,\ x_0+\frac nd,\ \ldots,\ x_0+(d-1)\frac nd.
\]
These $d$ residues are distinct.
\end{proof}

\begin{example}[A nonunit can give several solutions]
For $14x\congr21\pmod{35}$, extended Euclid gives
$7=35-2\cdot14$. Multiplying by $3$ gives $21=3\cdot35-6\cdot14$,
so $x\congr-6\congr4\pmod5$. The seven solutions modulo $35$ are
$4,9,14,19,24,29,34$. The same calculation rules out
$14x\congr20\pmod{35}$, because $7\ndivides20$.
\end{example}

\section{Reassembling remainders}
\label{sec:crt}

Suppose a number leaves remainder $2$ modulo $3$ and remainder $3$ modulo $5$. Listing possibilities gives $8$ modulo $15$. The Chinese remainder theorem says this always works for coprime clocks.

\begin{theorem}[Chinese remainder theorem]
If $n_1,\dots,n_r$ are pairwise coprime, then for any $a_1,\dots,a_r$, the system
\[
x\congr a_i\pmod{n_i}\qquad(i=1,\dots,r)
\]
has a solution, unique modulo $N=n_1\cdots n_r$.
\end{theorem}

\begin{proof}[Construction]
Set $N_i=N/n_i$. Since $\gcd(N_i,n_i)=1$, \emph{use extended Euclid} to
find $M_i,v_i$ with $N_iM_i+n_iv_i=1$. Then
\[
x=\sum_{i=1}^r a_iN_iM_i
\]
has every prescribed remainder: modulo $n_i$, its $i$th term is $a_i$ and every other term is $0$. If $x,y$ both work, each $n_i$ divides $x-y$; coprimality makes $N\divides x-y$.
\end{proof}

The quantities $e_i=N_iM_i$ are \emph{selectors}: $e_i$ is $1$ at its
own modulus and $0$ at every other modulus. The term $a_i e_i$ installs
one prescribed remainder without disturbing the others.

\begin{example}[Construct, reduce, and check]
Solve $x\congr2\pmod3$, $x\congr3\pmod5$, $x\congr2\pmod7$.
Here $N=105$, and the construction is
\begin{center}
\begin{tabular}{rrrrr}
\toprule
$n_i$ & $a_i$ & $N_i$ & $M_i$ & $e_i=N_iM_i$ \\
\midrule
3 & 2 & 35 & 2 & 70 \\
5 & 3 & 21 & 1 & 21 \\
7 & 2 & 15 & 1 & 15 \\
\bottomrule
\end{tabular}
\end{center}
For instance, $2\cdot35-23\cdot3=1$ certifies the first inverse.
Thus $x=2\cdot70+3\cdot21+2\cdot15=233\congr23\pmod{105}$.
Finally, $23$ has remainders $2,3,2$, as required. Every solution is
$23+105k$ for an integer $k$.
\end{example}

\begin{center}
\begin{tikzpicture}[node distance=8mm]
\node[draw=teal,fill=mist,rounded corners,minimum width=21mm] (z15) {$\Z/15\Z$};
\node[draw=navy!55,fill=white,rounded corners,below left=of z15,minimum width=21mm] (z3) {$\Z/3\Z$};
\node[draw=navy!55,fill=white,rounded corners,below right=of z15,minimum width=21mm] (z5) {$\Z/5\Z$};
\draw[-{Stealth},thick,teal] (z15)--(z3) node[midway,left,font=\scriptsize] {$x\mapsto x\bmod3$};
\draw[-{Stealth},thick,teal] (z15)--(z5) node[midway,right,font=\scriptsize] {$x\mapsto x\bmod5$};
\node[below=12mm of z15,muted] {One clock of size $15$ is equivalent to two independent clocks.};
\end{tikzpicture}
\end{center}

The map from one clock to its component clocks preserves addition and
multiplication, since reduction does. CRT says it is also a bijection.
For example, the selectors modulo $15$ are $10$ and $6$, giving
$x\congr10a+6b\pmod{15}$ for remainders $a$ modulo $3$ and $b$ modulo $5$.

\begin{proposition}[When two clocks share a factor]
The conditions $x\congr a\pmod m$, $x\congr b\pmod n$ are compatible
exactly when $a\congr b\pmod{\gcd(m,n)}$. If compatible, the solution is
unique modulo $\operatorname{lcm}(m,n)$.
\end{proposition}
\begin{proof}
Write $x=a+mt$. The remaining condition is $mt\congr b-a\pmod n$,
so the linear-congruence construction supplies a solution exactly when
$\gcd(m,n)\divides b-a$. Differences between two solutions are precisely
the common multiples of $m$ and $n$.
\end{proof}

\section{Chinese remainders become Lagrange interpolation}
\label{sec:crt-interpolation}

CRT builds an integer with prescribed remainders. Interpolation builds
a polynomial with prescribed values. Polynomial division makes these
the same kind of problem.

Let $K$ be a field, such as $\Q$, $\R$, or $\F_p$. Write $K[X]$ for
polynomials in an indeterminate $X$ with coefficients in $K$. Division
by $X-c$ gives
\[
 f(X)=(X-c)q(X)+r,
\]
where $r$ is constant. Substituting $X=c$ shows $r=f(c)$. In other words,
\[
 f(X)\congr b\pmod{X-c}\quad\Longleftrightarrow\quad f(c)=b.
\]
Here polynomial congruence means that the difference is divisible by the
polynomial modulus.

\subsection*{The same construction, with polynomial moduli}

The division algorithm works in $K[X]$: cancel the leading term using
division in $K$, subtract, and continue while the degree decreases.
Extended Euclid therefore works too. For coprime polynomials $A,B$,
it produces $U,V$ with $UA+VB=1$.

Consequently, for pairwise coprime monic moduli $m_1,\ldots,m_r$, set
$P=\prod_i m_i$ and $P_i=P/m_i$, and use extended Euclid to find
$U_iP_i+V_im_i=1$. The polynomial
\[
 f=\sum_i b_iP_iU_i
\]
has prescribed polynomial remainders $b_i$ modulo $m_i$.
The proof is the integer CRT proof: reduce at each modulus separately.
Uniqueness is modulo $P$, since the pairwise coprime factors $m_i$ all
divide the difference of two solutions. Dividing by the monic polynomial
$P$ selects one representative of degree less than $\deg P$.

For interpolation, take $m_i=X-c_i$ with the $c_i$ distinct. Two such
linear polynomials are coprime: subtracting them gives the nonzero
constant $c_j-c_i$, which is invertible in $K$.
In this special case, the inverse required by CRT is particularly simple.
For
\[
 P_i(X)=\prod_{j\ne i}(X-c_j),
\]
reduction modulo $X-c_i$ gives the nonzero number
$P_i(c_i)=\prod_{j\ne i}(c_i-c_j)$.
Thus define
\[
 \ell_i(X)=\frac{P_i(X)}{P_i(c_i)}
 =\prod_{j\ne i}\frac{X-c_j}{c_i-c_j}.
\]
These are polynomial selectors: $\ell_i(c_i)=1$ and
$\ell_i(c_j)=0$ for $j\ne i$.

\begin{theorem}[Lagrange interpolation]
Let $c_1,\ldots,c_r$ be distinct points of a field $K$.
Given values $b_i\in K$, the unique polynomial of degree less than $r$
with $f(c_i)=b_i$ for every $i$ is
\[
 f(X)=\sum_{i=1}^r b_i\ell_i(X).
\]
Every polynomial with these values has the form
$f(X)+h(X)\prod_i(X-c_i)$, with $h\in K[X]$.
\end{theorem}

\begin{proof}
The selector values prove existence. If $f,g$ have the required values,
each $X-c_i$ divides $f-g$, and hence so does their product. A nonzero
multiple of that degree-$r$ product cannot have degree less than $r$.
This proves uniqueness and the description of all solutions.
The same argument proves the useful \emph{root bound}: a nonzero
polynomial of degree $d$ over a field has at most $d$ distinct roots.
\end{proof}

\begin{center}
\small
\begin{tabularx}{\textwidth}{p{.18\textwidth}XX}
\toprule
 & Integer CRT & Lagrange interpolation \\
\midrule
Local data & $x\congr a_i\pmod{n_i}$ & $f(c_i)=b_i$ \\
Moduli & pairwise coprime $n_i$ & distinct linear factors $X-c_i$ \\
Selector & $N_i(N_i^{-1}\bmod n_i)$ & $P_i(X)/P_i(c_i)$ \\
Formula & $\sum_i a_i e_i$ & $\sum_i b_i\ell_i(X)$ \\
Ambiguity & multiples of $N$ & multiples of $\prod_i(X-c_i)$ \\
Choice & $0\le x<N$ & $\deg f<r$ \\
\bottomrule
\end{tabularx}
\end{center}

\begin{example}[Three values determine a quadratic]
Over $\Q$, prescribe $f(0)=1$, $f(1)=2$, $f(2)=5$. The selectors are
\[
 \ell_0=\frac{(X-1)(X-2)}2,\qquad
 \ell_1=-X(X-2),\qquad
 \ell_2=\frac{X(X-1)}2.
\]
Combining them gives $f=\ell_0+2\ell_1+5\ell_2=X^2+1$.
All other solutions are $X^2+1+X(X-1)(X-2)h(X)$.

The same computation works in $\F_5$, where the prescribed values are
$1,2,0$ and division by $2$ means multiplication by $3$.
The certificate $2\cdot3-5=1$ makes that division legitimate.
The result is again $X^2+1$, now with coefficients modulo $5$.
\end{example}

\begin{warningbox}
Distinct interpolation points must have invertible differences.
This is automatic in a field. It can fail for composite moduli:
over $\Z/8\Z$, the distinct points $0$ and $2$ have difference $2$,
which has no inverse. Indeed, $f(2)-f(0)$ is even for every polynomial,
so the data $f(0)=0$, $f(2)=1$ cannot be interpolated.
\end{warningbox}

\begin{checkpoint}
Both reconstruction formulas first make an object that vanishes at every
unwanted coordinate, then multiply by an inverse to make its remaining
coordinate equal to $1$. Identify those two steps before memorizing either
formula.
\end{checkpoint}


\section{Fast powers and patterns}

Modular reduction prevents numbers from exploding. To compute $7^{222}\pmod{13}$, square and reduce:
\[
7^2\congr10,\quad 7^4\congr9,\quad7^8\congr3,\quad7^{16}\congr9\pmod{13},
\]
then use the binary expansion $222=128+64+16+8+4+2$. Later, Fermat's theorem will make this even faster.

\begin{labbox}
Make a multiplication table modulo $n$ and highlight the rows that are permutations of all residues. Conjecture which row labels have inverses. Then compare with the gcd criterion.
\end{labbox}

\section{Fibonacci residues: why the pattern must repeat}
\label{sec:fibonacci-periodicity}

Set $F_0=0$, $F_1=1$, and $F_{n+2}=F_{n+1}+F_n$. Modulo an integer
$m\ge2$, the next term depends on the previous \emph{two} terms.
The appropriate state is therefore
\[
 S_n=(F_n,F_{n+1})\pmod m,
 \qquad T(a,b)=(b,a+b)\pmod m.
\]
There are only $m^2$ pairs, so a state must eventually repeat.
But finiteness alone only gives eventual periodicity. Here we can do better:
\[
 T^{-1}(u,v)=(v-u,u)\pmod m.
\]
Every step has a unique predecessor, using only subtraction.

\begin{theorem}[Fibonacci periodicity from the starting pair]
For every $m\ge2$, the Fibonacci sequence modulo $m$ is periodic from
the start. Its least period $\pi(m)$ is the smallest positive $t$ for which
\[
 (F_t,F_{t+1})\congr(0,1)\pmod m.
\]
In particular, $\pi(m)\le m^2-1$.
\end{theorem}

\begin{proof}
If $S_i=S_j$ with $i<j$, apply the inverse step $i$ times to get
$S_0=S_{j-i}$. Repeating the forward step now gives a period beginning
at $S_0$, not just after an initial segment. Conversely, a period of the
sequence must return both $F_0$ and $F_1$, so it returns the starting pair.
Before the first return all states are distinct. The pair $(0,0)$ is fixed
by $T$ and has no other predecessor, so the orbit of $(0,1)$ never visits
it. There are at most $m^2-1$ available states.
\end{proof}

\begin{example}[A zero is not yet a restart]
Modulo $3$, the successive states are
\[
\begin{array}{c|rrrrrrrrr}
n&0&1&2&3&4&5&6&7&8\\\hline
F_n&0&1&1&2&0&2&2&1&0\\
F_{n+1}&1&1&2&0&2&2&1&0&1
\end{array}
\]
The pair first returns at $n=8$, so $\pi(3)=8$.
Although $F_4\congr0$, its successor is $2$, not $1$.
Looking only for a repeated Fibonacci value would give the wrong answer.
\end{example}

\begin{proposition}[CRT combines periods]
If $\gcd(m,n)=1$, then
\[
 \pi(mn)=\operatorname{lcm}(\pi(m),\pi(n)).
\]
\end{proposition}

\begin{proof}
CRT says that a pair returns modulo $mn$ exactly when it returns modulo
both $m$ and $n$. The return times modulo $m$ are precisely the multiples
of $\pi(m)$: divide any return time by $\pi(m)$ and use the minimality of
the remainder. Likewise for $n$. The first common return time is their
least common multiple.
\end{proof}

For example, the orbit modulo $2$ has length $3$, so
$\pi(6)=\operatorname{lcm}(3,8)=24$.
The connection is operational: CRT reassembles the two coordinates of
every state as well as the arithmetic on the clocks.

\begin{labbox}
For a small modulus, start at $(0,1)$, repeatedly replace $(a,b)$ by
$(b,a+b\bmod m)$, and stop on the first return to $(0,1)$.
The theorem proves this procedure terminates and that the answer is exact.
Record pairs, not just zeros. These small experiments need no extension fields.
\end{labbox}

\subsection*{What we can prove now, and what comes next}

The reversible-pair argument proves periodicity for \emph{every} modulus,
prime or composite. The estimate $\pi(m)\le m^2-1$ is an upper bound;
it does not assert that $\pi(m)$ divides $m^2-1$.
For example, direct iteration gives $\pi(5)=20$, and $20\ndivides24$.

A structural computation uses the roots of $X^2-X-1$.
Their $\sqrt5$ exists in $\F_p$ for $p=11$, but not for $p=7$.
Appendix~\ref{app:finite-fields} first builds arithmetic in an extension
field, then uses Frobenius and multiplicative orders to compute exact
periods. Those tools belong to a later unit; pair iteration and CRT
require only the arithmetic developed here.


\begin{exercisebox}
\begin{enumerate}
  \item Find the inverse of $37$ modulo $101$ using extended Euclid. \exlevel{Fluency}
  \item Solve $18x\congr30\pmod{42}$ and list all solutions modulo $42$. \exlevel{Fluency}
  \item Prove that congruence modulo $n$ is an equivalence relation. \exlevel{Reasoning}
  \item Show that an integer is divisible by $9$ exactly when its decimal digit sum is divisible by $9$. \exlevel{Reasoning}
  \item Solve $x\congr2\pmod7$, $x\congr3\pmod8$, $x\congr4\pmod9$ by the constructive CRT proof. \exlevel{Reasoning}
  \item Determine every possible last two digits of a square. Explain how CRT can split the work into moduli $4$ and $25$. \exlevel{Exploration}
  \item Prove that $x^2\congr1\pmod p$ has only the solutions $x\congr\pm1$ when $p$ is prime. Find a composite modulus where there are more than two solutions. \exlevel{Challenge}
  \item Run the forward extended-gcd table for $83$ and $29$, recording
  $r=83s+29t$ in every row. Find $29^{-1}\pmod{83}$ and check it by
  multiplication. \exlevel{Fluency}
  \item Solve $18x\congr24\pmod{42}$ using a B\'ezout certificate and
  list all residues modulo $42$. Explain why $18x\congr25\pmod{42}$
  has no solution. \exlevel{Fluency}
  \item For the CRT selectors, prove $e_i^2\congr e_i$,
  $e_ie_j\congr0$ when $i\ne j$, and $\sum_i e_i\congr1\pmod N$.
  Explain how these identities recover componentwise arithmetic.
  \exlevel{Reasoning}
  \item Solve $x\congr2\pmod6$, $x\congr5\pmod9$. Give the modulus
  of uniqueness. Why is $x\congr1\pmod6$, $x\congr2\pmod9$
  incompatible? \exlevel{Reasoning}
  \item Construct the polynomial of degree at most $2$ over $\Q$ with
  $f(-1)=2$, $f(0)=1$, $f(2)=3$ using Lagrange selectors.
  State all polynomials with these values. \exlevel{Fluency}
  \item Work in $\F_7$. Use Lagrange interpolation to find the
  polynomial of degree at most $2$ with $f(0)=1$, $f(1)=4$, $f(3)=2$.
  Certify each denominator inverse by an integer B\'ezout identity.
  \exlevel{Reasoning}
  \item Prove that no polynomial over $\Z/8\Z$ can satisfy
  $f(0)=0$, $f(2)=1$. Which hypothesis of the interpolation theorem
  fails? \exlevel{Reasoning}
  \item List the Fibonacci pair orbits modulo $2$ and $3$. Why is the
  first positive zero of $F_n\pmod3$ not the period? \exlevel{Fluency}
  \item Use CRT and $\pi(2)=3$, $\pi(3)=8$, $\pi(5)=20$ to find
  $\pi(6)$ and $\pi(15)$. Explain why the least common multiple,
  rather than the product, is the appropriate operation.
  \exlevel{Reasoning}
  \item Iterate $(a,b)\mapsto(b,2a+b)$ modulo $2$, starting at $(1,0)$.
  Show that the orbit has an initial segment before it repeats.
  Locate the failure of the Fibonacci reversibility proof.
  \exlevel{Reasoning}
  \item List the squares modulo $7$ and show that $X^2-X-1$ has no
  root in $\F_7$. Explain why the real Binet formula cannot simply
  be reduced modulo $7$ without first enlarging the arithmetic system.
  \exlevel{Bridge}
\end{enumerate}
\end{exercisebox}

\chapter{A theorem you can rotate}
\chaptertag{Necklace colorings and Fermat's little theorem}

Why should a prime $p$ divide $a^p-a$ for every integer $a$? Instead of manipulating the expression, we count objects that naturally arrive in packets of $p$.

\section{Words on a circular loom}

Fix a prime $p$ and $a$ colors. First imagine $p$ labeled positions around a circle. A coloring is a word
\[
(c_0,c_1,\dots,c_{p-1}),
\]
and there are $a^p$ such words. Rotation by one step sends it to
\[
(c_{p-1},c_0,\dots,c_{p-2}).
\]

\begin{center}
\begin{tikzpicture}[scale=.82]
\foreach \shift/\rot/\lab in {0/0/{first},4.4/72/{rotate},8.8/144/{rotate again}}{
  \begin{scope}[xshift=\shift cm,rotate=\rot]
    \draw[navy!25,thick] (0,0) circle (1.25);
    \foreach \ang/\col in {90/teal,18/gold,-54/teal,-126/coral,162/gold}{
      \fill[\col] (\ang:1.25) circle (5.5pt);
    }
  \end{scope}
  \node[muted] at (\shift,-1.75) {\lab};
}
\draw[-{Stealth},teal,thick] (1.5,0)--(2.9,0);
\draw[-{Stealth},teal,thick] (5.9,0)--(7.3,0);
\end{tikzpicture}
\end{center}

Group words together when one is a rotation of another. Such a group is called a rotational \emph{orbit}. How large can an orbit be?

\begin{lemma}[Prime orbit dichotomy]
For $p$ prime, a length-$p$ word has either $1$ or $p$ distinct rotations. It has one exactly when it is constant.
\end{lemma}

\begin{proof}
If rotation by $k$ positions, with $1\le k<p$, fixes the word, then entries $i$ and $i+k$ have the same color (indices modulo $p$). Since $\gcd(k,p)=1$, repeatedly adding $k$ visits every position. Thus all entries have one color. A nonconstant word has no repeated rotation before the $p$th, so its orbit has size $p$.
\end{proof}

There are $a$ constant words, one for each color. The remaining $a^p-a$ words split into orbits of exactly $p$ words.

\begin{theorem}[Fermat's little theorem, necklace form]
If $p$ is prime, then for every integer $a$,
\[
a^p\congr a\pmod p.
\]
If $p\ndivides a$, then $a^{p-1}\congr1\pmod p$.
\end{theorem}

\begin{proof}
For $a\ge0$, the orbit count above shows $p\divides a^p-a$. This congruence depends only on the residue of $a$ modulo $p$, so it holds for every integer $a$. If $p\ndivides a$, then $a$ has an inverse modulo $p$, and we may cancel it from $a^p\congr a$.
\end{proof}

\begin{bigidea}
Divisibility by $p$ appears because symmetry packages nonexceptional objects into groups of $p$. The exceptions are exactly the colorings fixed by rotation. This pattern -- count, isolate fixed points, divide the rest into orbits -- reaches far beyond number theory.
\end{bigidea}

\section{A subtle choice: labeled words or necklaces?}

A physical necklace usually ignores where position $0$ is. Here we first count \emph{labeled} circular words, then use rotation to partition them. Calling each orbit a necklace is convenient, but the proof depends on keeping those two levels distinct.

\begin{warningbox}
When the length is composite, an orbit can have an intermediate size. The word $RBRB$ has two distinct rotations, not $1$ or $4$. Primality is doing real work.
\end{warningbox}

\section{Using the theorem}

To compute $3^{1000}\pmod7$, note that $3^6\congr1$ and $1000=6\cdot166+4$:
\[
3^{1000}\congr(3^6)^{166}3^4\congr81\congr4\pmod7.
\]

Fermat's theorem also explains the first layer of binomial divisibility. For $1\le k\le p-1$,
\[
\binom pk=\frac{p!}{k!(p-k)!}
\]
is divisible by $p$, because the denominator contains no factor $p$. Therefore
\[
(x+y)^p\congr x^p+y^p\pmod p.
\]
We will return to this ``freshman's dream'' after developing binomial coefficients systematically.

\begin{exercisebox}
\begin{enumerate}
  \item List the rotational orbits of binary words of length $3$. Identify the constant words. \exlevel{Fluency}
  \item Compute $11^{2026}\pmod{17}$ using Fermat's theorem. \exlevel{Fluency}
  \item In the orbit lemma, prove carefully that $0,k,2k,\dots,(p-1)k$ are all distinct modulo $p$. \exlevel{Reasoning}
  \item Give an orbit-counting proof that $p\divides a^p-a$ using linear strings and cyclic shift, without drawing a circle. \exlevel{Reasoning}
  \item Find every possible orbit size for binary words of length $6$, and give an example of each. \exlevel{Exploration}
  \item Prove that if $p$ is prime, then $(x+y)^p\congr x^p+y^p\pmod p$. Use induction to deduce $(x_1+\cdots+x_m)^p\congr x_1^p+\cdots+x_m^p$. \exlevel{Reasoning}
  \item If $n$ is composite, must there be an $a$ with $a^n\not\congr a\pmod n$? Test $n=341$ and $561$. This opens the door to pseudoprimes and Carmichael numbers. \exlevel{Challenge}
\end{enumerate}
\end{exercisebox}


\chapter{The reusable machine}
\label{chap:groups}
\chaptertag{A quick introduction to finite groups}

The necklace proof is vivid, but it was tailored to one symmetry. Group theory isolates the small collection of rules shared by rotations, invertible residues, permutations, and many other systems.

\section{One operation, four rules}

\begin{definition}
A \emph{group} is a set $G$ with an operation $\ast$ such that:
\begin{enumerate}
  \item $a\ast b\in G$ for all $a,b\in G$ (closure);
  \item $(a\ast b)\ast c=a\ast(b\ast c)$ (associativity);
  \item there is an identity $e$ with $e\ast a=a\ast e=a$;
  \item each $a$ has an inverse $a^{-1}$ with $a\ast a^{-1}=a^{-1}\ast a=e$.
\end{enumerate}
If also $a\ast b=b\ast a$, the group is \emph{abelian}.
\end{definition}

\begin{center}
\begin{tabular}{lll}
\toprule
Group & Operation & Identity / inverse \\
\midrule
$\Z/n\Z$ & addition & $0$ / $-a$ \\
$(\Z/n\Z)^\times$ & multiplication of units & $1$ / modular inverse \\
rotations of a regular $n$-gon & composition & no rotation / reverse rotation \\
permutations of a set & composition & do nothing / undo \\
\bottomrule
\end{tabular}
\end{center}

Not every familiar set works: all residues modulo $n$ under multiplication fail because nonunits have no inverse.

\section{Powers and cyclic subgroups}

For $g\in G$, the list
\[
e,g,g^2,g^3,\dots
\]
eventually repeats when $G$ is finite. The first positive $m$ with $g^m=e$ is the \emph{order} of $g$, written $\operatorname{ord}(g)$. The powers form the cyclic subgroup $\angles g$.

\begin{center}
\begin{tikzpicture}
\foreach \k in {0,...,5}{
  \node[circle,draw=teal,fill=mist,minimum size=8mm] (n\k) at ({90-60*\k}:1.65) {$g^{\k}$};
}
\foreach \k [evaluate=\k as \j using {int(mod(\k+1,6))}] in {0,...,5}{
  \draw[-{Stealth},navy!60] (n\k) to[bend left=12] (n\j);
}
\node at (0,0) {multiply by $g$};
\end{tikzpicture}

\small A cyclic group of order $6$: repeated multiplication traces a loop.
\end{center}

\section{Cosets and Lagrange's theorem}

If $H$ is a subgroup of $G$ and $g\in G$, the left coset
\[
gH=\{gh:h\in H\}
\]
is a translated copy of $H$. Two cosets are either equal or disjoint, and every coset has $\abs H$ elements because $h\mapsto gh$ is a bijection.

\begin{theorem}[Lagrange]
If $G$ is finite and $H\le G$, then $\abs H$ divides $\abs G$.
\end{theorem}

\begin{proof}
The distinct cosets of $H$ partition $G$ into equally sized blocks. If there are $r$ blocks, then $\abs G=r\abs H$.
\end{proof}

\begin{corollary}
For every $g$ in a finite group $G$, $\operatorname{ord}(g)\divides\abs G$ and hence $g^{\abs G}=e$.
\end{corollary}

\section{Fermat's theorem, now in one line}

For prime $p$, every nonzero residue has an inverse, so
\[
(\Z/p\Z)^\times=\{1,2,\dots,p-1\}
\]
is a group of order $p-1$. Lagrange immediately yields the standard group-theoretic proof:
\[
a^{p-1}\congr1\pmod p\qquad(p\ndivides a).
\]

The gain is not a shorter proof at any price. The gain is a theorem that can be reused.

\begin{theorem}[Euler]
If $\gcd(a,n)=1$, then
\[
a^{\varphi(n)}\congr1\pmod n,
\]
where $\varphi(n)=\abs{(\Z/n\Z)^\times}$.
\end{theorem}

\begin{proof}
Apply the corollary of Lagrange to the group $(\Z/n\Z)^\times$.
\end{proof}

\begin{bigidea}
The necklace proof uses orbits of a cyclic action; Lagrange uses cosets of a subgroup. Both prove divisibility by partitioning a finite set into equal blocks. Group theory turns that shared move into a portable language.
\end{bigidea}

\section{A second standard proof without naming Lagrange}

If $a\not\congr0\pmod p$, multiplication by $a$ permutes the nonzero residues:
\[
a,2a,\dots,(p-1)a\pmod p
\]
is just $1,2,\dots,p-1$ in another order. Multiply all entries:
\[
a^{p-1}(p-1)!\congr(p-1)!\pmod p.
\]
Since $(p-1)!$ is invertible, cancel it. This is the group proof with the machinery folded away.

\begin{exercisebox}
\begin{enumerate}
  \item List $(\Z/12\Z)^\times$ and make its multiplication table. Find the order of every element. \exlevel{Fluency}
  \item Find $\varphi(20)$ and use Euler's theorem to reduce $3^{2026}$ modulo $20$. \exlevel{Fluency}
  \item Prove that the identity of a group is unique and that each element has a unique inverse. \exlevel{Reasoning}
  \item Prove that two left cosets of $H$ are equal or disjoint. \exlevel{Reasoning}
  \item Show that a group of prime order is cyclic. \exlevel{Reasoning}
  \item Draw the subgroup lattice of $(\Z/12\Z,+)$. Check Lagrange's theorem at every node. \exlevel{Exploration}
  \item Prove $\varphi(mn)=\varphi(m)\varphi(n)$ when $\gcd(m,n)=1$ by using the Chinese remainder theorem. \exlevel{Challenge}
  \item Compare the necklace, permutation, and Lagrange proofs of Fermat's theorem. For each, identify the set being partitioned or permuted and the exact place primality enters. \exlevel{Synthesis}
\end{enumerate}
\end{exercisebox}

\chapter{Arithmetic in Pascal's triangle}
\chaptertag{Binomial coefficients and divisibility}

\begin{definition}
For $0\le k\le n$,
\[
\binom nk=\frac{n!}{k!(n-k)!}.
\]
It counts the $k$-element subsets of an $n$-element set.
\end{definition}

The counting meaning proves integrality without asking the fraction to look integral. It also makes identities visible.

\begin{proposition}[Pascal and symmetry]
\[
\binom nk=\binom n{n-k},\qquad
\binom nk=\binom{n-1}{k-1}+\binom{n-1}{k}.
\]
\end{proposition}

For Pascal's identity, distinguish $k$-subsets that contain a chosen element from those that do not.

\begin{center}
\begin{tikzpicture}[x=.92cm,y=.70cm]
\foreach \n in {0,...,6}{
  \foreach \k in {0,...,\n}{
    \pgfmathtruncatemacro{\val}{factorial(\n)/(factorial(\k)*factorial(\n-\k))}
    \node[font=\small] (p\n-\k) at ({\k-\n/2},-\n) {$\val$};
  }
}
\draw[teal,thick,rounded corners] ($(p4-1)+(-.3,.28)$) rectangle ($(p4-1)+(.3,-.28)$);
\draw[teal,thick,rounded corners] ($(p4-2)+(-.3,.28)$) rectangle ($(p4-2)+(.3,-.28)$);
\draw[coral,thick,rounded corners] ($(p5-2)+(-.34,.28)$) rectangle ($(p5-2)+(.34,-.28)$);
\draw[-{Stealth},teal] (p4-1)--(p5-2);
\draw[-{Stealth},teal] (p4-2)--(p5-2);
\node[coral,right=15mm of p5-2] {$4+6=10$};
\end{tikzpicture}
\end{center}

\section{Three ways to see the binomial theorem}

\begin{theorem}[Binomial theorem]
\[
(x+y)^n=\sum_{k=0}^n\binom nk x^{n-k}y^k.
\]
\end{theorem}

In expanding $(x+y)^n$, choose which $k$ of the $n$ factors contribute $y$. This proof also explains why the coefficients are integers and why their sum is $2^n$.

\begin{proposition}[Vandermonde]
For nonnegative integers $r,s,n$,
\[
\sum_k\binom rk\binom s{n-k}=\binom{r+s}{n}.
\]
\end{proposition}

Count $n$-person committees from two groups: either all at once, or according to how many come from the first group.

\section{Prime rows}

\begin{theorem}
If $p$ is prime and $1\le k\le p-1$, then
\[
p\divides\binom pk.
\]
\end{theorem}

\begin{proof}
The identity $k\binom pk=p\binom{p-1}{k-1}$ shows $p$ divides $k\binom pk$. Since $p\ndivides k$, Euclid's lemma gives $p\divides\binom pk$.
\end{proof}

Thus the $p$th row of Pascal's triangle becomes $1,0,\dots,0,1$ modulo $p$, which again yields
\[
(x+y)^p\congr x^p+y^p\pmod p.
\]

\section{Measuring divisibility one prime at a time}

\begin{definition}
For nonzero integer $m$, $v_p(m)$ is the exponent of $p$ in its prime factorization. Thus $v_2(40)=3$.
\end{definition}

Valuations turn multiplication into addition:
\[
v_p(ab)=v_p(a)+v_p(b),\qquad
v_p\!\left(\frac ab\right)=v_p(a)-v_p(b)
\]
when the quotient is an integer.

\begin{theorem}[Legendre]
\[
v_p(n!)=\floor{\frac np}+\floor{\frac n{p^2}}+\floor{\frac n{p^3}}+\cdots.
\]
\end{theorem}

\begin{proof}
Among $1,\dots,n$, exactly $\floor{n/p}$ contribute at least one $p$; $\floor{n/p^2}$ contribute one additional $p$; and so on.
\end{proof}

Consequently,
\[
v_p\binom nk=\sum_{j\ge1}\left(\floor{\frac n{p^j}}-\floor{\frac k{p^j}}-\floor{\frac{n-k}{p^j}}\right).
\]
Each parenthesis is $0$ or $1$: it records whether adding $k$ and $n-k$ creates a carry at the $p^j$ place.

\begin{theorem}[Kummer's carry rule]
$v_p\binom nk$ equals the number of carries when $k$ and $n-k$ are added in base $p$.
\end{theorem}

The formula above is already the proof once carries are read from right to left. It gives a remarkable visual test: $\binom nk$ is not divisible by $p$ precisely when subtracting $k$ from $n$ in base $p$ needs no borrowing.

\begin{example}
To find $v_2\binom{10}{3}$, add $3=0011_2$ and $7=0111_2$:
\[
\begin{array}{r}0011\\+0111\\\hline1010\end{array}
\]
There are three carries, so $v_2\binom{10}{3}=3$. Indeed, $\binom{10}{3}=120=2^3\cdot15$.
\end{example}

\begin{labbox}
Color entry $\binom nk$ when it is nonzero modulo $2$, then modulo $3$. Zoom out. The self-similar patterns are the visual shadow of base-$p$ digits and lead to Lucas's theorem.
\end{labbox}

\begin{exercisebox}
\begin{enumerate}
  \item Prove $\sum_k\binom nk=2^n$ and $\sum_k(-1)^k\binom nk=0$ for $n>0$. \exlevel{Fluency}
  \item Compute $v_2(100!)$, $v_5(100!)$, and the number of trailing zeros of $100!$. \exlevel{Fluency}
  \item Give a counting proof of $k\binom nk=n\binom{n-1}{k-1}$. \exlevel{Reasoning}
  \item Use Vandermonde's identity to prove $\sum_k\binom nk^2=\binom{2n}{n}$. \exlevel{Reasoning}
  \item Determine all $n$ for which every interior entry of row $n$ is even. Use Kummer's rule. \exlevel{Reasoning}
  \item For a fixed prime $p$, explore the first $p^2$ rows of Pascal's triangle modulo $p$. State a digit-based conjecture for when $\binom nk\not\congr0\pmod p$. \exlevel{Exploration}
  \item Prove that $\operatorname{lcm}(1,2,\dots,2n)$ is divisible by $\binom{2n}{n}$ by comparing $p$-adic valuations. \exlevel{Bridge}
\end{enumerate}
\end{exercisebox}


\chapter{How crowded are the primes?}
\chaptertag{Chebyshev's estimate from the middle of Pascal's triangle}

Euclid proved that the primes never end. A sharper question asks how many lie below $x$:
\[
\pi(x)=\#\{p\le x:p\text{ prime}\}.
\]
The prime number theorem says $\pi(x)\sim x/\log x$, but its proof lies beyond this course. Chebyshev's elementary argument already captures the correct scale.

\section{The central binomial squeeze}

The largest coefficient in $(1+1)^{2n}$ is central. Since there are $2n+1$ nonnegative coefficients with sum $4^n$,
\begin{equation}\label{eq:central-squeeze}
\frac{4^n}{2n+1}\le\binom{2n}{n}<4^n.
\end{equation}
This elementary inequality contains surprising information about primes.

Define two weighted prime counts:
\[
\vartheta(x)=\sum_{p\le x}\log p,
\qquad
\psi(x)=\sum_{p^k\le x}\log p.
\]
The first logs the product of primes at most $x$; the second logs the least common multiple:
\[
\psi(x)=\log\operatorname{lcm}(1,2,\dots,\floor x).
\]

\begin{center}
\begin{tikzpicture}[x=.075cm,y=1cm]
  \draw[-{Stealth},navy] (0,0)--(105,0) node[right] {$x$};
  \foreach \p in {2,3,5,7,11,13,17,19,23,29,31,37,41,43,47,53,59,61,67,71,73,79,83,89,97}{
    \pgfmathsetmacro{\h}{.35+ln(\p)/2.3}
    \draw[teal,line width=1.2pt] (\p,0)--(\p,\h);
  }
  \draw[coral,decorate,decoration={brace,amplitude=5pt}] (50,-.28)--(100,-.28)
    node[midway,below=6pt] {primes in $(n,2n]$ divide $\binom{2n}{n}$};
  \node[muted,above] at (50,2.45) {A logarithmic weight turns products into sums.};
\end{tikzpicture}
\end{center}

\section{An upper bound for the prime product}

Every prime $n<p\le2n$ occurs in the numerator $(n+1)(n+2)\cdots(2n)$ and not in $n!$. Therefore it divides $\binom{2n}{n}$, and
\[
\vartheta(2n)-\vartheta(n)
\le\log\binom{2n}{n}
<2n\log2.
\]
Adding these inequalities over dyadic intervals gives the following deliberately non-optimized bound.

\begin{lemma}\label{lem:theta-upper}
For $x\ge2$,
\[
\vartheta(x)\le 2(\log2)x.
\]
\end{lemma}

\begin{proof}
Split $(1,x]$ into $(x/2,x],(x/4,x/2],\dots$ and apply the preceding estimate, with integer endpoints rounded outward. The resulting geometric series is bounded by
\[
x\log2+\frac{x}{2}\log2+\frac{x}{4}\log2+\cdots=2x\log2.
\]
Endpoint rounding can only add primes already included in the neighboring interval.
\end{proof}

\section{A lower bound from an lcm}

The bridge exercise in the last chapter showed
\[
\binom{2n}{n}\divides\operatorname{lcm}(1,2,\dots,2n).
\]
Taking logarithms and using \eqref{eq:central-squeeze},
\begin{equation}\label{eq:psi-lower}
\psi(2n)\ge\log\binom{2n}{n}\ge2n\log2-\log(2n+1).
\end{equation}

Prime powers make $\psi$ larger than $\vartheta$, but not by too much:
\[
\psi(x)=\vartheta(x)+\vartheta(x^{1/2})+\vartheta(x^{1/3})+\cdots.
\]
By Lemma~\ref{lem:theta-upper}, all terms after the first contribute at most
\[
2\sqrt{x}\log x
\]
for $x\ge4$: there are at most $\log_2x$ terms, and every argument is at most $\sqrt x$. Combining this with \eqref{eq:psi-lower} shows that, once $x$ is sufficiently large,
\begin{equation}\label{eq:theta-lower}
\vartheta(x)\ge \frac{\log2}{2}x.
\end{equation}
The constant $1/2$ is chosen for convenience; better bookkeeping improves it.

\section{From weighted primes to ordinary counting}

The lower bound is immediate because every term in $\vartheta(x)$ is at most $\log x$:
\[
\vartheta(x)\le\pi(x)\log x
\quad\Longrightarrow\quad
\pi(x)\ge\frac{\log2}{2}\frac{x}{\log x}.
\]
For the upper bound, primes between $\sqrt x$ and $x$ each contribute at least $\tfrac12\log x$:
\[
\vartheta(x)\ge\bigl(\pi(x)-\pi(\sqrt x)\bigr)\frac{\log x}{2}.
\]
Using $\pi(\sqrt x)\le\sqrt x$ and Lemma~\ref{lem:theta-upper} gives
\[
\pi(x)\le4\log2\,\frac{x}{\log x}+\sqrt x
\le4\frac{x}{\log x}
\]
for all sufficiently large $x$.

\begin{theorem}[Chebyshev estimate]
For all sufficiently large $x$,
\[
\boxed{
\frac{\log2}{2}\frac{x}{\log x}
\le \pi(x)\le
4\frac{x}{\log x}}
\]
In particular, $\pi(x)\asymp x/\log x$: their ratio stays between two positive constants.
\end{theorem}

\begin{warningbox}
The symbol $\asymp$ does not mean $\sim$. Chebyshev's estimate gets the correct order of magnitude. The prime number theorem makes the much stronger claim
\[
\lim_{x\to\infty}\frac{\pi(x)}{x/\log x}=1.
\]
\end{warningbox}

\begin{bigidea}
Why did binomial coefficients see primes? Factorials contain every integer up to a boundary. Dividing $(2n)!$ by two copies of $n!$ cancels most small-scale information while exposing primes that cross the boundary from $n$ to $2n$. The central coefficient is both arithmetically rich and analytically easy to estimate.
\end{bigidea}

\begin{exercisebox}
\begin{enumerate}
  \item Compute $\pi(10^k)$ and $10^k/\log(10^k)$ for $k=1,2,3,4$. Track the ratio. \exlevel{Fluency}
  \item Prove both sides of the central binomial squeeze \eqref{eq:central-squeeze}. \exlevel{Reasoning}
  \item Prove that every prime $p$ with $n<p\le2n$ divides $\binom{2n}{n}$ exactly once. \exlevel{Reasoning}
  \item Fill in the valuation proof that $\binom{2n}{n}$ divides $\operatorname{lcm}(1,\dots,2n)$. \exlevel{Reasoning}
  \item Starting from the dyadic inequality, write a fully detailed proof of Lemma~\ref{lem:theta-upper}, including floors and ceilings. \exlevel{Challenge}
  \item Use a computer to graph $\pi(x)$, $x/\log x$, and the two Chebyshev bounds. Which statement is visible at moderate $x$, and which remains too coarse? \exlevel{Exploration}
  \item Show that Chebyshev's estimate implies the $n$th prime $p_n$ has order of magnitude $n\log n$. \exlevel{Synthesis}
\end{enumerate}
\end{exercisebox}

\chapter{Capstone studio}
\chaptertag{Choose a thread, extend it, and explain it}

The capstone is not a miniature research paper. It is a carefully explained mathematical journey: a motivating question, small examples, one substantial proof, and a clear account of how course ideas connect.

\section{Project menu}

\subsection*{A. Rational approximation and Pell equations}
Choose a nonsquare $d$. Compute the continued fraction of $\sqrt d$, identify convergents that solve $x^2-dy^2=\pm1$, and prove at least one recurring pattern. Include a number-line or lattice illustration.

\subsection*{B. Symmetry, orbits, and divisibility}
Generalize necklace counting to length $n$. Classify possible orbit sizes for one composite $n$, explain stabilizers in concrete language, and derive a divisibility identity. M\"obius inversion is an optional extension, not a prerequisite.

\subsection*{C. Pascal's triangle modulo a prime}
Visualize Pascal's triangle modulo $p$. Develop and prove a digit rule (Lucas's theorem is the natural destination), then explain how Kummer's carry rule and the visual pattern encode the same arithmetic.

\subsection*{D. Prime-counting laboratory}
Compare $\pi(x)$ with $x/\log x$, $\operatorname{Li}(x)$, and the Chebyshev bounds over a range your computer can handle. Explain which plots are evidence and which claims are proved. Reconstruct one link in Chebyshev's proof in full detail.

\section{Required deliverables}

\begin{enumerate}
  \item A four-to-six page mathematical exposition.
  \item At least one original diagram or data visualization.
  \item One theorem stated precisely and proved completely.
  \item A one-page ``proof map'' naming the definitions and earlier results used.
  \item A five-minute oral explanation, followed by questions.
\end{enumerate}

\section{Evaluation rubric}

\begin{center}
\small
\begin{tabularx}{\textwidth}{p{.22\textwidth}Xr}
\toprule
Criterion & Evidence of strong work & Weight \\
\midrule
Mathematical correctness & Claims are true, hypotheses are visible, examples are checked, and the main proof closes every logical gap. & 35\% \\
Exposition & The question motivates the definitions; notation is stable; paragraphs explain why steps matter. & 25\% \\
Connection & At least two course strands are linked in a mathematically meaningful way. & 15\% \\
Visual or computational insight & A diagram or computation reveals structure and is interpreted rather than merely displayed. & 15\% \\
Revision and response & The final version addresses peer/instructor feedback and the oral defense shows ownership. & 10\% \\
\bottomrule
\end{tabularx}
\end{center}

\begin{checkpoint}
Before submitting, underline every sentence that makes a mathematical claim. For each underline, ask: Is it a definition, an observed example, a cited theorem, or a claim I have proved? Ambiguous status is often the real source of unclear exposition.
\end{checkpoint}

\section{A final synthesis prompt}

In 500--700 words, respond to this statement:

\begin{quote}
Number theory repeatedly converts multiplication into something easier to organize: remainders, exponent vectors, additive logarithms, or equal-sized orbits.
\end{quote}

Use one precise example from three different chapters. Your response should include at least one equation and explain not only what the conversion does, but why it helps a proof.



\appendix
\chapter{Two-by-two matrices and M\"obius maps}\label{app:matrices}
\chaptertag{A compact toolkit for continued fractions}

This appendix develops only the matrix algebra used in Chapter~3. The central
idea is simple: a matrix with determinant $\pm1$ can be reversed without ever
leaving the integers.

\section{Multiplication and determinants}

A $2\times2$ matrix acts on a column vector by
\[
 \begin{pmatrix}a&b\\c&d\end{pmatrix}
 \begin{pmatrix}x\\y\end{pmatrix}
 =\begin{pmatrix}ax+by\\cx+dy\end{pmatrix}.
\]
Matrix multiplication is defined so that composition works as expected:
applying $B$ and then $A$ is the same as applying $AB$. Explicitly,
\[
 \begin{pmatrix}a&b\\c&d\end{pmatrix}
 \begin{pmatrix}e&f\\g&h\end{pmatrix}
 =\begin{pmatrix}ae+bg&af+bh\\ce+dg&cf+dh\end{pmatrix}.
\]

The determinant of $A=\begin{psmallmatrix}a&b\\c&d\end{psmallmatrix}$ is
$\det A=ad-bc$. It measures invertibility:
\[
 A^{-1}=\frac1{ad-bc}
 \begin{pmatrix}d&-b\\-c&a\end{pmatrix},
 \qquad\text{provided }ad-bc\ne0.
\]
The identity $\det(AB)=\det(A)\det(B)$ makes determinants especially useful
for long products.

\begin{definition}
The group $\operatorname{GL}_2(\Z)$ consists of the integer matrices with
determinant $\pm1$. Its subgroup $\operatorname{SL}_2(\Z)$ consists of those
with determinant $+1$.
\end{definition}

\begin{proposition}[Integer reversibility]
An integer $2\times2$ matrix has an inverse with integer entries if and only if
its determinant is $\pm1$.
\end{proposition}

\begin{proof}
If $\det A=\pm1$, the displayed inverse formula has integer entries. Conversely,
if both $A$ and $A^{-1}$ have integer entries, then
$1=\det(I)=\det(A)\det(A^{-1})$ is a product of two integers. Hence
$\det(A)=\pm1$.
\end{proof}

\begin{bigidea}
The columns of a matrix in $\operatorname{GL}_2(\Z)$ form an integer basis of
$\Z^2$: every integer vector has a unique expression as an integer linear
combination of those columns. This is the exact algebraic strength supplied by
the determinant identity for consecutive convergents.
\end{bigidea}

\section{Matrices as fractional linear maps}

Associate to an invertible matrix
$A=\begin{psmallmatrix}a&b\\c&d\end{psmallmatrix}$, with $ad-bc\ne0$,
the M\"obius map
\[
  T_A(t)=\frac{at+b}{ct+d}.
\]
A direct substitution shows
\[
  T_A\bigl(T_B(t)\bigr)=T_{AB}(t).
\]
Thus nested fractional expressions become matrix products. Multiplying $A$ by
any nonzero scalar leaves $T_A$ unchanged, but retaining the specific integer
matrix remembers determinant and divisibility information that the function
alone would forget.

The inverse map is $T_{A^{-1}}$, so a change of variable transports fixed
points by conjugation:
\[
 x=T_N(x),\quad y=T_M(x)
 \quad\Longrightarrow\quad y=T_{MNM^{-1}}(y).
\]
Section~\ref{sec:periodic-fixed-points} uses this identity to prove that
adding a finite prefix to a periodic continued fraction preserves
quadratic irrationality. Section~\ref{sec:galois} then proves the precise
criterion for pure periodicity.

When $ct+d\ne0$, the quotient rule gives
\[
  T_A'(t)=\frac{ad-bc}{(ct+d)^2}.
\]
Therefore the determinant tells us whether the map is increasing or decreasing.
This is the M\"obius derivative trick used to locate alternating convergents.

\section{Continued fractions as iterated products}

One partial quotient corresponds to
\[
  D(a)=\begin{pmatrix}a&1\\1&0\end{pmatrix},
  \qquad T_{D(a)}(t)=a+\frac1t.
\]
Consequently,
\[
 T_{D(a_0)\cdots D(a_k)}(t)
 =[a_0;a_1,\dots,a_k,t].
\]
The convergent recurrences are equivalent to the product identity
\[
 D(a_0)\cdots D(a_k)
 =\begin{pmatrix}p_k&p_{k-1}\\q_k&q_{k-1}\end{pmatrix}.
\]
Since every $D(a_i)$ has determinant $-1$, the product has determinant
$(-1)^{k+1}$. Hence
\[
 p_kq_{k-1}-p_{k-1}q_k=(-1)^{k+1}.
\]

This one product packages four facts:
\begin{itemize}
  \item its first column is the current convergent $(p_k,q_k)$;
  \item its second column is the preceding convergent $(p_{k-1},q_{k-1})$;
  \item its determinant says the two columns form an integer basis;
  \item its M\"obius action inserts an arbitrary remaining tail.
\end{itemize}

\begin{example}
For $[4;2,6,7]=415/93$,
\[
 D(4)D(2)D(6)D(7)
 =\begin{pmatrix}415&58\\93&13\end{pmatrix}.
\]
Its inverse is
\[
 \begin{pmatrix}13&-58\\-93&415\end{pmatrix}.
\]
The top row applied to the first column gives
$13(415)-58(93)=1$. In the opposite direction, multiplying a vector by the
inverse recovers its integer coordinates in the basis of consecutive
convergents.
\end{example}

\begin{exercisebox}
\begin{enumerate}
  \item Multiply three general matrices $D(a_0)D(a_1)D(a_2)$ and identify the
  two convergents in its columns. \exlevel{Fluency}
  \item Prove $\det(AB)=\det(A)\det(B)$ directly for $2\times2$ matrices.
  \exlevel{Reasoning}
  \item Find the inverse of
  $\begin{psmallmatrix}17&7\\12&5\end{psmallmatrix}$ and use it to express
  $(1,0)$ and $(0,1)$ as integer combinations of its columns.
  \exlevel{Fluency}
  \item Verify directly that $T_A\circ T_B=T_{AB}$. Explain why multiplying a
  matrix by $-1$ does not change its M\"obius map. \exlevel{Reasoning}
  \item Let $P=D(a_1)\cdots D(a_\ell)$. Show that a purely periodic tail
  $t=[\overline{a_1,\dots,a_\ell}]$ is a fixed point of $T_P$, and derive its
  quadratic equation from the four entries of $P$. \exlevel{Reasoning}
\end{enumerate}
\end{exercisebox}

\chapter{Streaming arithmetic with continued fractions}\label{app:gosper}
\chaptertag{Gosper's homographic and bihomographic algorithms}

Suppose a number is arriving one continued-fraction digit at a time.  Can we
transform it without first converting it to a decimal, and without waiting for
the entire infinite input?  Bill Gosper's answer is yes.  The method reads just
enough of the input to prove the next output digit, emits that digit, and then
continues.

\section{The problem---before the machinery}

There are two versions.

\begin{description}[leftmargin=2.7cm,style=nextline]
  \item[One input.] Given
  \[
    x=[a_0;a_1,a_2,\dots]
  \]
  as a stream, and a linear fractional transformation
  \[
    L(x)=\frac{\alpha x+\beta}{\gamma x+\delta},
  \]
  produce the continued fraction of $L(x)$ from left to right.

  \item[Two inputs.] Given continued-fraction streams for $x$ and $y$, produce
  the continued fraction of a two-variable fractional expression.  In
  particular, $x+y$, $x-y$, $xy$, and $x/y$ all have the form
  \[
    \Phi(x,y)=\frac{axy+bx+cy+d}{exy+fx+gy+h}.
  \]
\end{description}

The intended computation is \emph{exact}.  It should never guess a digit from
a floating-point approximation.  Instead, every emitted digit comes with a
proof that all inputs compatible with the digits seen so far would produce that
same digit.

\begin{bigidea}
A finite input prefix determines an \emph{interval} of possible values.  Map
that interval through the desired transformation.  If the whole image fits
inside one continued-fraction digit cell $[k,k+1)$, the next output digit is
certified to be $k$.  If it crosses an integer boundary, read more input and
shrink the interval.
\end{bigidea}

The original algorithm appeared as Item 101B of the 1972 MIT memo
\emph{HAKMEM}.  A modern refinement by Michael J. Collins allows exact
intervals themselves to appear in the output stream, resolving boundary cases
in which no next digit can be certified after finite time.

\section{What the unread tail means}

The subscript on a tail records the \emph{next} digit, not the first digit of
the original number.  After the prefix $a_0,\dots,a_{j-1}$ has been read, write
\begin{equation}
  x=[a_0;a_1,\dots,a_{j-1},t_j],
  \qquad
  t_j=[a_j;a_{j+1},\dots].                         \label{eq:tail-invariant}
\end{equation}
Here $t_j$ is the complete unread remainder of the continued fraction, used as
the last entry in the first expression.  Before $a_j$ is known, the only
universal information is
\[
  1\leq t_j\leq\infty
\]
for a positive simple continued fraction.  When the next digit $a_j=m$ arrives,
\begin{equation}
  t_j=m+\frac1{t_{j+1}},
  \qquad t_{j+1}=[a_{j+1};a_{j+2},\dots]\in[1,\infty].
  \label{eq:tail-shift}
\end{equation}
Thus knowing the single digit $m$ restricts the old tail to
$t_j\in(m,m+1]$; absorbing $m$ replaces it by the new unknown tail
$t_{j+1}\in[1,\infty]$.

\begin{labbox}[title={Known prefix versus unknown tail}]
At the start, $t_0=x$ and no input digits have been absorbed.  After $a_0$ is
read, the state incorporates $a_0$ and the unknown variable is
$t_1=[a_1;a_2,\dots]$.  After $a_1$ is read, the state incorporates
$a_0,a_1$ and the unknown variable is $t_2$, and so on.  The state changes;
the exact number represented does not.
\end{labbox}

Let
\[
  D(m)=\begin{pmatrix}m&1\\1&0\end{pmatrix},
  \qquad
  M_0=\begin{pmatrix}\alpha&\beta\\\gamma&\delta\end{pmatrix}.
\]
After $j$ input digits have been absorbed, set
\begin{equation}
  M_j=M_0D(a_0)D(a_1)\cdots D(a_{j-1})
      =\begin{pmatrix}A&B\\C&D\end{pmatrix}.       \label{eq:prefix-state}
\end{equation}
Then the output value still waiting to be written is exactly
\[
  F_{M_j}(t_j)=\frac{At_j+B}{Ct_j+D}.
\]
This is the invariant behind the algorithm.

\section{The geometry: map an interval, then inspect its floor}

For a general certified tail interval $I=[b_0,b_1]$, consider
\[
  F_M(t)=\frac{At+B}{Ct+D}.
\]
Provided $Ct+D$ never vanishes on $I$, the derivative
\begin{equation}
  F_M'(t)=\frac{AD-BC}{(Ct+D)^2}                    \label{eq:mobius-derivative-b}
\end{equation}
has constant sign.  Hence $F_M$ is monotone, and the image $F_M(I)$ is simply
the interval whose endpoints are $F_M(b_0)$ and $F_M(b_1)$, in the appropriate
order.  No point between them can escape those endpoint bounds.

\begin{center}
\begin{tikzpicture}[x=1cm,y=1cm]
  \draw[navy!55,line width=.8pt] (0,1.25)--(5.2,1.25);
  \fill[teal] (1,1.25) circle (2.5pt) node[below=5pt,text=ink] {$b_0$};
  \fill[teal] (4.2,1.25) circle (2.5pt) node[below=5pt,text=ink] {$b_1$};
  \draw[teal,line width=3pt] (1,1.25)--(4.2,1.25);
  \node[text=muted] at (2.6,1.75) {possible input tail $I$};

  \draw[-{Stealth[length=2.5mm]},coral,line width=1.2pt]
    (5.65,1.25) to[bend left=22] (7.1,1.25);

  \foreach \x/\lab in {7.4/$k$,9.0/$k+1$,10.6/$k+2$}{
    \draw[navy!35] (\x,.94)--(\x,1.58);
    \node[below=4pt,text=muted] at (\x,.94) {\lab};
  }
  \draw[gold,line width=5pt] (7.75,1.25)--(8.65,1.25);
  \fill[coral] (7.75,1.25) circle (2.4pt);
  \fill[coral] (8.65,1.25) circle (2.4pt);
  \node[text=muted] at (8.2,1.75) {$F_M(I)\subset[k,k+1)$};
\end{tikzpicture}

\small\color{muted}If both endpoint images lie in the same floor cell and
there is no pole between them, every possible output has floor $k$.
\end{center}

For the universal interval $[1,\infty]$, the two endpoint values are
\[
  F_M(1)=\frac{A+B}{C+D},
  \qquad
  F_M(\infty)=\frac AC.
\]
An endpoint at infinity may be open, and a finite continued fraction may end;
exact implementations retain this endpoint information.  The essential test is
unchanged:
\[
  F_M(I)\subseteq[k,k+1)quad\Longrightarrow\quad
  \text{the next output digit is }k.
\]
If the image crosses an integer boundary, or if a pole is possible, no digit is
yet justified.  We must refine the input interval.

\section{The two algebraic moves}

\paragraph{Absorb the next input digit.}
If $a_j=m$, substitute $t_j=m+1/t_{j+1}$ into the current map.  Clearing the
factor $t_{j+1}$ gives
\[
 F_M\left(m+\frac1{t_{j+1}}\right)
 =\frac{(Am+B)t_{j+1}+A}{(Cm+D)t_{j+1}+C}.
\]
Thus the exact state update is
\begin{equation}
  \boxed{\mathcal I_m(M)=M D(m)
  =\begin{pmatrix}Am+B&A\\Cm+D&C\end{pmatrix}.}    \label{eq:unary-absorb}
\end{equation}

\paragraph{Emit a certified output digit.}
If the image interval lies in $[k,k+1)$, write the digit $k$.  The unread part
of the output continued fraction is then $1/(F_M-k)$, represented by
\begin{equation}
  \boxed{\mathcal E_k(M)=
  \begin{pmatrix}0&1\\1&-k\end{pmatrix}M
  =\begin{pmatrix}C&D\\A-kC&B-kD\end{pmatrix}.}   \label{eq:unary-emit}
\end{equation}

\begin{checkpoint}[title={The homographic loop}]
Maintain the statement: ``the unread output is $F_M(t_j)$ for every
$t_j\in I$.''
\begin{enumerate}
  \item Map the current certified interval $I$ through $F_M$ by evaluating its
  endpoints.
  \item If the image lies in one floor cell $[k,k+1)$, emit $k$, apply
  $M\leftarrow\mathcal E_k(M)$, and replace the output interval by
  $1/(F_M(I)-k)$.
  \item Otherwise read $a_j=m$, use $t_j=m+1/t_{j+1}$, apply
  $M\leftarrow\mathcal I_m(M)$, and return to the universal interval
  $t_{j+1}\in[1,\infty]$.
  \item Repeat.  A common nonzero factor may be removed from the four matrix
  entries without changing the map.
\end{enumerate}
\end{checkpoint}

\section{Worked stream: the first digits of \texorpdfstring{$\pi/2$}{pi/2}}

Here $L(x)=x/2$, so
$M_0=\begin{psmallmatrix}1&0\\0&2\end{psmallmatrix}$, and the input is
\[
 \pi=[3;7,15,1,292,\dots].
\]
After reading $a_0=3$ and $a_1=7$, the unknown variable is
$t_2=[15;1,292,\dots]\in[1,\infty]$, and
\[
 M_2=M_0D(3)D(7)=\begin{pmatrix}22&3\\14&2\end{pmatrix}.
\]
The first row below maps the endpoint tails $1$ and $\infty$.  Each following
row applies \eqref{eq:unary-emit}; no further input digit is read.

\begin{center}
\small
\begin{tabular}{c c c}
\toprule
event & current matrix $M$ & image of the possible tail interval \\
\midrule
read $3,7$ & $\begin{psmallmatrix}22&3\\14&2\end{psmallmatrix}$
 & $[25/16,\,11/7)$: emit $1$ \\
emit $1$ & $\begin{psmallmatrix}14&2\\8&1\end{psmallmatrix}$
 & $(7/4,\,16/9]$: emit $1$ \\
emit $1$ & $\begin{psmallmatrix}8&1\\6&1\end{psmallmatrix}$
 & $[9/7,\,4/3)$: emit $1$ \\
emit $1$ & $\begin{psmallmatrix}6&1\\2&0\end{psmallmatrix}$
 & $(3,\,7/2]$: emit $3$ \\
\bottomrule
\end{tabular}
\end{center}

After these four emissions the state is
$\begin{psmallmatrix}2&0\\0&1\end{psmallmatrix}$: the unread output is
$2t_2$.  Only now do we read the next input digits, $15$ and $1$; absorbing
them maps the new unknown tail into $[31,32)$ and certifies $31$.  Therefore
\[
 \frac\pi2=[1;1,1,3,31,\dots].
\]
The calculation has used exact rational endpoints throughout.  The decimal
digits of $\pi$ never entered.

\section{Two inputs: rectangles replace intervals}

Now suppose prefixes of two inputs have been read.  Their unread tails, say
$u_j$ and $v_\ell$, range independently over certified intervals $I$ and $J$.
The unread output has the form
\begin{equation}
 B_{\mathcal B}(u,v)
 =\frac{auv+bu+cv+d}{euv+fu+gv+h},
 \qquad
 \mathcal B=(a,b,c,d;\,e,f,g,h).                  \label{eq:bihomographic}
\end{equation}
Thus the possible input is a rectangle $I\times J$, not a line segment.

\begin{center}
\begin{tikzpicture}[x=1cm,y=1cm]
  \fill[teal!10] (.6,.35) rectangle (4.2,2.45);
  \draw[teal,line width=1pt] (.6,.35) rectangle (4.2,2.45);
  \foreach \p/\lab/\pos in {
    {.6,.35}/{$(b_0,c_0)$}/{below left},
    {4.2,.35}/{$(b_1,c_0)$}/{below right},
    {.6,2.45}/{$(b_0,c_1)$}/{above left},
    {4.2,2.45}/{$(b_1,c_1)$}/{above right}}
    \fill[coral] (\p) circle (2.5pt) node[\pos,text=ink] {\lab};
  \node[text=muted] at (2.4,1.4) {$I\times J$};
  \draw[-{Stealth[length=2.5mm]},coral,line width=1.2pt]
    (4.85,1.4) to[bend left=18] (6.35,1.4);
  \draw[navy!55] (6.65,1.4)--(10.8,1.4);
  \draw[gold,line width=5pt] (7.35,1.4)--(9.85,1.4);
  \fill[coral] (7.35,1.4) circle (2.5pt);
  \fill[coral] (9.85,1.4) circle (2.5pt);
  \node[text=muted] at (8.6,1.95) {range bounded by corner images};
\end{tikzpicture}
\end{center}

Assume the denominator has no zero on the rectangle.  For a fixed $v$, the
map is linear fractional in $u$, so its extrema along that horizontal segment
occur at the two ends.  On either vertical edge it is linear fractional in
$v$, so the remaining extrema occur at the edge endpoints.  Consequently the
entire range is bounded by the images of the four corners.  In particular, for
$I=J=[1,\infty]$ these values are
\begin{equation}
 \frac ae,\qquad \frac{a+b}{e+f},\qquad
 \frac{a+c}{e+g},\qquad
 \frac{a+b+c+d}{e+f+g+h},                         \label{eq:four-corners}
\end{equation}
with the usual interpretation of infinite and open endpoints.

The four basic operations differ only in their starting coefficients:
\[
\begin{array}{c|c}
\text{desired output} & (a,b,c,d;\,e,f,g,h)\\ \hline
x+y &(0,1,1,0;\,0,0,0,1)\\
x-y &(0,1,-1,0;\,0,0,0,1)\\
xy  &(1,0,0,0;\,0,0,0,1)\\
x/y &(0,1,0,0;\,0,0,1,0)
\end{array}
\]

If the next $u$-digit is $m$, substitute $u=m+1/u'$; if the next $v$-digit is
$n$, substitute $v=n+1/v'$.  Collecting coefficients gives
\begin{align}
\mathcal I^u_m(\mathcal B)
 &=(ma+c,\,mb+d,\,a,\,b;\,me+g,\,mf+h,\,e,\,f),
                                                        \label{eq:binary-absorb-u}\\
\mathcal I^v_n(\mathcal B)
 &=(na+b,\,a,\,nc+d,\,c;\,ne+f,\,e,\,ng+h,\,g).
                                                        \label{eq:binary-absorb-v}
\end{align}
When the four corner images all lie in one floor cell $[k,k+1)$, emit $k$ and
invert the fractional part:
\begin{equation}
\mathcal E_k(\mathcal B)
=(e,f,g,h;\,a-ke,b-kf,c-kg,d-kh).                  \label{eq:binary-emit}
\end{equation}

The geometry and the logic are exactly as before.  If the image of the current
rectangle crosses an integer boundary, absorb from one input and replace the
rectangle by a thinner one; alternate inputs, or choose the one expected to
shrink the image more.  If the denominator may vanish inside the rectangle,
no output is safe until further input or separate algebra resolves the issue.

\begin{bigidea}
The four coefficients in the one-input problem, or eight in the two-input
problem, are bookkeeping for a moving interval image.  Matrix formulas make
the updates fast, but the proof of correctness is geometric: every possible
tail lies in the current interval or rectangle, and every emitted digit is
shared by its entire image.
\end{bigidea}

\section{The boundary problem and the interval repair}

The floor function jumps at every integer.  This creates a subtle failure mode.
Feed two \emph{independent} copies of the continued fraction of $\sqrt2$ to the
multiplication algorithm.  The actual product is $2$, but every finite pair of
independent enclosures allows one copy to move a little below $\sqrt2$ and the
other a little above it.  The resulting product rectangle can straddle $2$.
Its image then meets both floor cells $[1,2)$ and $[2,3)$, so the rule cannot
certify even the first digit.

This is an information problem, not numerical instability.  The two stream
variables have forgotten that they denote the same algebraic number.  Retaining
the relation $x=y$ reduces the expression to $x^2$, and the separate fact
$x^2=2$ resolves the equality immediately.

Collins's interval-valued refinement distinguishes a continued fraction from
its computational stream.  A stream item may be either
\begin{itemize}
  \item an explicit partial quotient, or
  \item an exact rational interval bounding the current unread value.
\end{itemize}
If no digit is forced, the algorithm may output a tighter certified interval
instead of waiting silently.  A consumer asks for an error tolerance
$\varepsilon$ and reads until the resulting value interval has width below
$\varepsilon$.  Exact endpoints and their open-or-closed status preserve the
no-roundoff character of Gosper's method while guaranteeing useful progress at
an integer boundary.

\begin{warningbox}[title={Digits and numerical accuracy are different requests}]
There may be no finite proof of the next continued-fraction digit even when a
very accurate numerical enclosure is easy.  A robust implementation therefore
asks for a certified error tolerance, not an unconditional number of output
digits.  Away from boundaries, intervals quickly collapse into ordinary
partial quotients and recover Gosper's original stream.
\end{warningbox}

\section{Exercises}

\begin{exercisebox}
\begin{enumerate}
  \item Suppose only the prefix $[2;3]$ of $x$ is known.  Write
  $x=[2;3,t_2]$ with $t_2\in[1,\infty]$, evaluate the two endpoint values,
  and prove that $x\in[9/4,7/3)$. \exlevel{Fluency}

  \item Prove the endpoint-image rule using
  \eqref{eq:mobius-derivative-b}.  Apply it to
  $F(t)=(22t+3)/(14t+2)$ on $[1,\infty]$ and certify its floor.
  \exlevel{Reasoning}

  \item Derive the absorption and emission updates
  \eqref{eq:unary-absorb} and \eqref{eq:unary-emit} directly by
  substitution.  Express each as matrix multiplication. \exlevel{Fluency}

  \item Continue the $\pi/2$ calculation: from
  $\begin{psmallmatrix}2&0\\0&1\end{psmallmatrix}$, absorb $15$ and $1$ and
  show that the new image lies in $[31,32)$. \exlevel{Reasoning}

  \item Prove the four-corner range rule for a bihomographic map on a finite
  rectangle with no denominator zero.  Why is checking four corners not enough
  if a pole lies inside? \exlevel{Reasoning}

  \item Derive \eqref{eq:binary-absorb-u} and
  \eqref{eq:binary-absorb-v}.  Verify that the four initial states really
  represent $+$, $-$, $\times$, and $/$. \exlevel{Fluency}

  \item Explain why independent intervals for two copies of $\sqrt2$ can
  straddle the boundary $2$ forever when computing $\sqrt2\sqrt2$.  State the
  extra algebraic information that resolves the equality. \exlevel{Reasoning}

  \item Design a lazy implementation whose output type is either
  \texttt{Digit(k)} or \texttt{Bound(lo,hi)}.  State invariants that make every
  emitted item a certificate, then test all four operations on rational and
  quadratic inputs. \exlevel{Project}
\end{enumerate}
\end{exercisebox}

\chapter{Venus writes Euclid in the sky}
\chaptertag{Three clocks, five alignments, and the Dresden Codex}
\label{app:venus}

Continued fractions are useful whenever two repeating motions almost return to
their starting arrangement at the same time.  Venus gives a particularly
visible example.  Small Fibonacci numbers describe a near-repeat of the modern
planetary motions, while related integers organize an exact calendar table in
the Dresden Codex.  We will build the astronomical vocabulary from scratch and
keep the approximate observation separate from the exact calendar arithmetic.

\section{Begin with two runners on circular tracks}

Imagine Earth and Venus as two runners circling the Sun.  Venus uses the inner
track and completes a lap sooner, so it repeatedly catches and passes Earth.
At the instant shown below, the Sun, Venus, and Earth lie on the same ray.

\begin{center}
\begin{tikzpicture}[scale=1.15]
  \draw[navy!28,line width=.9pt] (0,0) circle (1.25);
  \draw[navy!28,line width=.9pt] (0,0) circle (2.15);
  \draw[navy!22,dashed] (0,0)--(2.55,0);
  \fill[gold] (0,0) circle (4pt) node[below left=3pt,text=ink] {Sun};
  \fill[coral] (1.25,0) circle (4pt) node[below=4pt,text=ink] {Venus};
  \fill[teal] (2.15,0) circle (4.5pt) node[above=4pt,text=ink] {Earth};
  \draw[-{Stealth[length=2.2mm]},coral,line width=1.1pt]
    (120:1.25) arc[start angle=120,end angle=200,radius=1.25];
  \draw[-{Stealth[length=2.2mm]},teal,line width=1.1pt]
    (120:2.15) arc[start angle=120,end angle=180,radius=2.15];
  \node[coral,align=center] at (-1.55,-1.45) {shorter lap\\faster runner};
  \node[teal,align=center] at (-2.45,1.65) {longer lap\\slower runner};
  \node[text=muted,align=center] at (4.0,.05)
    {an alignment\\(a conjunction)};
  \draw[-{Stealth[length=2mm]},muted] (3.25,.05)--(2.35,.05);
\end{tikzpicture}
\end{center}

Here are the three time intervals we need.

\begin{center}
\small
\begin{tabularx}{\textwidth}{>{\bfseries}p{.16\textwidth}p{.30\textwidth}X}
\toprule
Symbol & Everyday meaning & Astronomical name and approximate value \\
\midrule
$E$ & time for Earth to make one Sun-centered lap
& Earth's \emph{sidereal period}: $365.256$ days \\
$V$ & time for Venus to make one Sun-centered lap
& Venus's \emph{sidereal period}: $224.7$ days \\
$S$ & time from one Earth--Venus alignment to the next
& Venus's \emph{synodic period}: $583.92$ days \\
\bottomrule
\end{tabularx}
\end{center}

``Sidereal'' means that a full lap is measured relative to the distant stars.
A ``conjunction'' is an alignment; in the alignment drawn above, Venus lies
between Earth and the Sun, so astronomers call it an \emph{inferior
conjunction}.  ``Synodic period'' simply means the waiting time until the same
kind of alignment occurs again.\footnote{For the modern mean values, see NASA,
\href{https://eclipse.gsfc.nasa.gov/OH/transit12.html}{\emph{2012 Transit of
Venus}} and \href{https://history.arc.nasa.gov/hist_pdfs/nasa_sp461.pdf}{NASA
SP--461, \emph{Pioneer Venus}}.  The real orbits are neither perfectly circular
nor perfectly uniform; rounded mean periods are all we need here.}

\section{The overtaking clock}

Earth completes $1/E$ laps per day, while Venus completes $1/V$.  Because Venus
is faster, it gains on Earth at the rate
\[
  \frac1V-\frac1E\quad\hbox{laps per day}.
\]
The next alignment occurs when Venus has gained exactly one whole lap.  Hence
\begin{equation}
  \frac1S=\frac1V-\frac1E.                         \label{eq:venus-beat}
\end{equation}
This is the same ``lapping time'' calculation used for two runners on a track.

Put $r=E/V$: it compares the length of an Earth lap in time with a Venus lap.
Multiplying \eqref{eq:venus-beat} by $E$ gives
\begin{equation}
  \frac SE=\frac1{r-1}=:T(r).                      \label{eq:venus-mobius}
\end{equation}
The measured ratio $E/V\approx1.6255$ has the simple, very accurate fraction
$13/8$ nearby.  If we replace the measured ratio temporarily by this ideal one,
then
\[
  T\!\left(\frac{13}{8}\right)=\frac85.
\]
In plain language:
\begin{itemize}
  \item during about eight Earth laps, Venus makes about thirteen laps; and
  \item in that time Venus gains $13-8=5$ laps on Earth, producing about five
  repeat alignments.
\end{itemize}
None of these counts is exact for the physical planets.  Their usefulness is
that $13/8$ is an excellent approximation with a tiny denominator.

\begin{checkpoint}[title={Why the golden ratio appears}]
The positive fixed point of the overtaking map $T(r)=1/(r-1)$ satisfies
\[
  r=\frac1{r-1},\qquad r^2-r-1=0,
\]
so it is $\varphi=(1+\sqrt5)/2$.  More generally,
\[
  T\!\left(\frac{F_{n+1}}{F_n}\right)
  =\frac{F_n}{F_{n-1}}.
\]
The calculation of one lapping time walks backward through consecutive
Fibonacci ratios.
\end{checkpoint}

\section{Why five alignments draw a star}

During one alignment-to-alignment interval, Earth travels about
$S/E\approx8/5$ laps.  One complete lap returns to the same direction, so only
the leftover $3/5$ of a lap changes the direction of the alignment ray.  The
five successive directions are therefore approximately
\[
  0,\quad \frac35,\quad \frac15,\quad \frac45,\quad \frac25
  \pmod 1.
\]
Connecting those directions in time order draws a five-pointed star.

\begin{center}
\begin{tikzpicture}[scale=1.2]
  \coordinate (v0) at (90:2);
  \coordinate (v1) at (162:2);
  \coordinate (v2) at (234:2);
  \coordinate (v3) at (306:2);
  \coordinate (v4) at (18:2);
  \draw[navy!30,line width=.8pt] (0,0) circle (2);
  \draw[navy!20,dashed] (0,0) -- (v0);
  \draw[coral,line width=1.8pt]
    (v0) -- (v3) -- (v1) -- (v4) -- (v2) -- cycle;
  \fill[gold] (0,0) circle (3pt) node[below=4pt,text=muted] {Sun};
  \foreach \v/\lab/\pos in {
    v0/$C_0$/above,
    v3/$C_1$/below right,
    v1/$C_2$/above left,
    v4/$C_3$/above right,
    v2/$C_4$/below left}
    \fill[teal] (\v) circle (2.7pt) node[\pos,text=ink] {\lab};
\end{tikzpicture}

\small\color{muted}Each $C_j$ marks the direction of an alignment, not the
path traveled by Venus.  The actual fifth return misses the starting direction
slightly.
\end{center}

The geometry of a perfectly regular pentagram contains the golden ratio
exactly.  The planetary star is only approximate: because $13/8$ is not the
exact period ratio, the alignment points drift and the whole pattern slowly
turns over repeated eight-year cycles.\footnote{The European Southern
Observatory's educational note
\href{https://www.eso.org/public/outreach/eduoff/vt-2004/Education/edu1app5.html}{\emph{The
Pentagram of Venus}} illustrates the five alignments and the small miss after
eight years.}

\section{The Euclidean algorithm inside \texorpdfstring{$13:8$}{13:8}}

The approximation contains the short Euclidean chain
\[
\begin{aligned}
  13&=1\cdot8+5, & 8&=1\cdot5+3,\\
   5&=1\cdot3+2, & 3&=1\cdot2+1,\\
   2&=2\cdot1.
\end{aligned}
\]
Consequently
\[
  \frac{13}{8}=[1;1,1,1,2]=[1;1,1,1,1,1],
  \qquad
  \frac85=[1;1,1,2]=[1;1,1,1,1].
\]
These are consecutive convergents of
$\varphi=[1;\overline1]$.  The orbit ratio $13/8$, the difference $5$, the
step $3/5$, and the pentagram are not independent coincidences: they are
different consequences of one good rational approximation.

\section{The Dresden Codex: an exact calendar model}

The modern calculation above starts with measured periods and obtains
approximations.  The Maya Venus Table in the Dresden Codex instead uses chosen
whole-day counts for calendrical bookkeeping.  In particular:
\begin{itemize}
  \item $584$ days is the table's rounded alignment-to-alignment Venus cycle;
  \item $365$ days is one \emph{Haab}, the Maya solar-year calendar; and
  \item $260$ days is one Tzolk'in, a repeating named-day cycle.
\end{itemize}

Five pages of the table present one $584$-day Venus cycle apiece.  The five
cycles make an eight-Haab block because
\begin{equation}
  5\cdot584=8\cdot365=2920.                         \label{eq:dresden-block}
\end{equation}
This equality is exact arithmetic inside the calendar model, not an assertion
that the physical Venus cycle is exactly $584$ days.\footnote{See the Saxon
State and University Library Dresden's description of the
\href{https://www.slub-dresden.de/en/explore/manuscripts/the-dresden-maya-codex/content}{Dresden
Maya Codex}, especially the Venus pages 24 and 46--50.  For the page structure
and cumulative day counts, see Akira Kaneko,
\href{https://classic.scielo.org.mx/scielo.php?pid=S0185-25742010000100002&script=sci_arttext}{``Identification
of Postclassic Maya Constellations from the Venus Pages of the Dresden
Codex,''} \emph{Estudios de Cultura Maya} \textbf{35} (2010).}

To make the $2920$-day block agree again with the $260$-day named-day cycle,
take a least common multiple.  Since $\gcd(2920,260)=20$,
\[
  \operatorname{lcm}(2920,260)=\frac{2920\cdot260}{20}=37960.
\]
At that common return,
\begin{equation}
  65\cdot584=104\cdot365=146\cdot260
  =13\cdot2920=37960.                              \label{eq:dresden-long}
\end{equation}

\begin{labbox}[title={The same numbers are doing different jobs}]
In the modern orbital approximation, $13$ counts Venus's Sun-centered laps in
about eight Earth laps.  In \eqref{eq:dresden-long}, $13$ counts complete
$2920$-day calendar blocks in a $37960$-day table.  The identity $13=8+5$ is a
delightful number-theoretic lens, but the manuscript alone does not show that
Maya astronomers formulated their table using Fibonacci numbers or continued
fractions.
\end{labbox}

Finally, rounding has a cost.  The modern mean alignment cycle is about
$583.92$ days, so the calendar value $584$ gains roughly
\[
  65(584-583.92)\approx5.2\text{ days}
\]
over $65$ cycles.  Scholars study correction schemes in the table.  Our exact
identity \eqref{eq:dresden-long} describes its integer framework, not a
perfectly uniform planet.

\begin{exercisebox}
\begin{enumerate}
  \item In your own words, distinguish a planet's Sun-centered lap time from
  the waiting time between two Earth--planet alignments.  Which are $E$, $V$,
  and $S$? \exlevel{Fluency}

  \item Derive \eqref{eq:venus-beat} by comparing lap rates.  Then derive
  \eqref{eq:venus-mobius} and verify that $T(13/8)=8/5$.
  \exlevel{Reasoning}

  \item Starting from a step of $3/5$ of a turn, list the first five positions
  modulo $1$.  Why does joining them chronologically make a star rather than
  the boundary of a pentagon? \exlevel{Reasoning}

  \item Run the Euclidean algorithm on $(13,8)$.  Read its quotients as the
  continued fraction of $13/8$, and explain why replacing the final $2$ by
  $1,1$ gives the same rational number. \exlevel{Reasoning}

  \item Factor $2920$ and $260$, compute their gcd and lcm, and recover every
  multiplier in \eqref{eq:dresden-long}. \exlevel{Fluency}

  \item Prove $T(F_{n+1}/F_n)=F_n/F_{n-1}$.  Show that the positive fixed point
  of $T$ is $\varphi$, and connect this with $[1;\overline1]$.
  \exlevel{Reasoning}

  \item Using $583.92$ days as the modern mean alignment cycle, compare the
  elapsed time for $5$, $65$, and $650$ cycles with the nominal $584$-day
  count.  Which error is negligible for a sketch but serious for a long
  calendar? \exlevel{Exploration}
\end{enumerate}
\end{exercisebox}

\chapter{Selected hints and solutions}
\chaptertag{For checking arguments, not replacing them}

Only selected exercises are treated here. A useful routine is: write a complete attempt, mark the exact step where you are stuck, then consult only the first relevant sentence below.

\section*{Chapter 1}

\textbf{1.} Euclid gives
\[
2024=2(748)+528,\quad748=528+220,\quad528=2(220)+88,
\]
\[
220=2(88)+44,\quad88=2(44),
\]
so the gcd is $44$.

\textbf{2.} The divisions are $527=341+186$, $341=186+155$, $186=155+31$, $155=5\cdot31$. Reverse:
\[
31=186-155=2(186)-341=2(527-341)-341=2\cdot527-3\cdot341.
\]

\textbf{4.} One solution is $35(-5)+22(8)=1$. All solutions are
\[
x=-5+22t,\qquad y=8-35t\qquad(t\in\Z).
\]
To prove completeness, subtract two solutions and use $35(x-x')=-22(y-y')$ with $\gcd(35,22)=1$.

\section*{Chapter 2}

\textbf{3.} Write $c=ar=bs$. Since $a\divides bs$ and $\gcd(a,b)=1$, repeated Euclid's lemma (or B\'ezout) gives $a\divides s$. Thus $s=at$ and $c=abt$.

\textbf{4.} If $\sqrt p=a/b$ in lowest terms, then $a^2=pb^2$. The exponent of $p$ on the left is even; on the right it is odd. Unique factorization forbids this.

\textbf{6.} The primes at most $100$ are
\[
2,3,5,7,11,13,17,19,23,29,31,37,41,43,47,53,59,61,67,71,73,79,83,89,97.
\]
Every composite at most $100$ has a prime factor at most $\sqrt{100}=10$, so
sieving by the primes $2,3,5,7$ crosses out every composite.

\textbf{7.} No: $5!+1=121=11^2$. Euclid's number only needs a prime divisor outside the proposed list; it need not itself be prime.

\textbf{9.} We already know that $a=\sqrt2$ is irrational. By the definition
of the logarithm, $a^b=3$. If $b=p/q$ were rational with $q>0$, then
$(\sqrt2)^{p/q}=3$; raising both sides to $2q$ would give
$2^p=3^{2q}$, contrary to unique factorization. Hence $b$ is irrational.
Unlike the proof by cases, this argument names a particular pair and proves
that it works; it does not leave the witness dependent on an undecided
alternative.

\section*{Chapter 3}

\textbf{1.} Euclid gives
\[
355/113=[3;7,16],
\]
with convergents $3$, $22/7$, and $355/113$.

\textbf{2.} The convergents are
\[
1,\;2,\;\frac32,\;\frac53,\;\frac85,\;\frac{13}{8}.
\]
Numerators and denominators are consecutive Fibonacci numbers; the recurrence with every $a_k=1$ proves it.

\textbf{4.} Since $a_n=(a_n-1)+1/1$ at the bottom level,
\[
[a_0;\dots,a_n]=[a_0;\dots,a_n-1,1].
\]
Requiring the last partial quotient to exceed $1$ selects a unique finite expansion.

\textbf{6.} With $\alpha\approx0.2421988$, the continued fraction begins
\[
 \alpha=[0;4,7,1,3,5,\dots],
\]
so its first four nonzero convergents are
\[
 \frac14,\qquad\frac7{29},\qquad\frac8{33},\qquad\frac{31}{128}.
\]
The Julian calendar uses $1/4$ leap day per year. A century-based cycle has
$q=100Q$, so we instead approximate
\[
 100\alpha\approx24.21988=[24;4,1,1,4,\dots].
\]
Its first nonintegral convergent is $97/4$; dividing the denominator scale by
$100$ gives $97/400$, the Gregorian rate. The corresponding mean years are
\[
 365+\frac14=365.25,
 \qquad 365+\frac{97}{400}=365.2425,
 \qquad 365+\frac{31}{128}=365.2421875.
\]
Relative to the fixed approximation $365.2421988$, their annual errors in days
are about $7.80\times10^{-3}$, $3.01\times10^{-4}$, and
$-1.13\times10^{-5}$, respectively. Thus $31/128$ is numerically better,
whereas $97/400$ fits a convenient rule organized by fours and centuries.

\textbf{8.} Euclid gives $527/341=[1;1,1,5]$, whose convergents are
$1,2,3/2,17/11$. The penultimate convergent is $p_2/q_2=3/2$, so the
remainder formula gives
\[
  \gcd(527,341)=r_2=2\cdot527-3\cdot341=31.
\]

\textbf{9.} The first six convergents of $\sqrt3=[1;\overline{1,2}]$ are
\[
1,\quad 2,\quad \frac53,\quad \frac74,\quad \frac{19}{11},\quad \frac{26}{15}.
\]
They satisfy
\[
1<\frac53<\frac{19}{11}<\sqrt3
<\frac{26}{15}<\frac74<2.
\]
For even $k$, $M_k'(t)<0$, so the finite-tail value lies above
$M_k(\infty)=c_k$; for odd $k$, the derivative is positive and the inequality
reverses. Applying this once to $x$ and once to $c_{k+2}$ gives the two chains.

\textbf{10.} For $k=0$ the formula is just
$D(a_0)=\begin{psmallmatrix}p_0&p_{-1}\\q_0&q_{-1}\end{psmallmatrix}$.
If it holds at $k$, right multiplication by $D(a_{k+1})$ replaces the first
column by $a_{k+1}(p_k,q_k)^T+(p_{k-1},q_{k-1})^T$ and moves the old first
column into second place. These are exactly the recurrences. The determinant
is $(-1)^{k+1}$, so the inverse formula introduces no denominators.

\textbf{11.} For $1/7$, the remainders are
$1,3,2,6,4,5,1,\dots$, producing $0.\overline{142857}$. For $5/12$, they are
$5,2,8,8,\dots$, producing $0.41\overline6$. The first repeated remainder
marks the start and length of the repeating block.

\textbf{15.} Since
\[
D(1)D(2)=\begin{pmatrix}3&1\\2&1\end{pmatrix},
\]
the repeating tail obeys $t=(3t+1)/(2t+1)$. Hence
$2t^2-2t-1=0$, and positivity selects $t=(1+\sqrt3)/2$.

\textbf{16.} For $x=1+\sqrt2$, the conjugate $x'=1-\sqrt2$ lies in
$(-1,0)$; for $x=(1+\sqrt5)/2$, the conjugate is $(1-\sqrt5)/2$, again in
$(-1,0)$. Both numbers exceed $1$, so Galois's criterion applies. Indeed their
expansions are $[\overline2]$ and $[\overline1]$, respectively.

\textbf{17.} Factoring the difference of squares gives
\[
 \abs{\sqrt D-p/q}
 =\frac{\abs{D-p^2/q^2}}{\sqrt D+p/q}
 =\frac{1}{q^2(\sqrt D+p/q)}.
\]
The denominator factor $\sqrt D+p/q$ exceeds $2$, so Legendre's criterion
forces $p/q$ to be a convergent of $\sqrt D$.

\textbf{18.} We have $\sqrt3=[1;\overline{1,2}]$, of even period length $2$;
one period therefore gives the $+1$ solution $2^2-3(1)^2=1$. The first columns
of the first three powers of
$\begin{psmallmatrix}2&3\\1&2\end{psmallmatrix}$ are
\[
 \begin{pmatrix}2\\1\end{pmatrix},\qquad
 \begin{pmatrix}7\\4\end{pmatrix},\qquad
 \begin{pmatrix}26\\15\end{pmatrix},
\]
and each satisfies $p^2-3q^2=1$.

\textbf{19.} A finite computation verifies only the displayed finite averages;
an asymptotic conclusion requires control of every later partial quotient.
Khinchin's theorem excludes only a measure-zero exceptional set and supplies no
known test proving that a named constant such as $\pi$ lies outside it. For $e$,
the repeating pattern of positions gives
\[
 a_{3j-2}=1,\qquad a_{3j-1}=2j,\qquad a_{3j}=1,
\]
so $a_1\cdots a_{3m}=2^m m!$. Taking the $(3m)$th root and applying
$m!\sim\sqrt{2\pi m}(m/e)^m$ yields
\[
 G_{3m}(e)=(2^m m!)^{1/(3m)}\sim(2m/e)^{1/3}\to\infty.
\]
The denominator recurrence gives
$q_{3m}\geq a_1\cdots a_{3m}=2^m m!$, so $q_{3m}^{2/(3m)}\to\infty$ too.
Thus $e$ is provably on the exceptional side of both laws.

\textbf{20.} The two matrices and their maps are
\[
 D(1)D(2)=\begin{pmatrix}3&1\\2&1\end{pmatrix},
 \quad t\longmapsto\frac{3t+1}{2t+1};
 \qquad
 D(2)D(1)=\begin{pmatrix}3&2\\1&1\end{pmatrix},
 \quad t\longmapsto\frac{3t+2}{t+1}.
\]
The rightmost matrix acts first. The fixed-point equations are
$2t^2-2t-1=0$ and $t^2-2t-2=0$, with positive roots
$(1+\sqrt3)/2$ and $1+\sqrt3$, respectively.

\textbf{21.} For $j=1$ and $j=2$, respectively,
\[
 MNM^{-1}
 =\begin{pmatrix}2&3\\1&2\end{pmatrix},
 \qquad
 MNM^{-1}
 =\begin{pmatrix}3&2\\1&1\end{pmatrix}.
\]
The positive fixed points are $\sqrt3$ and $1+\sqrt3$.
Only the second is reduced, since $-\sqrt3<-1$ but
$-1<1-\sqrt3<0$. Its pure period is $[\overline{2,1}]$.

\textbf{22.} A pole $ct+d=0$ would make $t$ rational if $c\ne0$;
if $c=0$, invertibility implies $d\ne0$. If $T_M(t)=r\in\Q$,
then $(a-rc)t=rd-b$. The coefficient $a-rc$ cannot vanish:
together with the equality it would force $ad-bc=0$.
Thus $t=(rd-b)/(a-rc)$ would be rational, a contradiction.
The fixed-point assertion must exclude scalar matrices.
For example, $T_I(t)=t$ for every $t$, including $\sqrt[3]{2}$,
which is not quadratic: $X^3-2$ has no rational root and is an
irreducible cubic over $\Q$.

\textbf{23.} Precisely $1+\sqrt2$, $(1+\sqrt5)/2$, and
$(1+\sqrt3)/2$ are reduced, so those three are purely periodic.
The conjugate of $\sqrt2$ is less than $-1$, and the last two
numbers fail $x>1$. All six are real quadratic irrationals, so
Lagrange's theorem makes all their continued fractions eventually
periodic, including the one with negative initial digit.

\textbf{24.} The successive triples are
\[
 (0,1,3),\ (3,4,1),\ (1,3,1),\ (2,3,1),\
 (1,4,1),\ (3,1,6),\ (3,4,1).
\]
Thus $\sqrt{13}=[3;\overline{1,1,1,1,6}]$.
After the first step, $|P_n|\le3<\sqrt{13}$ and
$0<Q_n\le4<2\sqrt{13}$; each corresponding surd exceeds $1$
and has conjugate in $(-1,0)$.

\textbf{25.} Here $-1/z'=1+\sqrt3$, whose floor is $2$.
The predecessor is $2+1/z=1+\sqrt3$.
The bounds allow $P=-1,0,1$ and $Q=1,2,3$, but the
divisibility condition and both reduced inequalities leave only
$(P,Q)=(1,1),(1,2)$. They represent $1+\sqrt3$ and
$(1+\sqrt3)/2$, and the algorithm interchanges them.
A finite map can have an initial segment before its cycle:
$0\mapsto1$, $1\mapsto1$. Unique reduced predecessors exclude
that behavior in the continued-fraction argument.

\textbf{26.} The period matrix and equation are
\[
 N=\begin{pmatrix}10&3\\7&2\end{pmatrix},
 \qquad 7x^2-8x-3=0,\qquad x=\frac{4+\sqrt{37}}7.
\]
Transposition gives $3z^2-8z-7=0$, with positive root
$z=(4+\sqrt{37})/3$. Finally,
$-1/x'=7/(\sqrt{37}-4)=(4+\sqrt{37})/3$, as predicted.

\textbf{27.} For $D=23$, the triples $(P_n,Q_n,a_n)$ are
\[
 (0,1,4),\ (4,7,1),\ (3,2,3),\ (3,7,1),\ (4,1,8),\ (4,7,1).
\]
Thus $L=4$, the palindrome is $(1,3,1)$, and the final digit is
$8=2a_0$. The first four convergents are $4,5,19/4,24/5$, giving
$24^2-23\cdot5^2=1$. Including the final $8$ instead produces
$211/44$, with remainder $211^2-23\cdot44^2=-7$.

\textbf{28.} The period is $\sqrt6=[2;\overline{2,4}]$.
The four rows $(k,p_k/q_k,Q_{k+1},p_k^2-6q_k^2)$ are
\[
 (0,2/1,2,-2),\ (1,5/2,1,1),\
 (2,22/9,2,-2),\ (3,49/20,1,1).
\]
Their last entries equal $(-1)^{k+1}Q_{k+1}$. Since $L=2$, exactly
the odd-indexed convergents solve $+1$, and none solves $-1$.

\textbf{29.} The least periods for $D=2,3,5,7,13,14$ are, respectively,
\[
 (2),\ (1,2),\ (4),\ (1,1,1,4),\ (1,1,1,1,6),\ (1,2,1,6).
\]
Only $D=2,5,13$ have odd length. Their least positive solutions of
the negative equation are $(1,1),(2,1),(18,5)$.
For $D=3$, reducing $p^2-3q^2=-1$ would give
$p^2\congr2\pmod3$, impossible because a square is $0$ or $1$ modulo $3$.

\textbf{30.} Multiplication by $2-\sqrt3$ sends
\[
 (p,q)\longmapsto(2p-3q,\,2q-p)
\]
and preserves norm $1$. For $q=1$, positivity forces $p=2$, and
the image is $(1,0)$. If $q>1$, then
$\sqrt3<p/q=\sqrt{3+1/q^2}<2$. Hence $2p-3q>0$ and
$0<2q-p<q$. Repeated descent strictly decreases the positive integer
second coordinate until $(2,1)$ and then $(1,0)$ are reached.
Reversing the steps proves the power formula.

\textbf{31.} The state algorithm gives
\[
 \sqrt{61}=[7;\overline{1,4,3,1,2,2,1,3,4,1,14}],\qquad L=11.
\]
The palindrome is the ten-digit block before $14$.
At indices $10$ and $21$, respectively, the convergents give
\[
 29718^2-61\cdot3805^2=-1,
\]
\[
 1766319049^2-61\cdot226153980^2=1.
\]
The second also follows by squaring $29718+3805\sqrt{61}$.
Theorem~\ref{thm:pell-convergents} excludes every smaller positive
solution of either sign.

\textbf{32.} Clear denominators to check the seven equalities
\begin{gather*}
 6W=5B+6Y,\quad20B=9S+20Y,\quad42S=13W+42Y,\\
 12w=7(B+b),\quad20b=9(S+s),\\
 42y=13(W+w),\quad30s=11(Y+y).
\end{gather*}
The displayed table at $t=1$ satisfies each equality and has total
$50\,389\,082$. The first three equations determine all bulls from
$W$. The cow cycle has coefficient
$(7/12)(9/20)(11/30)(13/42)\ne1$, so it determines all cows as well.
Euclid gives gcd $4657$ for the four bull counts and gcd $1$ after
including the cows. A B\'ezout combination of the scaled counts
therefore equals $t$, proving $t\in\Z$ whenever the counts are
integral. At $t=1$, $W+B=4\cdot4\,456\,749$ is not a square,
so the full problem is not solved.

\textbf{33.} If $cr^2=h(h+1)/2$, then
$(2h+1)^2=8cr^2+1$. Conversely, an odd $X>1$ gives the integer
$h=(X-1)/2$ and $h(h+1)/2=(X^2-1)/8=cr^2$.
Here $c=51\,285\,802\,909\,803$.
The identity $8c=4\,729\,494\cdot9314^2$ makes
$(X,r)\mapsto(X,Z)=(X,9314r)$ a bijection onto the positive
solutions of the smaller equation with $9314\divides Z$.
For the reverse map, use $r=Z/9314$; the equation then implies
$X^2\congr1\pmod8$, so $X$ is odd.

\textbf{34.} The recurrence in Section~\ref{sec:archimedes-cattle}
computes the two coefficients of $\varepsilon^j$ modulo $9314$
by induction. Its first vanishing second coefficient occurs at
$j=2329$, with first coefficient $9313\congr-1$.
Because all positive Pell solutions are powers and their second
coordinates increase, this first hit gives the smallest herd.
Write any $j\ge1$ as $j=2329k+r$ with $0\le r<2329$.
Then $\varepsilon^j\congr(-1)^k\varepsilon^r\pmod{9314}$,
so its second coefficient vanishes exactly when $r=0$.
Finally, $T$ is proportional to $Z_j^2$, while
$Z_j=(\varepsilon^j-\varepsilon^{-j})/(2\sqrt{4\,729\,494})$.
Squaring doubles the exponent: $2\cdot2329=4658$.

\section*{Chapter 4}

\textbf{1.} Extended Euclid yields
\[
1=11\cdot101-30\cdot37,
\]
so $37^{-1}\congr-30\congr71\pmod{101}$.

\textbf{2.} Divide $18x\congr30\pmod{42}$ by $6$ while reducing the modulus: $3x\congr5\pmod7$. Since $3^{-1}\congr5$, $x\congr4\pmod7$. Modulo $42$ the six solutions are $4,11,18,25,32,39$.

\textbf{5.} Here $N=504$. One convenient construction is
\[
x=2(72)(4)+3(63)(7)+4(56)(5)=3019\congr499\pmod{504}.
\]
The multipliers $4,7,5$ invert $72,63,56$ modulo $7,8,9$, respectively.

\textbf{8.} The rows $(r,s,t)$ are
\[
 (83,1,0),\ (29,0,1),\ (25,1,-2),\ (4,-1,3),\
 (1,7,-20),\ (0,-29,83).
\]
Thus $1=7\cdot83-20\cdot29$ and $29^{-1}\congr63\pmod{83}$.
Indeed, $29\cdot63=1827=22\cdot83+1$.

\textbf{9.} From $6=42-2\cdot18$, multiplication by $4$ gives
$24=4\cdot42-8\cdot18$. Hence $x\congr-8\congr6\pmod7$;
the residues modulo $42$ are $6,13,20,27,34,41$.
The gcd $6$ does not divide $25$, so the second congruence is impossible.

\textbf{10.} At every component modulus, $e_i$ has value $0$ or $1$,
exactly one selector has value $1$, and distinct selectors cannot both
have value $1$. CRT transfers the three identities back modulo $N$.
Consequently
$(\sum a_i e_i)(\sum b_i e_i)=\sum a_i b_i e_i$ modulo $N$;
addition is componentwise as well.

\textbf{11.} Put $x=2+6t$. Then $6t\congr3\pmod9$, or
$2t\congr1\pmod3$, so $t\congr2\pmod3$ and
$x\congr14\pmod{18}$. The second system prescribes different residues
modulo $\gcd(6,9)=3$, so it is incompatible.

\textbf{12.} The selectors at $-1,0,2$ give
\[
 f(X)=2\frac{X(X-2)}3-\frac{(X+1)(X-2)}2
       +3\frac{X(X+1)}6
     =\frac{2X^2-X+3}{3}.
\]
All solutions are $f(X)+(X+1)X(X-2)h(X)$ with $h\in\Q[X]$.

\textbf{13.} The denominators at $0,1,3$ are $3,5,6$ modulo $7$.
Their inverse certificates are
\[
 3\cdot5-7\cdot2=1,\quad 5\cdot3-7\cdot2=1,\quad
 6\cdot6-7\cdot5=1.
\]
Thus
\[
 f(X)=5(X-1)(X-3)+4\cdot3X(X-3)+2\cdot6X(X-1)
     =X^2+2X+1\quad\text{in }\F_7[X].
\]
Substitution gives the required values $1,4,2$.

\textbf{14.} In $f(2)-f(0)$ every term has a factor $2$, so this
difference cannot be $1$ modulo $8$. The coefficient ring is not a
field; specifically, the difference between the two interpolation
points is a nonunit.

\textbf{15.} Modulo $2$, the orbit is
$(0,1),(1,1),(1,0),(0,1)$, giving period $3$.
Modulo $3$, it is
\[
 (0,1),(1,1),(1,2),(2,0),(0,2),(2,2),(2,1),(1,0),(0,1),
\]
giving period $8$. At the first positive zero, $n=4$, the next residue
is $2$, not $1$, so the recurrence has not restarted.

\textbf{16.} The periods are $\operatorname{lcm}(3,8)=24$ and
$\operatorname{lcm}(8,20)=40$. A simultaneous return occurs exactly at
an index divisible by both component periods. For modulus $15$,
the product $8\cdot20=160$ is a period but is four times the least one.

\textbf{17.} The orbit is $(1,0)\mapsto(0,0)\mapsto(0,0)$.
Modulo $2$ the rule is $(a,b)\mapsto(b,b)$, which forgets $a$;
a state therefore does not determine a unique predecessor.

\textbf{18.} The squares modulo $7$ are $0,1,2,4$. A root $x$ would
satisfy $(2x-1)^2=5$, which is impossible. Binet's formula requires
the missing roots (equivalently, a square root of $5$); the extension
in Appendix~\ref{app:finite-fields} supplies arithmetic for them.

\section*{Chapter 5}

\textbf{1.} The binary length-$3$ words form four orbits:
\[
\{000\},\quad\{111\},\quad\{001,010,100\},\quad\{011,110,101\}.
\]

\textbf{3.} If $ik\congr jk\pmod p$, then $p\divides k(i-j)$. As $p\ndivides k$, Euclid's lemma gives $p\divides i-j$. For $0\le i,j<p$, this forces $i=j$.

\textbf{5.} Orbit sizes divide $6$, and binary examples realize $1,2,3,6$: $000000$, $010101$, $001001$, and $000001$.

\textbf{6.} Fermat gives $(a+b)^p=a+b=a^p+b^p$ in $\F_p$.
Applying it repeatedly proves the corresponding identity for any finite
sum. Fermat's statement about residues does not assert $u^p=u$ for
new elements of an extension; the polynomial identity from
Chapter~\ref{chap:binomial} proves additivity there instead.

\section*{Chapter 6}

\textbf{1.} $(\Z/12\Z)^\times=\{1,5,7,11\}$. Every nonidentity element squares to $1$, so their orders are $2$; the identity has order $1$.

\textbf{4.} If $aH\cap bH$ contains $x$, write $x=ah_1=bh_2$. Then $a=b h_2h_1^{-1}$, so every $ah\in aH$ equals $b(h_2h_1^{-1}h)\in bH$. Thus $aH\subseteq bH$; symmetry gives equality.

\textbf{5.} For $g\ne e$, the order of $g$ divides the prime $\abs G$, so it is $\abs G$. Hence $G=\angles g$.

\section*{Chapter 7}

\textbf{2.} Legendre gives
\[
v_2(100!)=50+25+12+6+3+1=97,
\]
\[
v_5(100!)=20+4=24.
\]
Each trailing zero consumes one $2$ and one $5$, so there are $24$.

\textbf{4.} Choose $n$ objects from two disjoint groups of $n$. If $k$ come from the first group, then $n-k$ come from the second. Summing over $k$ gives
\[
\sum_k\binom nk\binom n{n-k}=\binom{2n}{n}.
\]

\textbf{7.} Legendre gives
\[
v_p\binom{2n}{n}=\sum_{j\ge1}\left(\floor{\frac{2n}{p^j}}-2\floor{\frac n{p^j}}\right).
\]
Each summand is $0$ or $1$, and there are at most $\floor{\log_p(2n)}$ nonzero candidates. This is the exponent of $p$ in the lcm, so the valuation of the binomial coefficient is no larger for every prime $p$.

\textbf{8.} Fermat makes $X^p-X$ zero at every point of $\F_p$,
but its leading coefficient is $1$, so it is not the zero polynomial.
By contrast, expansion shows that every coefficient of
$(X+Y)^p-X^p-Y^p$ vanishes modulo $p$. This is a polynomial identity
and remains valid upon substitution in any field containing $\F_p$.

\section*{Chapter 8}

\textbf{2.} The sum of the $2n+1$ coefficients in row $2n$ is $4^n$. The central coefficient is largest, so it is at least the average $4^n/(2n+1)$ and strictly below the whole sum.

\textbf{3.} For $n<p\le2n$, the numerator $(n+1)\cdots(2n)$ contains exactly one multiple of $p$, while $n!$ contains none. Because $p^2>2n$ for $p>n\ge1$, the valuation is exactly one.

\textbf{7 (hint).} Insert $x=p_n$ into both Chebyshev bounds to compare $n$ with $p_n/\log p_n$. First obtain coarse polynomial bounds on $p_n$ so that $\log p_n$ can be compared with $\log n$, then substitute back.

\section*{Appendix B: continued-fraction arithmetic}

\textbf{1.} With $t_2\in[1,\infty]$,
\[
 x=2+\frac1{3+1/t_2}.
\]
This expression increases with $t_2$.  At $t_2=1$ it is $9/4$; as
$t_2\to\infty$ it approaches $7/3$.  For an infinite continued fraction the
second endpoint is not attained, so $x\in[9/4,7/3)$.

\textbf{2.} The denominator is positive for $t\geq1$, and
\[
 F'(t)=\frac{22\cdot2-3\cdot14}{(14t+2)^2}
       =\frac2{(14t+2)^2}>0.
\]
Thus the range is determined by the endpoint values
$F(1)=25/16$ and $F(\infty)=11/7$.  The whole interval lies in $[1,2)$, so
the certified output digit is $1$.

\textbf{3.} Substituting $t=m+1/u$ into
$F_M(t)=(At+B)/(Ct+D)$ and clearing the factor $u$ gives
\[
 \frac{(Am+B)u+A}{(Cm+D)u+C},
\]
which is represented by $MD(m)$.  After emitting $k$, the unread output is
$1/(F_M-k)$, represented by
$\begin{psmallmatrix}0&1\\1&-k\end{psmallmatrix}M$.

\textbf{4.} From $M=\begin{psmallmatrix}2&0\\0&1\end{psmallmatrix}$,
absorption gives
\[
 MD(15)=\begin{pmatrix}30&2\\1&0\end{pmatrix},
 \qquad
 MD(15)D(1)=\begin{pmatrix}32&30\\1&1\end{pmatrix}.
\]
The latter map has range $[31,32)$ on $[1,\infty)$, so the next output digit
is $31$.

\textbf{5.} For each fixed $v$, the map in $u$ is linear fractional and hence
monotone unless it is constant; its extrema on a horizontal slice occur on the
two vertical edges.  On either edge, the remaining function of $v$ is again
linear fractional, so its extrema occur at the two endpoints.  Therefore all
global extrema are corner values.  If the denominator vanishes inside the
rectangle, the map is not continuous there and can become unbounded, so corner
values alone give no range certificate.

\textbf{7.} Treating the two copies as independent discards the relation
$x=y$.  Every finite continued-fraction prefix gives a nontrivial interval
containing $\sqrt2$, so the product box permits values slightly below and
slightly above $2$.  Since the floor jumps at $2$, it is never certified from
such boxes.  Retaining the dependency $x=y$ changes the expression to $x^2$;
using the known equation $x^2=2$ resolves it immediately.  Without that extra
algebraic relation, an interval-valued output can still certify arbitrarily
small numerical error.

\section*{Appendix C: Venus and astronomical cycles}

\textbf{1.} $E$ is Earth's Sun-centered lap time and $V$ is Venus's
Sun-centered lap time.  The quantity $S$ measures something different: the
waiting time between successive Earth--Venus alignments.

\textbf{2.} Venus and Earth move through $1/V$ and $1/E$ laps per day.
Their relative phase therefore increases at $1/V-1/E$ revolutions per day, so
one relative revolution takes $S$ days with $1/S=1/V-1/E$.  If $r=E/V$,
then
\[
  \frac ES=\frac EV-1=r-1,
  \qquad \frac SE=\frac1{r-1}.
\]
Substitution gives $T(13/8)=1/(13/8-1)=8/5$.

\textbf{3.} Repeated addition of $3/5$ modulo $1$ gives
\[
  0,\ \frac35,\ \frac65\congr\frac15,\
  \frac45,\ \frac75\congr\frac25,
\]
then returns to $0$.  Each step skips one neighboring vertex of the regular
pentagon, so the chronological edges are the diagonals and form a pentagram.

\textbf{5.} The factorizations are
\[
  2920=2^3\cdot5\cdot73,
  \qquad 260=2^2\cdot5\cdot13.
\]
Thus the gcd is $20$ and the lcm is
$2^3\cdot5\cdot13\cdot73=37960$.  Division by $584,365,260,$ and $2920$
gives the multipliers $65,104,146,$ and $13$.

\textbf{6.} The Fibonacci recurrence gives
\[
 T\!\left(\frac{F_{n+1}}{F_n}\right)
 =\frac{1}{F_{n+1}/F_n-1}
 =\frac{F_n}{F_{n+1}-F_n}
 =\frac{F_n}{F_{n-1}}.
\]
The positive solution of $r=1/(r-1)$ is
$r=(1+\sqrt5)/2=\varphi$.  Its identity
$\varphi=1+1/\varphi$ is exactly the infinite continued fraction
$[1;\overline1]$.

\textbf{7.} The nominal count gains $0.08$ day per cycle, hence $0.4$ day
after $5$ cycles, $5.2$ days after $65$, and $52$ days after $650$.

\section*{Appendix \ref{app:finite-fields}: extension arithmetic and periods}

\textbf{1.} The product is $13+7s=6$. Since
$1-5\cdot2^2=-19=2$ in $\F_7$, its inverse is
$4(1-2s)=4+6s$. Multiplication gives
$(1+2s)(4+6s)=64+14s=1$.

\textbf{2.} The elements are $0,1,t,t+1$. The polynomial
$X^2+X+1$ has no root in $\F_2$ and is irreducible. Directly,
$t(t+1)=t^2+t=1$, so the nonzero inverses are $1,t+1,t$,
respectively. Squaring fixes $0,1$ and exchanges $t,t+1$.

\textbf{3.} Write $u=a+bs$. Then $u^7=a-bs$, which equals $u$
exactly when $2b=0$, hence $b=0$. Applying the map twice gives
$u^{49}=u$. The example $s^7=-s\ne s$ exhibits its nonidentity action.

\textbf{4.} Both $6$ and $24$ satisfy $X^2-X-1=0$ modulo $29$.
We have $6^7=-1$, $6^{14}=1$, and $6^2=7\ne1$, so $6$ has
order $14$. Also $24^7=1$ and $24\ne1$, so $24$ has order $7$.
Their least common multiple is $14$, agreeing with factor removal.

\textbf{5.} Frobenius interchanges the roots, so
$\alpha^{p+1}=\beta^{p+1}=-1$. Binet gives $F_{p+1}=0$ and
$F_{p+2}=(-\alpha+\beta)/(\alpha-\beta)=-1$.
For $p=3$ this is precisely the state $(0,2)$ at index $4$;
the full return occurs at index $8$.

\textbf{6 (check values).} For primes $3,7,11,13,29$, valid starting
multiples are $8,16,10,28,28$ and the respective least periods are
$8,16,10,28,14$. The exceptional primes $2,5$ have periods $3,20$.
The modular arithmetic notebook gives implementations of direct pair
iteration, fast doubling, and factor removal for comparison.

\chapter{Reference sheets}
\chaptertag{Definitions, algorithms, and proof moves}

\section*{Core notation}

\begin{center}
\begin{tabularx}{\textwidth}{p{.22\textwidth}X}
\toprule
Notation & Meaning \\
\midrule
$a\divides b$ & $b=ak$ for some integer $k$ \\
$\gcd(a,b)$ & greatest positive common divisor \\
$a\congr b\pmod n$ & $n\divides(a-b)$ \\
$(\Z/n\Z)^\times$ & invertible residue classes modulo $n$ \\
$\varphi(n)$ & number of units modulo $n$ \\
$v_p(m)$ & exponent of $p$ in $m$ \\
$\pi(x)$ & number of primes at most $x$ \\
$\vartheta(x)$ & $\sum_{p\le x}\log p$ \\
$\psi(x)$ & $\sum_{p^k\le x}\log p$ \\
$f(x)\asymp g(x)$ & their ratio stays between positive constants eventually \\
\bottomrule
\end{tabularx}
\end{center}

\section*{Algorithms at a glance}

\textbf{Extended Euclid.} Divide until zero; the last nonzero remainder is the gcd; substitute backward to obtain $ax+by=\gcd(a,b)$.

\textbf{Sieve of Eratosthenes.} To find the primes through $N$, repeatedly
circle the least uncrossed integer and cross out its larger multiples. Stop
after the sieving primes $p\le\sqrt N$; a composite at most $N$ must have a
prime factor in that range.

\textbf{Existence versus a witness.} The classical dichotomy for
$x=(\sqrt2)^{\sqrt2}$ proves that one of
$(\sqrt2,\sqrt2)$ and $(x,\sqrt2)$ has a rational power, but does not say
which. The explicit choice
$a=\sqrt2$, $b=\log_{\sqrt2}3$ gives $a^b=3$, and unique factorization proves
that $b$ is irrational.

\textbf{Continued fraction.} Repeatedly take the integer part, subtract it, and invert. Convergents satisfy
\[
p_k=a_kp_{k-1}+p_{k-2},\qquad q_k=a_kq_{k-1}+q_{k-2}.
\]
The recurrences are the columns of the matrix product
\[
\begin{pmatrix}a_0&1\\1&0\end{pmatrix}\cdots
\begin{pmatrix}a_k&1\\1&0\end{pmatrix}
=\begin{pmatrix}p_k&p_{k-1}\\q_k&q_{k-1}\end{pmatrix},
\]
whose determinant is $(-1)^{k+1}$. Each factor lies in
$\operatorname{GL}_2(\Z)$; a product with an even number of factors lies in
$\operatorname{SL}_2(\Z)$.
For $A/B=[a_0;\dots,a_n]$, the Euclidean remainders obey
\[
r_k=(-1)^k(q_kA-p_kB).
\]
For irrational $x$, the zero-indexed convergents form alternating monotone bounds:
\[
c_0<c_2<c_4<\cdots<x<\cdots<c_5<c_3<c_1.
\]
Legendre's test says that a reduced fraction satisfying
$\abs{x-p/q}<1/(2q^2)$ must be a convergent. Lagrange's theorem says that a
simple continued fraction is eventually periodic exactly for a quadratic
irrational; Galois says it is purely periodic exactly when $x>1$ and
$-1<x'<0$.

\textbf{Fixed points and prefixes.} For $\det M\ne0$,
$T_M(t)=(at+b)/(ct+d)$ satisfies $T_M\circ T_N=T_{MN}$.
An irrational fixed point of a nonscalar rational $M$ obeys the genuine
quadratic $ct^2+(d-a)t-b=0$. If $x=T_N(x)$ and $y=T_M(x)$, then
\[
 y=T_{MNM^{-1}}(y).
\]
A finite continued-fraction prefix therefore preserves quadratic
irrationality by conjugating the period matrix.

\textbf{Reduced quadratic states.} Write $x_n=(P_n+\sqrt D)/Q_n$ with
$Q_n\divides D-P_n^2$. With $a_n=\floor{x_n}$, the exact update is
\[
 P_{n+1}=a_nQ_n-P_n,\qquad Q_{n+1}=(D-P_{n+1}^2)/Q_n.
\]
Every quadratic irrational eventually has a reduced tail.
For reduced states, $|P_n|<\sqrt D$ and $0<Q_n<2\sqrt D$,
so only finitely many states occur. A reduced successor $z$ determines
its reduced predecessor:
$a=\floor{-1/z'}$ and $u=a+1/z$.
Finiteness gives eventual repetition; unique predecessors make a
reduced starting state purely periodic.

\textbf{Typical information rate.} For almost every $x$, the Khinchin--L\'evy
law gives
\[
 \lim_{n\to\infty}q_n^{1/n}=e^{\pi^2/(12\log2)}\approx3.275823,
 \qquad
 \lim_{n\to\infty}\abs{x-p_n/q_n}^{-1/n}
 =e^{\pi^2/(6\log2)}\approx10.731.
\]
So a continued-fraction term earns about $1.0306$ decimal digits on average.
For the first $100{,}000$ partial quotients of $\pi$, computation gives
$G_n(\pi)\approx2.67956$ and $q_n^{2/n}\approx10.67835$; neither limiting law
has been proved for $\pi$. In contrast,
$e=[2;1,2,1,1,4,1,1,6,\dots]$ has
$G_{3m}(e)\sim(2m/e)^{1/3}\to\infty$ and $q_{3m}^{2/(3m)}\to\infty$.

\textbf{Gosper's streaming arithmetic.} Represent a unary rational
transformation by $M=\begin{psmallmatrix}A&B\\C&D\end{psmallmatrix}$.  Reading
an input digit $m$ and emitting a certified output digit $k$ update the state by
\[
 M\longmapsto M\begin{pmatrix}m&1\\1&0\end{pmatrix},
 \qquad
 M\longmapsto\begin{pmatrix}0&1\\1&-k\end{pmatrix}M.
\]
For two inputs use
$(axy+bx+cy+d)/(exy+fx+gy+h)$; its eight integer coefficients are closed
under the same absorb-or-emit operations.  Exact interval bounds certify a
digit when every possible value has the same floor.  Interval-valued output is
the safe fallback at integer boundaries where a digit-only process can stall.

\textbf{Calendar design.} If a tropical year is $365+\alpha$ days, a cycle
with $p$ leap days in $q$ years asks for $p/q\approx\alpha$. For
$\alpha\approx0.2421988$, the convergents begin
\[
 \frac14,\quad\frac7{29},\quad\frac8{33},\quad\frac{31}{128}.
\]
The Julian calendar uses $1/4$. Requiring $q=100Q$ means approximating
$100\alpha$; its convergent $97/4$ yields the Gregorian rate $97/400$.

\textbf{Planetary lapping time.} If an inner planet takes $V$ days for one
Sun-centered lap and Earth takes $E>V$ days, the waiting time $S$ between
successive alignments (the \emph{synodic period}) satisfies
\[
  \frac1S=\frac1V-\frac1E,
  \qquad \frac SE=T(E/V),\qquad T(r)=\frac1{r-1}.
\]
For Earth and Venus, $E/V\approx13/8$ and $S/E\approx8/5$; the five
successive conjunction directions advance by about $3/5$ of a turn and draw a
near-pentagram.  The Dresden Codex uses the exact nominal calendar identities
\[
  5\cdot584=8\cdot365=2920,
  \qquad
  65\cdot584=104\cdot365=146\cdot260=37960.
\]

\textbf{Square-root periods.} For a nonsquare integer $D\ge2$,
\[
 \sqrt D=[a_0;\overline{a_1,\ldots,a_{L-1},2a_0}],
 \qquad a_j=a_{L-j}\quad(1\le j<L),
\]
where $L$ is the least period length. The block before $2a_0$ is a
palindrome. The complete quotients satisfy
$Q_n=1$ exactly when $L\divides n$, for $n\ge1$.

\textbf{Pell bridge.} If $p^2-Dq^2=\pm1$ with $p,q>0$, then
\[
 \abs{\sqrt D-p/q}=\frac{1}{q^2(\sqrt D+p/q)}<\frac1{2q^2},
\]
so Legendre makes $p/q$ a convergent. The exact identity
\[
 p_{n-1}^2-Dq_{n-1}^2=(-1)^nQ_n
\]
shows that the Pell convergents have indices $mL-1$ and sign $(-1)^{mL}$.
The first $+1$ solution is at $L-1$ for even $L$ and $2L-1$ for odd $L$;
the $-1$ equation is solvable exactly when $L$ is odd.
Stop before the final $2a_0$. If $(u,v)$ is the first $+1$ solution,
every positive $+1$ solution is a power of $u+v\sqrt D$, or equivalently
is generated by $\begin{psmallmatrix}u&Dv\\v&u\end{psmallmatrix}$.

\textbf{Cattle reduction.} Archimedes' seven proportions give the total
$50\,389\,082t$. The square condition forces $t=4\,456\,749r^2$;
the triangular condition is equivalent to
\[
 X^2-4\,729\,494Z^2=1,\qquad Z=9314r.
\]
The divisibility condition $9314\divides Z$ must be retained. The first
admissible power of the fundamental unit is $2329$; the resulting
smallest total has $206\,545$ digits. See
Section~\ref{sec:archimedes-cattle} and Vardi's exposition.

\textbf{Base-$b$ expansion.} Long division evolves by
$br_n=qd_{n+1}+r_{n+1}$ with $0\le r_{n+1}<q$. Repeated remainders make every
rational expansion eventually periodic. A reduced fraction $m/q$ is purely
periodic exactly when $\gcd(q,b)=1$; the reverse direction follows from
Euler's theorem.

\textbf{Linear congruence.} For $ax\congr b\pmod n$, first compute $d=\gcd(a,n)$. There are solutions iff $d\divides b$. Divide by $d$, reduce the modulus to $n/d$, and invert the new coefficient.

\textbf{CRT.} With $N=\prod n_i$, take $N_i=N/n_i$ and choose $M_i=N_i^{-1}\pmod{n_i}$. Then $x=\sum a_iN_iM_i\pmod N$.

\section*{Proof moves worth naming}

\begin{itemize}
  \item \textbf{Minimal counterexample / strong induction}: factor a composite object into smaller ones.
  \item \textbf{B\'ezout insertion}: replace $1$ by $ax+by$ to transfer divisibility.
  \item \textbf{Exponent comparison}: use unique factorization or $v_p$ to compare products and quotients.
  \item \textbf{M\"obius monotonicity}: compute the derivative of $(at+b)/(ct+d)$ to determine how a continued-fraction tail moves its value.
  \item \textbf{Conjugation}: transport a fixed-point equation through a change of variable using $MNM^{-1}$.
  \item \textbf{Finite states with reversible steps}: bound integer states to force eventual repetition, then recover predecessors to remove any initial segment.
  \item \textbf{Equal-block partition}: cosets and rotation orbits turn a count into divisibility.
  \item \textbf{Count two ways}: binomial identities become transparent when both sides count the same set.
  \item \textbf{Logarithmic weighting}: products of primes become the sums $\vartheta$ and $\psi$.
  \item \textbf{Dyadic decomposition}: split a range into $(x/2,x],(x/4,x/2],\dots$ and sum a geometric series.
\end{itemize}

\section*{The two Fermat proofs}

\begin{center}
\small
\begin{tabularx}{\textwidth}{p{.18\textwidth}XX}
\toprule
 & Necklace proof & Group proof \\
\midrule
Objects & $a^p$ labeled colorings & nonzero residues modulo $p$ \\
Structure & cyclic rotation orbits & group $(\Z/p\Z)^\times$ \\
Exceptional part & $a$ constant colorings & identity/powers inside cyclic subgroups \\
Equal blocks & nonconstant orbits have size $p$ & element orders divide $p-1$ \\
Conclusion & $a^p\congr a\pmod p$ & $a^{p-1}\congr1\pmod p$ for $p\ndivides a$ \\
\bottomrule
\end{tabularx}
\end{center}


\backmatter
\chapter*{Further reading}
\addcontentsline{toc}{chapter}{Further reading}
\markboth{Further reading}{}
\begin{itemize}[parsep=0pt,itemsep=2pt,topsep=4pt]
  \item G. H. Hardy and E. M. Wright, \emph{An Introduction to the Theory of Numbers}.
  \item I. Niven, H. Zuckerman, and H. Montgomery, \emph{An Introduction to the Theory of Numbers}.
  \item J. Stillwell, \emph{Elements of Number Theory}.
  \item R. M. Corless,
  \href{https://www.cecm.sfu.ca/organics/papers/corless/confrac/html/node6-an.shtml}{\emph{Continued Fractions and Chaos}, section on periodic continued fractions}
  (the distinction between Lagrange's theorem and Galois's criterion).
  \item G. Panti,
  \href{https://arxiv.org/abs/0712.2996}{\emph{A General Lagrange Theorem}}
  (2007; a broader perspective beyond simple continued fractions).
  \item I. Vardi,
  \href{https://doi.org/10.1080/00029890.1998.12004887}{``Archimedes' Cattle Problem,''}
  \emph{American Mathematical Monthly} \textbf{105}, no.~4 (1998), 305--319
  (Pell equations, divisibility conditions, and a compact description of
  the enormous least solution;
  \href{https://www.maths.tcd.ie/~lebed/Cattle\%20problem\%20Ilan\%20Vardi.pdf}{preprint}).
  \item A. Ya. Khinchin, \emph{Continued Fractions} (for the metric theory behind the ``base 11'' aside).
  \item P. L\'evy, ``Sur le d\'eveloppement en fraction continue d'un nombre choisi au hasard,'' \emph{Compositio Mathematica} \textbf{3} (1936), 286--303.
  \item G. Lochs, ``Vergleich der Genauigkeit von Dezimalbruch und Kettenbruch,'' \emph{Abh. Math. Sem. Univ. Hamburg} \textbf{27} (1964), 142--144.
  \item D. H. Bailey, J. M. Borwein, and R. E. Crandall, ``On the Khintchine Constant,'' \emph{Mathematics of Computation} \textbf{66} (1997), 417--431.
  \item H. Cohn, ``A Short Proof of the Simple Continued Fraction Expansion of $e$,'' \emph{American Mathematical Monthly} \textbf{113} (2006), 57--62.
  \item F. Eisenbrand, ``Pope Gregory, the calendar, and continued fractions,'' in \emph{Discrete Mathematics}, DOI: 10.4171/DMS/6/14.
  \item R. W. Gosper, ``Continued Fraction Arithmetic,'' Item 101B in M. Beeler,
  R. W. Gosper, and R. Schroeppel, \emph{HAKMEM}, MIT Artificial Intelligence
  Memo 239 (1972), 39--44.
  \item M. J. Collins, ``Arithmetic on Continued Fractions,'' \emph{Qeios}
  \textbf{7} (2025), DOI: 10.32388/DD1PU3.2.
  \item NASA Goddard Space Flight Center,
  \href{https://eclipse.gsfc.nasa.gov/OH/transit12.html}{``2012 Transit of Venus''}
  (orbital and synodic periods of Venus).
  \item European Southern Observatory,
  \href{https://www.eso.org/public/outreach/eduoff/vt-2004/Education/edu1app5.html}{``The Pentagram of Venus''}.
  \item Saxon State and University Library Dresden,
  \href{https://www.slub-dresden.de/en/explore/manuscripts/the-dresden-maya-codex/content}{``The Dresden Maya Codex: Content''}.
  \item A. Kaneko, ``Identification of Postclassic Maya Constellations from
  the Venus Pages of the Dresden Codex,'' \emph{Estudios de Cultura Maya}
  \textbf{35} (2010), 33--74.
  \item J.-P. Serre, \emph{A Course in Arithmetic} (for a more advanced next step).
\end{itemize}

\vfill
\begin{center}\color{muted}\small
Course packet designed as an adaptable teaching draft. Classical results
and examples are attributed where appropriate; exposition, exercises, and
diagrams prepared for this packet.
\end{center}

\end{document}
